% Mechanically extracted after \mainmatter from the frozen standalone master.

\markboth{Exposé XV. Complements on Sub-tori}{Exposé XV. Complements on Sub-tori}
% BEGIN inputs/XV/tex/components/00_expose_XV_title_and_introduction.tex
% Authority: Expo15.pdf, local p. 1 / combined p. 871 / printed p. 349.

\chapter*{Exposé XV\\Complements on Sub-tori of a Group Prescheme\\
Application to Smooth Groups}
\addcontentsline{toc}{chapter}{Exposé XV. Complements on Sub-tori of a Group Prescheme}
\SGARefTarget{sga3:xv:expose}

\begin{center}
\textit{by M. Raynaud}\footnote{Source-edition note: xy version dated
1 December 2008.}\textsuperscript{,}\footnote{This Exposé and the two following ones
(Exposés XVI and XVII) do not correspond to oral lectures delivered at
the seminar. With some additional material, they develop the substance
of concise manuscript notes written by A. Grothendieck in 1964 for the
present seminar.}
\end{center}

\SGASection{0}{Introduction}
\SGARefTarget{sga3:xv:section:0}

This Exposé supplements and partly recasts Exposés XI and XII; knowledge
of Exposés XIII and XIV is not indispensable. Continuing the effort
begun in XII, we shall work with \(S\)-group preschemes that are not
necessarily affine and not necessarily separated over \(S\).

\SGARefLink{sga3:xv:section:1}{Sections 1}, 2, 3, and 4 are devoted to the study of sub-tori of a group
prescheme. We obtain infinitesimal lifting theorems (§2) and a global
lifting theorem (§4), in which points of finite order (§1) play an
essential role.

\SGARefLink{sga3:xv:section:5}{Sections 5}, 6, and 7 are independent of the preceding sections. The
consideration of infinitesimal neighbourhoods leads to the
representability of the functor of smooth subgroups equal to their
connected normalizer (§5). In §§6 and 7 we turn more particularly to
Cartan subgroups.

Finally, in §8 we give a necessary and sufficient condition for the
functor of sub-tori of a smooth group, or the functor of maximal tori,
to be representable.

% END inputs/XV/tex/components/00_expose_XV_title_and_introduction.tex

% BEGIN inputs/XV/tex/components/01_expose_XV_section1_opening.tex
% Authority: Expo15.pdf, local pp. 1--3 / combined pp. 871--873 /
% printed pp. 350--353.
% Stops immediately before §1.2, Lemma 1.4.

\SGASection{1}{Lifting finite subgroups}
\SGARefTarget{sga3:xv:section:1}

\SGASubsection{1.1}{Finite, smooth, central subgroups of multiplicative type}
\SGARefTarget{sga3:xv:subsection:1:1}

\SGAStatement{Proposition 1.1}{Let \(S\) be an affine scheme, let
\(S_0\) be a closed subscheme of \(S\) defined by a square-zero ideal,
let \(G\) be an \(S\)-group prescheme, and let \(H_0\) be a subgroup
scheme of \(G_0=G\times_S S_0\) that is smooth and of finite type over
\(S_0\), and of multiplicative type. A subgroup scheme of \(G\), of
multiplicative type, lifting \(H_0\) exists if and only if there exists
a subscheme \(H'\) of \(G\), flat over \(S\), lifting \(H_0\).}
\SGARefTarget{sga3:xv:proposition:1:1}

The necessity of the condition is clear; let us prove its sufficiency.
The group \(H_0\) of multiplicative type is quasi-isotrivial (\SGARefLink{sga3:x:corollary:4-5}{X~4.5}).
By Exposé \SGARefLink{sga3:x:proposition:2:1}{X~2.1}, there exist an \(S\)-group \(H\) of multiplicative
type and an \(S_0\)-isomorphism of groups
\[
  u_0:H\times_S S_0\xrightarrow{\sim}H_0.
\]
Since \(H'\) is flat over \(S\) and \(H'\times_S S_0\) is of finite
presentation over \(S_0\) (Exposé \SGARefLink{sga3:ix:proposition:2:1:item:b:01}{IX~2.1(b)}), \(H'\) is of finite
presentation over \(S\). Its fibres are smooth, so \(H'\) is smooth
over \(S\) (EGA IV~17.5.1). Since \(S\) is affine, \(u_0\) therefore
lifts to an \(S\)-morphism of preschemes
\[
  u:H\longrightarrow H'.
\]
It then follows from Exposé \SGARefLink{sga3:iii:theorem:2-1}{III~2.1} and Exposé \SGARefLink{sga3:ix:theorem:3:1}{IX~3.1} that the composite
\[
  v_0:H\times_S S_0\xrightarrow{\sim}H_0\longrightarrow G_0
\]
also lifts to an \(S\)-homomorphism
\[
  v:H\longrightarrow G.
\]
Since \(v_0\) is an immersion, so is \(v\). The image of \(H\) under
\(v\) is therefore a subgroup scheme of \(G\), of multiplicative type,
lifting \(H_0\).

\SGAStatement{Proposition 1.2}{Let \(S\) be a prescheme, let \(S_0\)
be a subprescheme of \(S\) defined by a locally nilpotent sheaf of
ideals, let \(G\) be an \(S\)-group prescheme flat and of finite
presentation over \(S\), and let \(H_0\) be a subgroup scheme of
\(G_0=G\times_S S_0\) that is smooth, finite over \(S_0\), of
multiplicative type, and central. Then there exists a unique subgroup
scheme \(H\) of \(G\), of multiplicative type, lifting \(H_0\).
Moreover, \(H\) is central. (See XVII, Appendix III, 1.)}
\SGARefTarget{sga3:xv:proposition:1:2}

\SGAInlineRefTarget{sga3:xv:proposition:1-2-bis}\SGAStatement{Proposition 1.2 bis}{Let \(S,G,S_0,G_0\) be as above,
let \(H\) be an \(S\)-group scheme of multiplicative type, smooth and
finite over \(S\), and let
\(u_0:H_0=H\times_S S_0\to G_0\) be a central homomorphism. Then
\(u_0\) lifts uniquely to a homomorphism \(u:H\to G\). Moreover,
\(u\) is central.}
\SGARefTarget{sga3:xv:proposition:1:2:bis}

The existence of the lifting \(u\) in \SGARefLink{sga3:xv:proposition:1:2:bis}{Proposition 1.2 bis} follows
readily from \SGARefLink{sga3:xv:proposition:1:2}{Proposition 1.2} by considering the graph of \(u_0\). The
lifting \(u\) is unique and central by Exposé \SGARefLink{sga3:ix:corollary:3:4}{IX~3.4} and Exposé \SGARefLink{sga3:ix:theorem:5:1}{IX~5.1}.

\SGAProof{Proof of \SGARefLink{sga3:xv:proposition:1:2}{Proposition 1.2}}
The uniqueness of \(H\), and the fact that \(H\) is central, follow
from Exposé \SGARefLink{sga3:ix:theorem:5:6:item:b}{IX~5.6(b)} and Exposé \SGARefLink{sga3:ix:corollary:3:4:bis}{IX~3.4 bis}. In view of uniqueness, to
prove existence we may assume that \(S\) is affine and that \(S_0\) is
defined by a square-zero ideal; by \SGARefLink{sga3:xv:proposition:1:1}{Proposition 1.1} it is enough to find
a subscheme of \(G\), flat over \(S\), lifting \(H_0\).

Since \(H_0\) is smooth and finite over \(S_0\), after possibly
restricting \(S\) we may assume that there is an integer \(n>0\),
invertible on \(S\), such that \(H_0={}_nH_0\). Consider the \(n\)-th
power morphism on \(G\):
\[
  u:G\longrightarrow G,\qquad x\longmapsto x^n.
\]
We continue to denote by \({}_nG\) the ``kernel of \(u\)'', that is,
the inverse image under \(u\) of the identity section of \(G\).
Notice that \(u\) is not, in general, a group homomorphism. We first
admit the following lemma.

\SGAStatement{Lemma 1.3}{Let \(k\) be a field, let \(G\) be a group
scheme locally of finite type over \(k\), let \(n\) be an integer prime
to the characteristic of \(k\), and let \(u\) be the \(n\)-th power
morphism on \(G\). Then \(u\) is étale at every point \(x\) of \(G\)
that belongs to the centre of \(G\).}
\SGARefTarget{sga3:xv:lemma:1:3}

Since \(G\) is flat and of finite presentation over \(S\), the
preceding lemma and EGA IV~17.8.2 show that, if \(x\in G\) lies above
\(s\in S\) and belongs to the centre of \(G_s\), then \(u\) is étale
at \(x\). If in addition \(x\in {}_nG\), then \({}_nG\) is étale over
\(S\) at \(x\). By hypothesis, \(H_0\) is central and is contained in
\({}_nG_0\); it is therefore contained in the largest open subset
\(V\) of \({}_nG\) that is étale over \(S\). Since \(H_0\) and
\(V\times_S S_0=V_0\) are étale over \(S_0\), \(H_0\) is an open
subscheme of \(V_0\) (EGA IV~17.8.7 and 17.9.1). The open subprescheme
of \(V\) having the same underlying space as \(H_0\) is consequently
a subscheme of \(G\), flat over \(S\), lifting \(H_0\).

It remains to prove \SGARefLink{sga3:xv:lemma:1:3}{Lemma 1.3}. The largest open subset of \(G\) on
which \(u\) is étale is invariant under extension of the base field
(EGA IV~17.8.2), so we may reduce to the case in which \(x\) is
\(k\)-rational. Let \(t_x\) denote translation by \(x\), a
\(k\)-automorphism of \(G\). Since \(x\) belongs to the centre of
\(G\), the diagram
\[
\begin{tikzcd}[column sep=large,row sep=large]
G \arrow[r,"u"] \arrow[d,"t_x"'] & G \arrow[d,"t_x^n"] \\
G \arrow[r,"u"] & G .
\end{tikzcd}
\]
is commutative. It is therefore enough to prove that \(u\) is étale
at the identity; this was established in VII\textsubscript{A}~\SGARefLink{sga3:xv:definition:8-4}{8.4}.

% END inputs/XV/tex/components/01_expose_XV_section1_opening.tex

% BEGIN inputs/XV/tex/components/02_expose_XV_section1_remainder_and_section2.tex
% Authority: Expo15.pdf, local pp. 3--14 / combined pp. 873--884 /
% printed pp. 353--374.
% Begins at §1.2, Lemma 1.4, and stops immediately before §3.

\SGASubsection{1.2}{Global lifting of finite groups}
\SGARefTarget{sga3:xv:subsection:1:2}

\SGAStatement{\SGAInlineRefTarget{sga3:xv:lemma:1-4}Lemma 1.4}{Let \(A\) be a local ring, separated and
complete for the topology defined by its maximal ideal \(\mathfrak m\),
let
\[
  S=\operatorname{Spec}(A),\qquad
  S_n=\operatorname{Spec}(A/\mathfrak m^n).
\]
Then, for every prescheme \(X\) (respectively, for every \(S\)-prescheme
\(X\)), the canonical map
\[
  \tag{$\ast$}\label{sga3:xv:lemma:1-4:equation:ast}
  \operatorname{Hom}(S,X)
  \longrightarrow
  \varprojlim_n\operatorname{Hom}(S_n,X)
\]
(respectively,
\[
  \tag{$\ast\ast$}\label{sga3:xv:lemma:1-4:equation:ast-ast}
  \Gamma(X/S)
  \longrightarrow
  \varprojlim_n\Gamma(X_n/S_n),
  \qquad X_n=X\times_S S_n)
\]
is bijective.}
\SGARefTarget{sga3:xv:lemma:1:4}

Statement \((\ast\ast)\) is an easy consequence of \((\ast)\). Let us
prove \((\ast)\). Let
\[
  u_n:S_n\longrightarrow X\qquad(n\in\mathbb N)
\]
be a coherent system of morphisms. The image \(y\) of the closed point
of \(S_n\) is independent of \(n\), and \(u_n\) factors through
\(\operatorname{Spec}\mathcal O_y\). On passing to the inverse limit,
the morphisms \(u_n\) define a ring homomorphism
\[
  \widetilde u:\mathcal O_y
  \longrightarrow
  \varprojlim_n(A/\mathfrak m^n)=A.
\]
This proves that \((\ast)\) is surjective; it is injective as soon as
\(A\) is separated for the \(\mathfrak m\)-adic topology.

% Authority page transition: local p. 4 / combined p. 874 / printed p. 354.
\SGAStatement{\SGAInlineRefTarget{sga3:xv:corollary:1-5}Corollary 1.5}{Let \(A\) be a complete noetherian local
ring, let \(\mathfrak m\) be its maximal ideal, let
\[
  S=\operatorname{Spec}(A),\qquad
  S_n=\operatorname{Spec}(A/\mathfrak m^n),
\]
let \(X\) be a finite scheme over \(S\), and let \(Y\) be an
\(S\)-prescheme. Then the canonical map
\[
  \operatorname{Hom}_S(X,Y)
  \longrightarrow
  \varprojlim_n\operatorname{Hom}_{S_n}(X_n,Y_n)
\]
(where \(X_n=X\times_S S_n\), and similarly for \(Y_n\)) is
bijective.}
\SGARefTarget{sga3:xv:corollary:1:5}

Indeed, it follows from EGA II~6.2.5 that \(X\) is a finite disjoint
sum of local \(S\)-schemes finite over \(S\). We are thus reduced to
the case in which \(X\) itself is the spectrum of a complete
noetherian local ring. But
\[
  \operatorname{Hom}_S(X,Y)=\Gamma(Z/X),
  \qquad Z=Y\times_S X,
\]
and the preceding lemma applies.

\SGAStatement{\SGAInlineRefTarget{sga3:xv:proposition:1-6}Proposition 1.6}{Let \(A,S,S_n\) be as above, and let
\(G\) and \(M\) be two \(S\)-group preschemes, with \(M\) finite over
\(S\). Then:
\begin{enumerate}
\item[\normalfont a)]\SGAInlineRefTarget{sga3:xv:proposition:1-6:item:a} The canonical map
\[
  \operatorname{Hom}_{S\text{-}\mathrm{gr}}(M,G)
  \longrightarrow
  \varprojlim_n
  \operatorname{Hom}_{S_n\text{-}\mathrm{gr}}(M_n,G_n)
\]
is bijective.
\item[\normalfont b)] If \(M\) is of multiplicative type and \(G\) is
smooth over \(S\), the canonical map
\[
  \varphi:\operatorname{Hom}_{S\text{-}\mathrm{gr}}(M,G)
  \longrightarrow
  \operatorname{Hom}_{S_0\text{-}\mathrm{gr}}(M_0,G_0)
\]
is surjective. Moreover, if
\(\varphi(u)=\varphi(u')=u_0\), then \(u\) and \(u'\) are conjugate by
an element of \(G(S)\) that reduces to the identity element of
\(G(S_0)\).
\item[\normalfont c)] If \(M\) is of multiplicative type and smooth
over \(S\), if \(G\) is flat and of finite type over \(S\), and if
\(u_0:M_0\to G_0\) is a central homomorphism, then \(u_0\) lifts
uniquely to a homomorphism
\[
  u:M\longrightarrow G.
\]
Moreover, \(u\) is central if \(G\) has connected fibres.
\end{enumerate}}
\SGARefTarget{sga3:xv:proposition:1:6}

\SGAProof{Proof}
\emph{a)} This follows from \SGARefLink{sga3:xv:corollary:1-5}{1.5}, from the fact that
\(M\times_S M\) is finite over \(S\), and from the characterization
of group morphisms as the morphisms for which the familiar diagram
\[
\begin{tikzcd}[column sep=large,row sep=large]
M\times_S M \arrow[r,"u\times u"] \arrow[d]
  & G\times_S G \arrow[d] \\
M \arrow[r,"u"'] & G
\end{tikzcd}
\]
is commutative.

\emph{b)} By Exposé \SGARefLink{sga3:ix:theorem:3:6}{IX~3.6}, one can construct a coherent system of
homomorphisms \(u_n:M_n\to G_n\) lifting a homomorphism
\(u_0:M_0\to G_0\). In view of part a), this proves the first assertion
of b).

If \(u\) and \(u'\) are two liftings of \(u_0\), then \(u_n\) and
\(u'_n\) are conjugate by an element \(g_n\in G(S_n)\) lifting the
identity element of \(G(S_0)\) (Exposé \SGARefLink{sga3:ix:theorem:3:6}{IX~3.6}). The same reference
also implies that the \(g_n\) may be chosen coherently; by \SGARefLink{sga3:xv:lemma:1-4}{1.4} they
therefore arise from a section \(g\in G(S)\). The morphisms \(u\) and
\(\operatorname{Int}(g)u'\) then agree modulo \(\mathfrak m^n\) for
every \(n\), and hence agree by \SGARefLink{sga3:xv:corollary:1-5}{1.5}.

\emph{c)} The existence and uniqueness of \(u\) follow from part a)
and \SGARefLink{sga3:xv:proposition:1:2:bis}{1.2 bis}. If \(G\) has connected fibres, \(u\) is central by
Exposé \SGARefLink{sga3:ix:theorem:5:6:item:a}{IX~5.6(a)}.

% Authority page transition: local p. 5 / combined p. 875 / printed pp. 356--357.
\SGASection{2}{Infinitesimal lifting of sub-tori}
\SGARefTarget{sga3:xv:section:2}

\SGASubsection{2.1}{Statement of the theorem}
\SGARefTarget{sga3:xv:subsection:2:1}

We shall give an infinitesimal lifting theorem for sub-tori of a group
prescheme that does not require smoothness hypotheses (in contrast
with Exposé \SGARefLink{sga3:ix:theorem:3:6:bis}{IX~3.6 bis}) and that also answers a very natural
question:\footnote{Editor's note: The idea of approximating a torus by
its finite subgroups occurs in Grothendieck's proof of the
connectedness of centralizers of tori; see §\SGARefLink{sga3:xv:lemma:4:6}{4.6} of \emph{BIBLE}.}
is it enough to be able to lift ``sufficiently many'' points of finite
order of a sub-torus \(T_0\) in order to ensure that \(T_0\) itself can
be lifted (infinitesimally, of course)?

\SGAStatement{Theorem 2.1}{Let \(S\) be a noetherian affine scheme,
let \(S_0\) be a closed subscheme of \(S\) defined by a square-zero
ideal \(J\), let \(G\) be an \(S\)-group prescheme of finite type, let
\(G_0=G\times_S S_0\), let \(T_0\) be a sub-torus of \(G_0\), and let
\(q>0\) be an integer invertible on \(S\). Suppose that, for every
integer \(n\) equal to a power of \(q\), there exists a subscheme
\(M_n\) of \(G\), flat over \(S\), such that
\[
  M_n\times_S S_0={}_nT_0.
\]
Then there exists a sub-torus \(T\) of \(G\) such that
\(T\times_S S_0=T_0\).}
\SGARefTarget{sga3:xv:theorem:2:1}

\SGARefLink{sga3:xv:theorem:2:1}{Theorem 2.1} will be useful to us through the following two
corollaries.

\SGAStatement{\SGAInlineRefTarget{sga3:xv:corollary:2-2}Corollary 2.2}{Let \(S\) be a locally noetherian
prescheme, let \(S_0\) be a closed subprescheme of \(S\) defined by a
locally nilpotent sheaf of ideals, let \(G\) be an \(S\)-group
prescheme of finite type, let \(T_0\) be a sub-torus of
\(G_0=G\times_S S_0\), and let \(q>0\) be an integer invertible on
\(S\). As \(n\) ranges over the powers of \(q\), let \((M_n)\) be a
coherent system of \(S\)-subgroup schemes of \(G\), of multiplicative
type, lifting the \({}_nT_0\). (The system \((M_n)\) is said to be
coherent if
\[
  M_m={}_m(M_n)
\]
whenever \(m\) divides \(n\).) Then there exists a unique sub-torus
\(T\) of \(G\) such that
\[
  T\times_S S_0=T_0
  \qquad\text{and}\qquad
  {}_nT=M_n
\]
for every \(n\).}
\SGARefTarget{sga3:xv:corollary:2:2}

\SGAStatement{Corollary 2.3}{Let \(G\) be an \(S\)-group prescheme,
flat and of finite presentation over \(S\), let \(S_0\) be a closed
subprescheme of \(S\) defined by a locally nilpotent sheaf of ideals
of finite type, and let \(T_0\) be a central torus of
\(G_0=G\times_S S_0\). Then there exists a unique sub-torus \(T\) of
\(G\) lifting \(T_0\). Moreover, \(T\) is central.}
\SGARefTarget{sga3:xv:corollary:2:3}

\SGAStatement{Remark 2.4}{We leave it to the reader to formulate the
analogues of statements \SGARefLink{sga3:xv:theorem:2:1}{2.1}, \SGARefLink{sga3:xv:corollary:2-2}{2.2}, and \SGARefLink{sga3:xv:corollary:2:3}{2.3} in which, instead of lifting
a sub-torus of \(G_0\), one is given a torus \(T\) over \(S\) and asks
to lift a morphism
\[
  u_0:T_0\longrightarrow G_0.
\]
One reduces to the preceding cases by considering the graph of
\(u_0\).}
\SGARefTarget{sga3:xv:remark:2:4}

Let us show how \SGARefLink{sga3:xv:corollary:2-2}{Corollaries 2.2} and \SGARefLink{sga3:xv:corollary:2:3}{2.3} follow from \SGARefLink{sga3:xv:theorem:2:1}{Theorem 2.1}.

\SGAProof{Proof of \SGARefLink{sga3:xv:corollary:2-2}{Corollary 2.2}}
The uniqueness of \(T\) follows from Exposé \SGARefLink{sga3:ix:corollary:4:8:item:b}{IX~4.8(b)} and Exposé
\SGARefLink{sga3:ix:remark:4:10}{IX~4.10}. To prove existence, we may therefore reduce to the case in
which \(S\) is affine, hence noetherian, and \(S_0\) is defined by a
square-zero ideal.

\SGAStatement{Lemma 2.5}{Let \(G\) be an \(S\)-group prescheme of
finite presentation, and let \(H\) be a subgroup scheme of \(G\), of
multiplicative type and smooth over \(S\). Then
\(\operatorname{Centr}_G(H)\) is representable by an \(S\)-group
subprescheme of \(G\), of finite presentation.}
\SGARefTarget{sga3:xv:lemma:2:5}

When \(G\) is separated over \(S\), the lemma follows from Exposé
\SGARefLink{sga3:viii:examples:6:5:item:e:01}{VIII~6.5(e)}, without any smoothness hypothesis on \(H\). In the present
case, observe that the assertion is local for the fpqc topology. We
may therefore assume that \(H\) is diagonalizable, hence of the form
\(D_S(M)\). We may also assume that \(S\) is affine and then, by EGA
IV~8, noetherian. Since \(H\) is smooth and of finite type over \(S\),
the order \(q\) of the torsion subgroup of \(M\) is invertible on
\(S\) (Exposé \SGARefLink{sga3:viii:proposition:2:1:item:e:01}{VIII~2.1(e)}). It is then immediate (cf. Exposé \SGARefLink{sga3:ix:remark:4:10}{IX~4.10})
that the subgroups \({}_nH\), with \(n\) ranging over the powers of
\(q\), are schematically dense in \(H\) (Exposé \SGARefLink{sga3:ix:definition:4:1}{IX~4.1}). But
\({}_nH\) is a completely split covering of \(S\), that is, it is
isomorphic to a finite disjoint sum of copies of \(S\). Consequently
\[
  \operatorname{Centr}_G({}_nH)=Z_n
\]
is representable as the intersection of the centralizers in \(G\) of
the sections of \({}_nH\) over \(S\). It remains to apply the following
lemma.

\SGAStatement{\SGAInlineRefTarget{sga3:xv:lemma:2-5-bis}Lemma 2.5 bis}{Let \(S\) be a noetherian prescheme, let
\(G\) be an \(S\)-group prescheme of finite type, let \(H\) be a
subgroup of \(G\) of multiplicative type, and let \((H_i)\) be an
increasing filtered family of subgroups of \(G\) of multiplicative
type. Suppose that
\[
  Z_i=\operatorname{Centr}_G(H_i)
\]
is representable by a group subprescheme of \(G\). Then the family
\((Z_i)\) is stationary. If, in addition, the \(H_i\) are schematically
dense in \(H\), then
\[
  Z_i=\operatorname{Centr}_G(H)
\]
for all sufficiently large \(i\).}
\SGARefTarget{sga3:xv:lemma:2:5:bis}

To see that the family \((Z_i)\) is stationary, it is enough to prove
that the family of underlying sets \(\operatorname{ens}(Z_i)\) is
stationary. Indeed, the stationary value will then be a closed subset
of an open subset \(U\) of \(G\). After replacing \(G\) by \(U\), we
are reduced to studying a decreasing filtered family of closed
subpreschemes of a noetherian prescheme. An elementary constructibility
argument reduces us to the case in which \(S\) is integral. We then
have to show that \(\operatorname{ens}(Z_i)\) is stationary over some
nonempty open subset of \(S\). The generic fibre of \(G\) is separated
(Exposé VI\textsubscript{A}~0.3); after restricting \(S\), we may
therefore assume that \(G\) is separated over \(S\) (EGA IV~8). Then
\(Z_i\) is closed in \(G\) (Exposé VIII~6.5(e)).

To prove the last assertion of the lemma, let \(Z\) denote the
stationary value of the family \((Z_i)\). It is clear that
\(\operatorname{Centr}_G(H)\) is a subfunctor of \(Z\); we show that
\(Z\) centralizes \(H\). Let \(E\) be the subprescheme of
\(H\times_S Z\) that is the kernel of the pair of morphisms
\[
\begin{aligned}
  H\times_S Z&\rightrightarrows G,\\
  (h,c)&\longmapsto c,\\
  (h,c)&\longmapsto hch^{-1}.
\end{aligned}
\]
The prescheme \(E\) contains \(H_i\times_S Z\) for every \(i\). The
\(H_i\) are flat over \(S\), so EGA IV~11.10.9 shows that, for every
point \(s\) of \(S\), the \((H_i)_s\) are schematically dense in
\(H_s\), and \(H_i\times_S Z\) is schematically dense in
\(H\times_S Z\). Since \(G_s\) is separated, \(E_s\) is closed in
\((H\times_S Z)_s\), and is therefore equal to it. Thus \(E\) is
closed in \(H\times_S Z\), and hence equals \(H\times_S Z\). This says
that \(Z\) centralizes \(H\).

Returning to the proof of \SGARefLink{sga3:xv:corollary:2-2}{2.2}, \SGARefLink{sga3:xv:lemma:2:5}{Lemma 2.5} shows that
\[
  Z_n=\operatorname{Centr}_G(M_n)
\]
is representable, and \SGARefLink{sga3:xv:lemma:2-5-bis}{Lemma 2.5 bis} shows that the decreasing family
of subpreschemes \((Z_n)\) is stationary; let \(Z\) be its stationary
value. The group \(Z\) contains \(T_0\) and all the \(M_n\). After
replacing \(G\) by \(Z\), we may therefore suppose that the \(M_n\)
are central.

The hypotheses of \SGARefLink{sga3:xv:theorem:2:1}{Theorem 2.1} are now satisfied, so there exists a
sub-torus \(T\) of \(G\) lifting \(T_0\). The groups \({}_nT\) and
\(M_n\) are two liftings of \({}_nT_0\); they are therefore conjugate
(Exposé \SGARefLink{sga3:ix:theorem:3:2:bis}{IX~3.2 bis}), and consequently coincide because \(M_n\) is
central. This torus \(T\) has the required properties.

\SGAProof{Proof of \SGARefLink{sga3:xv:corollary:2:3}{Corollary 2.3}}
The uniqueness of \(T\) follows from Exposé \SGARefLink{sga3:ix:theorem:5:1:bis}{IX~5.1 bis}, and its
centrality follows from Exposé \SGARefLink{sga3:ix:theorem:5:6:item:b}{IX~5.6(b)}. By the usual procedure, this
observation reduces us first to the case in which \(S\) is affine
(and hence \(S_0\) is defined by a nilpotent ideal of finite type),
and then to the noetherian case. After restricting \(S\) if necessary,
we may suppose that there exists an integer \(q\) invertible on \(S\).
\SGARefLink{sga3:xv:corollary:2:3}{Corollary 2.3} then follows from \SGARefLink{sga3:xv:corollary:2-2}{2.2} and \SGARefLink{sga3:xv:proposition:1:2}{1.2}.

\SGAStatement{Remark 2.6}{\SGARefLink{sga3:xv:corollary:2:3}{Corollary 2.3} remains true if the torus
\(T_0\) is replaced by a smooth central subgroup of multiplicative
type of \(G_0\).}
\SGARefTarget{sga3:xv:remark:2:6}

% Authority page transition: local p. 7 / combined p. 877 / printed p. 362.
\SGASubsection{2.2}{Proof of \SGARefLink{sga3:xv:theorem:2:1}{Theorem 2.1}}
\SGARefTarget{sga3:xv:subsection:2:2}

\medskip\noindent\textbf{a) Reduction to the case
\(\boldsymbol{T_0=\mathbb G_{m,S_0}}\).}

By \SGARefLink{sga3:xv:proposition:1:1}{1.1}, we may assume that \(M_n\) is a subgroup of multiplicative
type. Using Exposé \SGARefLink{sga3:ix:theorem:3:2:bis}{IX~3.2 bis}, we may suppose that the family
\((M_n)\) is coherent (\SGARefLink{sga3:xv:corollary:2-2}{2.2}). The centralizer \(Z_n\) of \(M_n\) in
\(G\) is representable (\SGARefLink{sga3:xv:lemma:2:5}{2.5}), and the filtered family \((Z_n)\) is
stationary (\SGARefLink{sga3:xv:lemma:2-5-bis}{2.5 bis}). After replacing \(G\) by \(Z_n\) for sufficiently
large \(n\), we may therefore suppose that \(T_0\) and the \(M_n\) are
central. The uniqueness of a lifting of \(T_0\) is then guaranteed
by Exposé \SGARefLink{sga3:ix:theorem:5:1:bis}{IX~5.1 bis}.

Proceeding as in the proof of \SGARefLink{sga3:xv:proposition:1:1}{1.1}, we may suppose that there exist an
\(S\)-torus \(T\) and an \(S_0\)-isomorphism
\[
  u_0:T\times_S S_0\xrightarrow{\sim}T_0.
\]
Lifting \(u_0\) is equivalent to lifting \(T_0\). By uniqueness, it is
enough to prove the existence of a lifting of \(u_0\) after an affine
faithfully flat extension \(S'\to S\) of finite type (fpqc descent).
This allows us to suppose that
\[
  T=\mathbb G_{m,S}^{\,r}
\]
(Exposé \SGARefLink{sga3:x:corollary:4-5}{X~4.5}). If the restriction of \(u_0\) to each factor
\(\mathbb G_m\) lifts to an \(S\)-morphism---necessarily central---then
these liftings immediately give a lifting of \(u_0\). In short, we may
suppose that \(T_0=\mathbb G_{m,S_0}\).

\medskip\noindent\textbf{b) Definition of the obstruction to the
existence of a lifting of \(\boldsymbol{T_0}\).}

By \SGARefLink{sga3:xv:proposition:1:1}{1.1}, to prove \SGARefLink{sga3:xv:theorem:2:1}{Theorem 2.1} it is enough to find a subscheme of \(G\),
flat over \(S\), lifting \(T_0\). We shall see that the obstruction to
the existence of such a lifting can be defined as an element of a
certain \(\operatorname{Ext}^1\) group of sheaves of modules.

Let \(U\) be an open subset of \(G\) in which \(T_0\) is closed, and
again denote by \(U\) (respectively, \(U_0\)) the open subscheme of
\(G\) (respectively, \(G_0\)) having this open subset as its underlying
space. The sheaf \(\mathcal O_{T_0}\), regarded as a sheaf on \(U\), is
therefore a quotient of \(\mathcal O_{U_0}\). Let
\[
  h:\mathcal O_{U_0}\longrightarrow\mathcal O_{T_0}
\]
be the canonical epimorphism.

\SGAStatement{\SGAInlineRefTarget{sga3:xv:lemma:2-7}Lemma 2.7}{The canonical map
\[
  h'=\operatorname{id}_J\otimes h:
  J\otimes_{S_0}\mathcal O_{U_0}
  \longrightarrow
  J\otimes_{S_0}\mathcal O_{T_0}
\]
factors, necessarily uniquely, as \(i\circ j_U\), where \(j_U\) is
the canonical epimorphism
\[
  J\otimes_{S_0}\mathcal O_{U_0}
  \longrightarrow
  J\mathcal O_U\simeq J\mathcal O_{U_0}.
\]}
\SGARefTarget{sga3:xv:lemma:2:7}

We must show that \(h'(K)=0\), where \(K\) is the kernel of \(j_U\).
For every integer \(n\) equal to a power of \(q\), there is an
epimorphism
\[
  h_n:\mathcal O_{T_0}\longrightarrow\mathcal O_{{}_nT_0}.
\]
Since \({}_nT_0\) lifts to a scheme \(M_n\) flat over \(S\), the
canonical morphism
\[
  j_n:J\otimes_{S_0}\mathcal O_{{}_nT_0}
  \xrightarrow{\sim}J\mathcal O_{M_n}
\]
is an isomorphism. The commutative diagram
\[
\begin{tikzcd}[column sep=large,row sep=large]
K \arrow[r]
  & J\otimes_{S_0}\mathcal O_{U_0}
      \arrow[r,"j_U"] \arrow[d,"h'"]
  & J\mathcal O_U \arrow[r,hook] \arrow[dl,"i"]
  & \mathcal O_U \\
  & J\otimes_{S_0}\mathcal O_{T_0}
      \arrow[d,"h'_n"] & & \\
  & J\otimes_{S_0}\mathcal O_{{}_nT_0}
      \arrow[r,"j_n"',"\sim"] & J\mathcal O_{M_n}
      \arrow[r,hook] & \mathcal O_{M_n}
\end{tikzcd}
\]
shows that \(h'(K)\) is contained in \(\ker h'_n\) for every \(n\).
It is therefore contained in \(\bigcap_n\ker h'_n\), and it remains
to prove that this intersection is zero. Now
\[
  \mathcal O_{T_0}=\mathcal O_{S_0}[\mathbb Z],
  \qquad
  \mathcal O_{{}_nT_0}
     =\mathcal O_{S_0}[\mathbb Z/n\mathbb Z].
\]
Let
\[
  a=\sum_{m\in\mathbb Z}a_m\otimes m
  \in J\otimes_{S_0}\mathcal O_{T_0}.
\]
The \(a_m\) are sections of \(J\), almost all zero. Choose \(n\) large
enough that the indices \(m\) for which \(a_m\ne0\) have distinct
images in \(\mathbb Z/n\mathbb Z\). If \(a\in\ker h'_n\), then
necessarily \(a=0\). Thus \(\bigcap_n\ker h'_n=0\), proving \SGARefLink{sga3:xv:lemma:2-7}{2.7}.

% Authority page transition: local p. 9 / combined p. 879 / printed p. 365.
Let \(K_0\) be the kernel of \(h\). Consider the diagram
\[
\begin{tikzcd}[column sep=large,row sep=large]
  &&& 0 \arrow[d] \\
  & J\mathcal O_U \arrow[d,"\sim"] && K_0 \arrow[d] \\
0 \arrow[r]
  & J\mathcal O_{U_0} \arrow[r] \arrow[d,"i"']
  & \mathcal O_U \arrow[r] & \mathcal O_{U_0}
      \arrow[r] \arrow[d,"h"'] & 0 \\
  & J\otimes_{S_0}\mathcal O_{T_0}
  && \mathcal O_{T_0} \arrow[d] \\
  &&& 0 .
\end{tikzcd}
\]
The sheaf \(\mathcal O_U\) defines an element
\[
  \Phi\in
  \operatorname{Ext}^{1}_{\mathcal O_U}
  (\mathcal O_{U_0},J\mathcal O_{U_0}).
\]
Let \(\Psi\) be the element of
\[
  \operatorname{Ext}^{1}_{\mathcal O_U}
  (K_0,J\otimes_{S_0}\mathcal O_{T_0})
\]
obtained from \(\Phi\), by the bifunctoriality of
\(\operatorname{Ext}^{1}\), through the morphisms
\[
  K_0\longrightarrow\mathcal O_{U_0},
  \qquad
  i:J\mathcal O_{U_0}
    \longrightarrow J\otimes_{S_0}\mathcal O_{T_0}.
\]
It follows from Exposé \SGARefLink{sga3:iii:lemma:4-1}{III~4.1} and the infinitesimal criterion of
flatness (cf. Exposé \SGARefLink{sga3:iii:proposition:4-3}{III~4.3}) that a subscheme of \(U\), flat over
\(S\) and lifting \(T_0\), exists if and only if \(\Psi=0\).

Since \(T_0\) is affine, Exposé \SGARefLink{sga3:iii:proposition:4-5}{III~4.5} and \SGARefLink{sga3:xv:lemma:4:6}{4.6} show that it is enough
to construct, locally on \(U\), a subscheme flat over \(S\) lifting
\(T_0\). It therefore suffices to show that the image of \(\Psi\) in
the sheaf
\[
  \mathscr E=
  \mathcal Ext^{1}_{\mathcal O_U}
  (K_0,J\otimes_{S_0}\mathcal O_{T_0})
\]
is zero. We again denote this element of
\(\Gamma(U,\mathscr E)\) by \(\Psi\).

\medskip\noindent\textbf{c) Reduction to the case in which \(S\) is
local artinian with algebraically closed residue field.}

Since \(K_0\) and \(J\otimes_{S_0}\mathcal O_{T_0}\) are coherent
sheaves, so is \(\mathscr E\); moreover, \(\mathscr E\) is supported
on \(T_0\), because the same is true of
\(J\otimes_{S_0}\mathcal O_{T_0}\). To show that the section \(\Psi\)
of \(\mathscr E\) is zero, it is enough to show that, for every point
\(u\) of \(T_0\), its image in the fibre of \(\mathscr E\) at \(u\)
is zero. Formation of the \(\mathcal Ext^i\) of coherent sheaves
commutes with flat extensions of the base,\footnote{Editor's note:
cf. EGA \(0_{\mathrm{III}}\), 12.3.5.} so we are reduced to proving the
existence of a lifting of
\[
  T_0\cap\operatorname{Spec}\mathcal O_u
\]
flat over \(S\). Let \(s\) be the image of \(u\) in \(S\). We may
replace \(S\) by \(\operatorname{Spec}\mathcal O_s\) and \(G\) by
\(G\times_S\operatorname{Spec}\mathcal O_s\).

After another faithfully flat extension if necessary, we may suppose
that \(\mathcal O_s\) has an algebraically closed residue field
(EGA \(0_{\mathrm{III}}\), 10.3.1). Let \(\mathfrak m\) be the maximal
ideal of \(\mathcal O_s\), and set
\[
  S_n=\operatorname{Spec}(\mathcal O_s/\mathfrak m^n).
\]
Suppose that, for every \(n>0\), the obstruction to lifting
\[
  T_0\times_S S_n=(T_0)_n
\]
to a subscheme of \(G_n=G\times_S S_n\), flat over \(S_n\), has been
shown to vanish, and let \(T_n\) be the unique flat lifting of
\((T_0)_n\) over \(S_n\) that is a sub-torus of \(G_n\). If
\(n>n'\), it is clear that
\[
  T_n\times_{S_n}S_{n'}=T_{n'}.
\]
If \(u\) is a point of \(T_0\) lying over \(s\), the lemma below,
applied to the coherent system of liftings
\[
  T_n\cap\operatorname{Spec}\mathcal O_u
\]
of
\((T_0)_n\cap\operatorname{Spec}\mathcal O_u\), shows that
\(T_0\cap\operatorname{Spec}\mathcal O_u\) indeed has a lifting flat
over \(\mathcal O_s\). We are therefore reduced to proving that
\(\Psi=0\) when \(S=S_n\) is the spectrum of an artinian local ring
with algebraically closed residue field, and to proving the following
lemma.

\SGAStatement{\SGAInlineRefTarget{sga3:xv:lemma:2-8}Lemma 2.8}{Let \(A\to B\) be a local homomorphism of
noetherian local rings, let \(\mathfrak m\) be the maximal ideal of
\(A\), let \(J\) be a square-zero ideal of \(A\), let \(M\) be a
finite-type \(B\)-module, and set
\[
  A_0=A/J,\qquad B_0=B/JB,\qquad M_0=M/JM.
\]
Let \(N_0\) be a quotient \(B_0\)-module of \(M_0\), flat over \(A_0\).
For every \(n>0\), let \(A_n,A_{0,n}\), and so forth denote the objects
obtained by the base extension
\[
  A\longrightarrow A/\mathfrak m^n=A_n,
\]
and let \(\overline{J}_n\) be the image of \(J\) in \(A_n\). For every \(n>0\),
let \(N_n\) be a quotient \(B_n\)-module of \(M_n\), flat over \(A_n\),
lifting \(N_{0,n}\), and suppose that, for \(n\ge n'\), \(N_{n'}\) is
obtained from \(N_n\) by the base extension \(A_n\to A_{n'}\). Then
there exists a quotient \(B\)-module \(N\) of \(M\), flat over \(A\),
lifting \(N_0\).}
\SGARefTarget{sga3:xv:lemma:2:8}

\SGAProof{Proof of \SGARefLink{sga3:xv:lemma:2-8}{2.8}}
Let \(P_0\) be the kernel of the epimorphism \(M_0\to N_0\). For every
\(n>0\), there is a commutative diagram with exact horizontal rows:
\[
\tag{$\ast_n$}\label{sga3:xv:lemma:2-8:equation:ast-n}
\begin{tikzcd}[column sep=2em,row sep=large]
&&&& 0 \arrow[d] & \\
&&&& P_{0,n} \arrow[d] & \\
& 0 \arrow[r]
  & \overline{J}_nM_n \arrow[r] \arrow[dd]
  & M_n \arrow[r] \arrow[dd]
  & M_{0,n} \arrow[r] \arrow[dd] & 0 \\
\overline{J}_n\otimes_{A_{0,n}}M_{0,n}
      \arrow[urr] \arrow[drr] &&&&& \\
& 0 \arrow[r]
  & \overline{J}_n\otimes_{A_{0,n}}N_{0,n} \arrow[r]
  & N_n \arrow[r] \arrow[d]
  & N_{0,n} \arrow[r] \arrow[d] & 0 \\
&&& 0 & 0 &
\end{tikzcd}
\]
Moreover, by hypothesis, the diagram \((\ast)_{n+1}\) reduces modulo
\(\mathfrak m^n\) to \((\ast)_n\).

% Authority page transition: local pp. 10--11 / combined pp. 880--881 /
% printed pp. 367--368.
The Artin--Rees lemma (Bourbaki, \emph{Algèbre commutative}, Chapter
III, §3, Corollary 1) shows that the filtrations on \(JM\),
\(J\otimes_{A_0}M_0\), and \(J\otimes_{A_0}N_0\), defined by the
kernels of the respective morphisms
\[
\begin{gathered}
  JM\longrightarrow \overline{J}_nM_n,\\
  J\otimes_{A_0}M_0
    \longrightarrow \overline{J}_n\otimes_{A_{0,n}}M_{0,n},\\
  J\otimes_{A_0}N_0
    \longrightarrow \overline{J}_n\otimes_{A_{0,n}}N_{0,n},
\end{gathered}
\]
are \(\mathfrak mB\)-good. Passing to the inverse limit in the
diagrams \((\ast)_n\), and writing \(\widehat Q\) for the separated
completion of a \(B\)-module \(Q\) for the \(\mathfrak mB\)-adic
topology, we obtain a commutative diagram in which the two horizontal
rows remain exact:
\[
\tag{$\widehat\ast$}\label{sga3:xv:lemma:2-8:equation:widehat-ast}
\begin{tikzcd}[column sep=large,row sep=large]
&&&& 0 \arrow[d] & \\
&&&& \widehat P_0 \arrow[d] & \\
& 0 \arrow[r]
  & \widehat{JM} \arrow[r] \arrow[dd]
  & \widehat M \arrow[r] \arrow[dd]
  & \widehat M_0 \arrow[r] \arrow[dd] & 0 \\
\widehat{J\otimes_{A_0}M_0}
      \arrow[urr] \arrow[drr] &&&&& \\
& 0 \arrow[r]
  & \widehat{J\otimes_{A_0}N_0} \arrow[r]
  & \varprojlim_n N_n \arrow[r] \arrow[d]
  & \widehat N_0 \arrow[r] \arrow[d] & 0 \\
&&& 0 & 0 &
\end{tikzcd}
\]
Now \((B,\mathfrak mB)\) is a Zariski ring and
\(J\otimes_{A_0}N_0\) is of finite type over \(B\), hence separated
for the \(\mathfrak mB\)-adic topology. The diagram
\((\widehat\ast)\) then shows that the morphism
\[
  J\otimes_{A_0}M_0\longrightarrow J\otimes_{A_0}N_0
\]
induced by the epimorphism \(M_0\to N_0\) factors through \(JM\):
\[
\begin{tikzcd}[column sep=large,row sep=large]
J\otimes_{A_0}M_0 \arrow[rr,"\mathrm{can.}"] \arrow[dr]
  && JM \arrow[dl] \\
& J\otimes_{A_0}N_0 &
\end{tikzcd}
\]

Under these conditions, Exposé \SGARefLink{sga3:iii:lemma:4-1}{III~4.1} and \SGARefLink{sga3:iii:proposition:4-3}{4.3} show that a quotient
\(B\)-module \(N\) of \(M\), flat over \(A\) and lifting \(N_0\),
exists if and only if a certain element
\[
  g\in E=\operatorname{Ext}^{1}_B
  (P_0,J\otimes_{A_0}N_0)
\]
vanishes. More precisely, the exact sequence
\[
  0\longrightarrow JM\longrightarrow M
  \longrightarrow M_0\longrightarrow0
\]
defines an element \(f\in F\), where
\[
  F=\operatorname{Ext}^{1}_B(M_0,JM),
\]
and \(g\) is the image of \(f\) under the natural morphism \(F\to E\)
arising by bifunctoriality from
\[
  P_0\longrightarrow M_0,
  \qquad
  JM\longrightarrow J\otimes_{A_0}N_0.
\]
The diagram \((\widehat\ast)\) and Exposé \SGARefLink{sga3:iii:lemma:4-1}{III~4.1} show that the image
of \(g\) in
\[
  \widehat E\simeq
  \operatorname{Ext}^{1}_{\widehat B}
  (\widehat P_0,\widehat{J\otimes_{A_0}N_0})
\]
is zero. But \(E\) is of finite type over \(B\), and hence separated
for the \(\mathfrak mB\)-adic topology. Consequently \(g=0\), which
completes the proof of \SGARefLink{sga3:xv:lemma:2-8}{2.8}.

\medskip\noindent\textbf{d) Study of \(\boldsymbol{\mathscr E}\).}

We now suppose that \(S\) is the spectrum of an artinian local ring
\(A\) with algebraically closed residue field \(k\). Let \(A_0\) be
the ring of \(S_0\), and let \(\mathfrak m_0\) be the maximal ideal
of \(A_0\). Since \(S_0\) is artinian, \(T_0\) is closed in \(G\)
(Exposé VI\textsubscript{B}~\SGARefLink{cor:sga3-VIB-1.4.2}{1.4.2}). We may therefore take \(U=G\), so
that
\[
  \mathscr E=
  \mathcal Ext^{1}_{\mathcal O_G}
  (K_0,J\otimes_{S_0}\mathcal O_{T_0}).
\]
Let \(B_0\) be the ring of the affine \(S_0\)-scheme
\(T_0=\mathbb G_{m,S_0}\), and let \(\overline B_0\) be the ring of its special
fibre
\[
  T_0\times_{S_0}\operatorname{Spec}(k)=\mathbb G_{m,k}.
\]
The sheaf \(\mathscr E\) is a coherent
\(\mathcal O_{T_0}\)-module, hence is defined by a finite-type
\(B_0\)-module, which we denote by \(E\).

Consider the graded module associated with \(E\) and the
\(\mathfrak m_0B_0\)-adic filtration:
\[
  E_r=\mathfrak m_0^rE/\mathfrak m_0^{r+1}E.
\]
Each \(E_r\) is a finite-type \(\overline B_0\)-module, and \(E_r=0\) for all
sufficiently large \(r\), since \(S_0\) is artinian.

Let \(g\) be a section of \(G\) over \(S\) whose restriction to
\(S_0\) is a section \(g_0\) of \(T_0\). Left translation by \(g\)
defines an ``automorphism of the situation'' from the point of view
of the obstruction problem. In particular, \(g\) induces an
\(S_0\)-automorphism of the sheaf \(\mathscr E\) that fixes the
obstruction \(\Psi\). More precisely, \(g\) defines a semilinear
automorphism of the \(B_0\)-module \(E\), relative to the
\(A_0\)-automorphism of \(B_0\) defined by translation by \(g_0\) in
\(T_0\). Reducing modulo \(\mathfrak m_0^{r+1}\), \(g\) therefore
defines a semilinear automorphism of \(E_r\), relative to the
\(k\)-automorphism of \(\overline B_0\) defined by translation by
\[
  g_0\times_{S_0}\operatorname{Spec}(k)
\]
in \(T_0\times_{S_0}\operatorname{Spec}(k)\).

\SGAStatement{Lemma 2.9}{For every integer \(r\ge0\), \(E_r\) is a
locally free \(\overline B_0\)-module.}
\SGARefTarget{sga3:xv:lemma:2:9}

Let \(x\) be a point of \(T_0\), let \(\kappa(x)\) be its residue
field, and let
\[
  (E_r)_x=E_r\otimes_{\overline B_0}\kappa(x)
\]
be the ``fibre'' of \(E_r\) at \(x\). Denote by \(\ell(x)\) the rank
of \((E_r)_x\) over \(\kappa(x)\), and by \(\ell\) the maximum of
\(\ell(x)\) as \(x\) ranges over the points of \(T_0\). Let \(L\) be
the largest closed subscheme of
\[
  \operatorname{Spec}\overline B_0=\mathbb G_{m,k}
\]
over which \(E_r\) is locally free of rank \(\ell\) (TDTE IV, \SGARefLink{sga3:xv:lemma:3:6}{Lemma
3.6}). Let \(\beta\in L(k)\), which exists because \(L\) is of finite
type over the algebraically closed field \(k\), and let
\(\alpha\in\mathbb G_{m,k}(k)\) be a point whose order is a power
\(n\) of \(q\). The point \(\alpha\) is \(k\)-rational. Since, by
hypothesis, \({}_nT_0\) lifts to a subscheme \(M_n\) étale over \(S\),
there is a section \(a\) of \(G\) over \(S\) lifting \(\alpha\) and
whose restriction to \(S_0\) is a section of \(T_0\).

The preceding observations show that the fibres of \(E_r\) at the
points \(\beta\) and \(\alpha\beta\) of \(\mathbb G_{m,k}(k)\) are
\(k\)-isomorphic. The points whose order is a power of \(q\) are dense
in \(\mathbb G_{m,k}\), as are their translates by \(\beta\). Since
\(L\) is closed in \(\mathbb G_{m,k}\) and \(\mathbb G_{m,k}\) is
reduced, \(L=\mathbb G_{m,k}\), and \(E_r\) is locally free of rank
\(\ell\) over \(\overline B_0\).

\medskip\noindent\textbf{e) Completion of the proof of \SGARefLink{sga3:xv:theorem:2:1}{Theorem 2.1}.}

% Authority page transition: local p. 13 / combined p. 883 / printed p. 372.
Let \(K_0(n)\) be the sheaf of ideals of \(\mathcal O_{G_0}\)
defining the closed subscheme \({}_nT_0\). The sheaf
\(K_0\)\footnote{Editor's note: Introduced after \SGARefLink{sga3:xv:lemma:2-7}{Lemma 2.7}.} is
therefore a subsheaf of \(K_0(n)\). For brevity, set
\[
  R=J\otimes_{S_0}\mathcal O_{T_0},
  \qquad
  R(n)=J\otimes_{S_0}\mathcal O_{{}_nT_0}.
\]
The sheaf \(R(n)\) is a quotient of \(R\), and we have the diagram
\[
\begin{tikzcd}[column sep=large,row sep=large]
&& K_0 \arrow[d] \\
&& K_0(n) \arrow[d] \\
0 \arrow[r]
  & J\mathcal O_G \arrow[r] \arrow[d]
  & \mathcal O_G \arrow[r]
  & \mathcal O_{G_0} \arrow[r] & 0 \\
  & R \arrow[d] \\
  & R(n)
\end{tikzcd}
\]
Using the bifunctoriality of
\(\operatorname{Ext}^{1}_{\mathcal O_G}\), we obtain the commutative
diagram
\[
\begin{tikzcd}[column sep=huge,row sep=large]
\mathcal{E}xt^{1}_{\mathcal O_G}
  (\mathcal O_{G_0},J\mathcal O_G)
  \arrow[dr] \\
&
\mathcal{E}xt^{1}_{\mathcal O_G}(K_0(n),R)
  \arrow[r] \arrow[d]
&
\mathcal{E}xt^{1}_{\mathcal O_G}(K_0,R)=\mathscr E
  \arrow[d] \\
&
\mathcal{E}xt^{1}_{\mathcal O_G}(K_0(n),R(n))
  \arrow[r]
&
\mathcal{E}xt^{1}_{\mathcal O_G}(K_0,R(n)).
\end{tikzcd}
\]
Again let \(\Phi\) denote the element of
\(\operatorname{Ext}^{1}_{\mathcal O_G}
(\mathcal O_{G_0},J\mathcal O_G)\) defined by the exact sequence
\[
  0\longrightarrow J\mathcal O_G
  \longrightarrow\mathcal O_G
  \longrightarrow\mathcal O_{G_0}
  \longrightarrow0,
\]
so that \(\Psi\) is the image of \(\Phi\) in \(\mathscr E\). Since
\({}_nT_0\) lifts to a subscheme \(M_n\) of \(G\), flat over \(S\),
the image of \(\Phi\) in
\[
  \operatorname{Ext}^{1}_{\mathcal O_G}(K_0(n),R(n))
\]
is zero (Exposé \SGARefLink{sga3:iii:lemma:4-1}{III~4.1}). A fortiori, the image of \(\Phi\) in
\(\operatorname{Ext}^{1}_{\mathcal O_G}(K_0,R(n))\), which is also
the image of \(\Psi\), is zero.

\SGAStatement{Lemma 2.10}{For every integer \(n\), the canonical
morphism
\[
  \operatorname{Ext}^{1}_{\mathcal O_G}(K_0,R)
    \otimes_{B_0}\mathcal O_{{}_nT_0}
  \longrightarrow
  \operatorname{Ext}^{1}_{\mathcal O_G}(K_0,R(n))
\]
is injective.}
\SGARefTarget{sga3:xv:lemma:2:10}

Indeed, the affine scheme \(T_0=\mathbb G_{m,S_0}\) has ring
\[
  B_0=A_0[T,T^{-1}].
\]
The subscheme \({}_nT_0\) is defined by the vanishing of
\[
  h(n)=T^n-1.
\]
This section is regular (EGA \(0_{\mathrm{IV}}\), 15.2.2) and remains
regular after every base change \(S'_0\to S_0\). Hence there is an
exact sequence of sheaves
\[
  0\longrightarrow\mathcal O_{T_0}
  \xrightarrow{\,h(n)\,}\mathcal O_{T_0}
  \longrightarrow\mathcal O_{{}_nT_0}
  \longrightarrow0.
\]
Since \({}_nT_0\) is flat over \(S\), tensoring over
\(\mathcal O_{S_0}\) with \(J\) gives the exact sequence
\[
  0\longrightarrow R\xrightarrow{\,h(n)\,}R
  \longrightarrow R(n)\longrightarrow0,
\]
and then the exact sequence of Ext groups
\[
  \cdots\longrightarrow
  \operatorname{Ext}^{1}_{\mathcal O_G}(K_0,R)
  \xrightarrow{\,h(n)\,}
  \operatorname{Ext}^{1}_{\mathcal O_G}(K_0,R)
  \longrightarrow
  \operatorname{Ext}^{1}_{\mathcal O_G}(K_0,R(n)),
\]
which proves the lemma.

The foregoing shows that, for every integer \(n\) equal to a power of
\(q\), the image of \(\Psi\) in
\(\mathscr E/h(n)\mathscr E\) is zero. To prove that \(\Psi=0\), it
is enough to show that
\[
  \Psi\in\mathfrak m_0^r\mathscr E
  \quad\Longrightarrow\quad
  \Psi\in\mathfrak m_0^{r+1}\mathscr E.
\]
Let \(\overline\Psi\) be the image of \(\Psi\) in \(E_r\). There is
an element \(\Psi(n)\in\mathscr E\) such that
\[
  \Psi=\Psi(n)h(n)
  \qquad(n\text{ a power of }q).
\]
We observed that the image \(\overline{h(n)}\) of \(h(n)\) in \(\overline B_0\)
is again \(\overline B_0\)-regular. Since \(E_s\) is locally free over \(\overline B_0\)
for every \(s\) (\SGARefLink{sga3:xv:lemma:2:9}{2.9}), multiplication by \(\overline{h(n)}\) on
\(E_s\) is injective. It follows immediately that
\(\Psi(n)\in\mathfrak m_0^r\mathscr E\). Let
\(\overline{\Psi(n)}\) be its image in \(E_r\); then
\[
  \overline\Psi=\overline{h(n)}\,\overline{\Psi(n)}
  \qquad(n\text{ a power of }q).
\]
Thus the set \(F\) of points of \(\mathbb G_{m,k}\) at which
\(\overline\Psi\) vanishes contains the dense set of points whose
order is a power of \(q\). Moreover, \(F\) is closed because \(E_r\)
is locally free over \(\overline B_0\), and \(\mathbb G_{m,k}\) is reduced.
Therefore \(\overline\Psi=0\).

This completes the proof of \SGARefLink{sga3:xv:theorem:2:1}{Theorem 2.1}.

% Exact stop: immediately before §3, “Characterization of a sub-torus
% by its underlying set”, authority local p. 14 / combined p. 884 /
% printed p. 374.

% END inputs/XV/tex/components/02_expose_XV_section1_remainder_and_section2.tex

% BEGIN inputs/XV/tex/components/03_expose_XV_section3.tex
% Authority: Expo15.pdf, local pp. 14--26 / combined pp. 884--896 /
% printed pp. 374--398.
% Begins at §3, “Characterization of a sub-torus by its underlying
% set”, and stops immediately before §4.

\SGASection{3}{Characterization of a sub-torus by its underlying set}
\SGARefTarget{sga3:xv:section:3}

\SGASubsection{3.1}{Statement of the theorem}
\SGARefTarget{sga3:xv:subsection:3:1}

\noindent\textbf{Notation.}---If \(X\) is a prescheme,
\(\operatorname{ens}(X)\) denotes the underlying set of \(X\). If
\(X\) and \(S'\) are two \(S\)-preschemes,
\[
  X_{S'}=X\times_S S',
  \qquad u:X_{S'}\longrightarrow X
\]
is the canonical morphism, and if \(E\) is a subset of
\(\operatorname{ens}(X)\), we denote by \(E_{S'}\) (or
\(E\times_S S'\)) the subset \(u^{-1}(E)\) of
\(\operatorname{ens}(X_{S'})\).

\SGAStatement{\SGAInlineRefTarget{sga3:xv:theorem:3-1}Theorem 3.1}{Let \(S\) be a locally noetherian
prescheme, let \(G\) be an \(S\)-group prescheme of finite type, let
\(q>0\) be an integer invertible on \(S\), and let \(E\) be a subset
of \(\operatorname{ens}(G)\). Consider the following assertions about
\(E\):
\begin{enumerate}
\item[\normalfont\SGAInlineRefTarget{sga3:xv:theorem:3-1:item:i}(i)] The set \(E\) is the underlying set of a
sub-torus \(T\) of \(G\).
\item[\normalfont\SGAInlineRefTarget{sga3:xv:theorem:3-1:item:ii}(ii)]
  \begin{enumerate}
  \item[\normalfont a)] For every point \(s\) of \(S\), there is a
  sub-torus \(T_s\) of \(G_s\) such that
  \[
    E\cap\operatorname{ens}(G_s)=E_s
  \]
  is the underlying set of \(T_s\).
  \item[\normalfont b)] As \(n\) ranges over the powers of \(q\),
  there is a coherent family (cf.\ \SGARefLink{sga3:xv:corollary:2-2}{2.2}) \(M_n\) of subgroup schemes
  of \(G\), of multiplicative type, such that for every point \(s\)
  of \(S\),
  \[
    (M_n)_s={}_nT_s.
  \]
  \end{enumerate}
\item[\normalfont\SGAInlineRefTarget{sga3:xv:theorem:3-1:item:iii}(iii)]
  \begin{enumerate}
  \item[\normalfont a)] As in \textup{\SGARefLink{sga3:xv:theorem:3-1:item:ii}{(ii)} a)}.
  \item[\normalfont b)] The set \(E\) is locally closed in
  \(\operatorname{ens}(G)\), and the dimension of the fibers of \(E\)
  over \(S\) is locally constant.
  \end{enumerate}
\item[\normalfont\SGAInlineRefTarget{sga3:xv:theorem:3-1:item:iv}(iv)]
  \begin{enumerate}
  \item[\normalfont a)] As in \textup{\SGARefLink{sga3:xv:theorem:3-1:item:ii}{(ii)} a)}.
  \item[\normalfont b)] For every \(S\)-scheme \(S'\) that is the
  spectrum of a complete discrete valuation ring with algebraically
  closed residue field, \(E_{S'}\) is the underlying set of a
  sub-torus of \(G_{S'}\).
  \end{enumerate}
\end{enumerate}
Then the following implications hold:
\begin{enumerate}
\item[\normalfont A)] \(\textup{(i)}\Longleftrightarrow
\textup{(ii)}\Longrightarrow\textup{(iii)}
\Longrightarrow\textup{(iv)}\).
\item[\normalfont B)] If \(G\) is separated over \(S\), then
\(\textup{(iii)}\Longleftrightarrow\textup{(iv)}\), and moreover
\(E\) is closed.
\item[\normalfont C)] Conditions \textup{\SGARefLink{sga3:xv:theorem:3-1:item:i}{(i)}}, \textup{\SGARefLink{sga3:xv:theorem:3-1:item:ii}{(ii)}}, and
\textup{\SGARefLink{sga3:xv:theorem:3-1:item:iii}{(iii)}} (and also \textup{\SGARefLink{sga3:xv:theorem:3-1:item:iv}{(iv)}} if \(G\) is separated over
\(S\)) are equivalent in the following two cases:
  \begin{enumerate}
  \item[\normalfont First case:]
    \begin{enumerate}
    \item[\normalfont a)] The prescheme \(S\) is reduced, or \(G\)
    is flat over \(S\); and
    \item[\normalfont b)] for every point \(s\) of \(S\), \(T_s\)
    is a central torus of \(G_s\).
    \end{enumerate}
  \item[\normalfont Second case:] \(S\) is normal.
  \end{enumerate}
\end{enumerate}
Moreover, in the two cases above, the torus \(T\) with underlying
set \(E\) is unique.}
\SGARefTarget{sga3:xv:theorem:3:1}

% Authority page transition: local p. 15 / combined p. 885 /
% printed pp. 375--376.
\SGAStatement{\SGAInlineRefTarget{sga3:xv:remark:3-2}Remarks 3.2}{\SGAInlineRefTarget{sga3:xv:remark:3-2:item:a}\textup{a)} When \(S\) is reduced, there
is no need in \textup{\SGARefLink{sga3:xv:theorem:3-1:item:ii}{(ii)}} to assume that the family \(M_n\) is
coherent, since this condition follows from the other hypotheses.
Indeed, if \(m\) divides \(n\), the subgroups \({}_mM_n\) and \(M_m\)
are étale over \(S\), hence reduced, and have the same underlying
space; they therefore coincide.

\textup{b)} If \(S\) is no longer assumed normal, it is in general no
longer true that \(\textup{(iii)}\Rightarrow\textup{(i)}\), even when
\(S\) is reduced and geometrically unibranch and \(G\) is a smooth
group scheme over \(S\). Indeed, consider the Borel subgroup of
\(\operatorname{SL}_{2,S}\) consisting of matrices of the form
\[
  \begin{pmatrix}t&u\\0&t^{-1}\end{pmatrix},
\]
where \(S\) is the affine curve over a field \(k\) with ring
\[
  k[x,y]/(y^2-x^3).
\]
Define \(E\) as follows. Above the cusp of \(S\), where \(x=y=0\),
take the diagonal torus \(u=0\). Above the complementary open subset
\(x\ne0\), take the torus obtained from the diagonal torus by
conjugation by
\[
  \begin{pmatrix}1&y/x\\0&1\end{pmatrix}.
\]
The resulting set \(E\) satisfies \textup{\SGARefLink{sga3:xv:theorem:3-1:item:iii}{(iii)} a)}; it is also
closed in \(G\), and the reduced subscheme having \(E\) as its
underlying set has equations
\[
  \begin{cases}
    xu+y(t-t^{-1})=0,\\
    u^2-x(t-t^{-1})^2=0.
  \end{cases}
\]
It is therefore not a sub-torus of \(G\), since its fiber above the
cusp is not reduced.}
\SGARefTarget{sga3:xv:remarks:3:2}

\noindent\textbf{Plan of the proof of \SGARefLink{sga3:xv:theorem:3-1}{3.1}.} In a first part we shall
establish the following ``easy'' implications:
\[
\begin{tikzcd}[column sep=huge,row sep=large]
 & \textup{(ii)} \arrow[dr,Rightarrow] & \\
\textup{(i)} \arrow[ur,Rightarrow] \arrow[dr,Rightarrow] && \textup{(iv)} \\
 & \textup{(iii)} \arrow[ur,Rightarrow] &
\end{tikzcd}
\]
and
\[
  \textup{(iv)}
  \Longrightarrow
  [\,\textup{(iii)}\text{ and }E\text{ is closed}\,]
  \quad\text{if \(G\) is separated over \(S\)}.
\]
The more delicate implications will be proved in three stages:
\begin{enumerate}
\item[\normalfont I)] Reduction of
\(\textup{(iii)}\Rightarrow\textup{(i)}\), under the hypotheses of
C), first case, to the case in which \(S\) is normal.
\item[\normalfont II)] \(\textup{(iii)}\Rightarrow\textup{(ii)}\)
when \(S\) is normal.
\item[\normalfont III)] \(\textup{(ii)}\Rightarrow\textup{(i)}\).
\end{enumerate}

\SGASubsection{3.2}{Proof of the ``easy'' assertions in \SGARefLink{sga3:xv:theorem:3-1}{3.1}}
\SGARefTarget{sga3:xv:subsection:3:2}

The implications \(\textup{(i)}\Rightarrow\textup{(ii)}\) and
\(\textup{(i)}\Rightarrow\textup{(iii)}\) are trivial.

\medskip
\noindent\(\textup{(iii)}\Rightarrow\textup{(iv)}\).
After replacing \(S\) by \(S'\), we may assume that \(S\) is the
spectrum of a discrete valuation ring. Let \(t\) be the generic point
of \(S\), and let \(s\) be its closed point.

Since \(E\) is locally closed, there is a reduced subprescheme of
\(G\) with underlying space \(E\); we also denote it by \(E\). Its
generic fiber \(E_t\) is therefore the sub-torus \(T_t\) of \(G_t\),
of rank \(r\), by \textup{\SGARefLink{sga3:xv:theorem:3-1:item:iii}{(iii)} a)}. Let \(E'\) be the schematic
closure of \(E_t\) in \(E\). The prescheme \(E'\) is irreducible, and
its closed fiber \(E'_s\) is nonempty (it contains the unit section
of \(G_s\)); hence \(E'_s\) has dimension \(r\) (EGA IV~14.3.10).
But \(E'_s\) is a closed subscheme of \(E_s\), while \(E_s\) has
dimension \(r\) by \textup{\SGARefLink{sga3:xv:theorem:3-1:item:iii}{(iii)} b)} and is irreducible because it
has the same underlying space as \(T_s\). Thus \(E'_s\) and \(E_s\)
have the same underlying space, and consequently \(E'=E\). This
proves that \(E\) is flat over \(S\).

% Authority page transition: local pp. 16--17 / combined pp. 886--887 /
% printed pp. 377--379.
Let \(E''\) now be the schematic closure of \(E_t\) in \(G\). Then
\(E''\) is a subgroup prescheme of \(G\), flat over \(S\)
(Exposé \SGARefLink{sga3:viii:lemma:7:1}{VIII~7.1}), which contains \(E'\), and therefore \(E\). Its
closed fiber \(E''_s\) is an algebraic subgroup of \(G_s\) of
dimension \(r\) (\emph{loc.\ cit.}). Since \(T_s\) is a closed
irreducible subset of \(G_s\) of dimension \(r\), \(E_s\) has the
same underlying set as the identity component \((E''_s)^0\) of
\(E''_s\). Let \((E'')^0\) denote the ``identity component'' of
\(E''\), that is, the open subgroup of \(E''\) complementary to the
union of the irreducible components of \(E''_s\) that do not contain
the origin. Then
\[
  E=\operatorname{ens}\bigl((E'')^0\bigr).
\]
Since \(E\) and \(E''\) are reduced, one even has
\(E=(E'')^0\). Finally, \(E\) is a subgroup prescheme of \(G\), flat
and of finite type over \(S\), with connected fibers, hence separated
(Exposé VI\textsubscript{B}~\SGARefLink{sga3:xv:theorem:5-2}{5.2}); its generic fiber is the torus
\(T_t\), and its reduced closed fiber is the torus \(T_s\). It follows
that \(E\) is a torus (Exposé \SGARefLink{sga3:x:theorem:8-8}{X~8.8}).

\medskip
\noindent\(\textup{(ii)}\Rightarrow\textup{(iv)}\).
Again we may assume that \(S\) is the spectrum of a discrete valuation
ring, and we retain the notation introduced above. The schematic
closure \(T\) of \(T_t\) in \(G\) is a subgroup prescheme of \(G\),
flat over \(S\), and it contains \(M_n\) for every \(n\) that is a
power of \(q\). Consequently, its closed fiber is a closed subscheme
of \(G_s\) containing \((M_n)_s\) for every \(n\), hence containing
\(T_s\). For dimension reasons, the ``identity component'' \(T^0\)
of \(T\) has \(E\) as its underlying set, and the preceding argument
shows that \(T^0\) is a sub-torus of \(G\).

\medskip
\noindent\(\textup{(iv)}\Rightarrow
[\,\textup{(iii)}\text{ and }E\text{ is closed}\,]\), if \(G\) is
separated over \(S\). Let us show that \(E\) is
closed.\footnote{Editor's note: The implication
\(\textup{(iv)}\Rightarrow\textup{(iii)}\) will appear in the course
of the proof.} We first prove the following lemma.

\SGAStatement{Lemma 3.3}{If the conditions stated in \SGARefLink{sga3:xv:theorem:3-1}{3.1}
\textup{\SGARefLink{sga3:xv:theorem:3-1:item:iv}{(iv)}} are satisfied, \(E\) is a locally constructible subset
of \(\operatorname{ens}(G)\).}
\SGARefTarget{sga3:xv:lemma:3:3}

By the usual method, we are reduced to the case in which \(S\) is
noetherian and integral, with generic point \(\eta\). After shrinking
\(S\), we may assume (Exposé VI\textsubscript{B}~\SGARefLink{par:sga3-VIB-10.10}{10.10}) that there is
a subgroup scheme \(H\) of \(G\), flat over \(S\) and with connected
fibers, whose generic fiber \(H_\eta\) is \(T_\eta\). To prove \SGARefLink{sga3:xv:lemma:3:3}{Lemma
3.3} it is enough to show that \(\operatorname{ens}(H)=E\).

Let \(s\) be a point of \(S\). By EGA II~7.1.9 there is an
\(S\)-scheme \(S'\), the spectrum of a discrete valuation ring which
we may assume complete with algebraically closed residue field, whose
generic point \(t'\) maps to \(\eta\) and whose closed point \(s'\)
maps to \(s\). By \textup{\SGARefLink{sga3:xv:theorem:3-1:item:iv}{(iv)} b)}, there is a sub-torus \(T_{S'}\)
of \(G_{S'}\) having \(E_{S'}\) as its underlying space. The two
subgroup preschemes \(T_{S'}\) and
\(H_{S'}=H\times_S S'\) of \(G_{S'}\) are flat over \(S'\), have
connected fibers, and have the same generic fiber \(T_{t'}\).
Therefore both coincide with the identity component of the schematic
closure of \(T_{t'}\) in \(G_{S'}\). Hence
\(\operatorname{ens}(H_s)=E_s\), and therefore
\(\operatorname{ens}(H)=E\), proving the lemma.

Now that \(E\) is known to be locally constructible, to see that it
is closed it is enough to prove that every specialization in \(G\)
of a point of \(E\) again belongs to \(E\). By the usual technique,
we are reduced to the case in which \(S\) is the spectrum of a
complete discrete valuation ring with algebraically closed residue
field. In that case the sub-torus of \(G\) with underlying space
\(E\), supplied by \textup{\SGARefLink{sga3:xv:theorem:3-1:item:iv}{(iv)} b)}, is closed in \(G\), because
\(G\) is separated (Exposé \SGARefLink{sga3:viii:proposition:7:12}{VIII~7.12}).

\SGASubsection{3.3}{Continuation of the proof of \SGARefLink{sga3:xv:theorem:3-1}{3.1}}
\SGARefTarget{sga3:xv:subsection:3:3}

\noindent\textbf{I) Reduction of
\(\textup{(iii)}\Rightarrow\textup{(i)}\), case C, first case, to
the case in which \(S\) is normal.}

\smallskip
\noindent\textit{a) Reduction to the case in which \(S\) is affine
and reduced.}\quad
We therefore assume that, for every point \(s\) of \(S\), \(E_s\)
is the underlying space of a central sub-torus of \(G_s\). The
uniqueness of \(T\) then follows from Exposé \SGARefLink{sga3:ix:theorem:5:1:bis}{IX~5.1 bis}, and moreover
\(T\) will necessarily be central (\emph{loc.\ cit.}). In view of
this uniqueness, in order to prove the existence of \(T\) we may
suppose that \(S\) is noetherian and affine, with ring \(A\). If
\(S\) is not reduced, \(G\) is flat over \(S\) by hypothesis. By
\SGARefLink{sga3:xv:corollary:2:3}{2.3} it is then enough to solve the problem for \(S_{\mathrm{red}}\)
and \(G\times_S S_{\mathrm{red}}\). We may therefore suppose in
addition that \(S\) is reduced.

\smallskip
\noindent\textit{b) Reduction to the case in which \(A\) is an
algebra of finite type over \(\mathbb Z\).}\quad
We shall use the following lemma.

\SGAStatement{Lemma 3.4}{Let \(k\) be a field, \(G\) an algebraic
\(k\)-group, and \(E\) a subset of \(G\). Let \(k'/k\) be a field
extension, and suppose that \(E_{k'}\) is the underlying set of a
sub-torus \(T_{k'}\) of \(G_{k'}\). If \(T_{k'}\) is central in
\(G_{k'}\), or if \(k\) is perfect, then \(E\) is the underlying set
of a sub-torus \(T\) of \(G\).}
\SGARefTarget{sga3:xv:lemma:3:4}

\SGAProof{Proof}
By fpqc descent it is enough to prove that the two inverse images of
\(T_{k'}\) in \(G_{k''}\), where
\[
  k''=k'\otimes_k k',
\]
coincide. They have the same underlying set. If \(T_{k'}\) is
central, they therefore coincide by Exposé \SGARefLink{sga3:ix:theorem:5:1:bis}{IX~5.1 bis}. If \(k\) is
perfect, \(k''\) is reduced; the two inverse images, being smooth
over \(k''\), are reduced and consequently coincide.

\SGAStatement{\SGAInlineRefTarget{sga3:xv:remark:3-5}Remarks 3.5}{\textup{(1)} \SGARefLink{sga3:xv:lemma:3:4}{Lemma 3.4} applies in
particular when the residue fields of \(S\) are perfect.

\textup{\SGARefLink{sga3:xv:theorem:5:2}{(2)}} Under the hypotheses of the first case in \SGARefLink{sga3:xv:theorem:3-1}{Theorem 3.1}
C, for every \(S\)-scheme \(S'\) as in \SGARefLink{sga3:xv:theorem:3-1}{3.1} \textup{\SGARefLink{sga3:xv:theorem:3-1:item:iv}{(iv)} b)}, the
sub-torus of \(G_{S'}\) having \(E_{S'}\) as its underlying space is
central.}
\SGARefTarget{sga3:xv:remarks:3:5}

\SGAStatement{Lemma 3.6}{Let \(S\) be the projective limit of a
filtered inverse system of affine schemes \(S_i\), and let \(H\) be
an \(S\)-group of multiplicative type and of finite type. There
exist an index \(i\), a group \(H_i\) of multiplicative type over
\(S_i\), and an \(S\)-isomorphism
\[
  H_i\times_{S_i}S\xrightarrow{\sim}H.
\]
If \(H\) is isotrivial, \(H_i\) may be chosen isotrivial.}
\SGARefTarget{sga3:xv:lemma:3:6}

\SGAProof{Proof}
Since \(H\) is of finite presentation over \(S\), there are an index
\(i\), an \(S_i\)-group prescheme \(H_i\), of finite presentation
over \(S_i\), and an \(S\)-isomorphism
\[
  H_i\times_{S_i}S\xrightarrow{\sim}H
\]
(Exposé VI\textsubscript{B}~\SGARefLink{prop:sga3-VIB-10.16}{10.16}). After increasing \(i\), we may
suppose that \(H_i\) is of multiplicative type
(Exposé \SGARefLink{sga3:ix:proposition:2:11}{IX~2.11}). It remains to prove that \(H_i\) can be chosen
quasi-isotrivial, and isotrivial when \(H\) is isotrivial.

We may thus suppose that \(H\) is trivialized by an étale,
surjective morphism \(S'\to S\) of finite presentation. For \(i\)
sufficiently large, there are an \(S_i\)-prescheme \(S'_i\), an
étale, surjective morphism \(S'_i\to S_i\) of finite presentation,
and an \(S\)-isomorphism
\[
  S'_i\times_{S_i}S\xrightarrow{\sim}S'
\]
(EGA IV~17.16). For \(j\) sufficiently large set
\[
  S'_j=S'_i\times_{S_i}S_j,\qquad
  H'_j=H_j\times_{S_j}S'_j,\qquad
  H'=H\times_S S'.
\]
By the choice of \(S'\), there are a finitely generated abelian
group \(M\) and an \(S'\)-isomorphism of group schemes
\[
  D_{S'}(M)\xrightarrow{\sim}H'.
\]
Since the \(S'_i\) are quasi-compact and
\(S'=\varprojlim S'_i\), Exposé VI\textsubscript{B}~10 shows that
there are an index \(j\) and an \(S'_j\)-isomorphism of group
schemes
\[
  D_{S'_j}(M)\xrightarrow{\sim}H'_j.
\]
This says precisely that \(H_j\) is a quasi-isotrivial group of
multiplicative type. When \(H\) is isotrivial, one argues similarly
using a trivialization of \(H\) by a finite étale morphism
\(S'\to S\).

We can now carry out reduction b). The ring \(A\) of \(S\) is the
filtered direct limit of its subrings \(A_i\) of finite type over
\(\mathbb Z\). Put \(S_i=\operatorname{Spec}A_i\), and let
\[
  u_j:S\longrightarrow S_j,\qquad
  u_{jk}:S_k\longrightarrow S_j\quad(k\geq j)
\]
be the transition morphisms. Since \(S\) is noetherian and \(G\) is
of finite presentation over \(S\), there are an index \(i\) and an
\(S_i\)-group prescheme \(G_i\), of finite type over \(S_i\), such
that \(G\) is \(S\)-isomorphic to \(G_i\times_{S_i}S\). Likewise,
since \(E\) is locally closed in \(G\), we may suppose that there is
a locally closed subset \(E_i\) of \(G_i\) such that
\[
  E=E_i\times_{S_i}S
\]
(EGA IV~8.3.11). For \(j>i\), put
\[
  G_j=G_i\times_{S_i}S_j,\qquad
  E_j=E_i\times_{S_i}S_j,
\]
and let \(Q_j\) be the set of points \(s\) of \(S_j\) for which
\((E_j)_s\) is the underlying set of a central sub-torus of
\((G_j)_s\).

It follows from \SGARefLink{sga3:xv:lemma:3:4}{3.4} that
\[
  Q_k=u_{jk}^{-1}(Q_j)\quad(k\geq j),
\]
and by hypothesis
\[
  \operatorname{ens}(S)=u_j^{-1}(Q_j)\quad(j\geq i).
\]
I further claim that \(Q_j\) is ind-constructible
(EGA IV~1.9.4). Indeed, since \(S_j\) is noetherian, it is enough
(EGA IV~1.9.10) to see the following. Let \(S\) be a noetherian
integral scheme with generic point \(\eta\), and suppose that
\(E_\eta\) is the underlying set of a central sub-torus \(T_\eta\)
of \(G_\eta\). Then there is a neighborhood \(U\) of \(\eta\) such
that \(E_s\) has the same property for every \(s\in U\). After
shrinking \(S\), we may suppose that the center \(Z\) of \(G\) is
representable and that there is a subgroup scheme \(T\) of \(G\)
whose generic fiber is \(T_\eta\)
(Exposé VI\textsubscript{B}, \S10). The two locally closed subsets
\(\operatorname{ens}(T)\) and \(E\) of \(\operatorname{ens}(G)\)
coincide on the generic fiber. We can therefore find a neighborhood
\(U\) of \(\eta\) such that, above \(U\),
\[
  \operatorname{ens}(T)=E
\]
(EGA IV~8.3.11). For the same reasons, we may suppose that \(T\) is
central above \(U\); it is a torus by 3.6.

Now that \(Q_j\) is known to be ind-constructible, EGA IV~8.3.4
implies that \(Q_j=\operatorname{ens}(S_j)\) for \(j\) sufficiently
large. We may then replace \(S,G,E\) by \(S_j,G_j,E_j\), which
reduces us to the case in which \(S\) is an affine reduced scheme of
finite type over \(\mathbb Z\).

\smallskip
\noindent\textit{c) Reduction to the case in which \(S\) is the
spectrum of a complete reduced noetherian local ring.}\quad
Because the torus with underlying space \(E\) is unique, the usual
technique (EGA IV~8) and Lemma 3.6 allow us to replace \(S\) by the
spectrum of the local ring \(A\) at a point of \(S\). Let
\(\widehat S\) be the spectrum of the completion \(\widehat A\) of
\(A\) for the topology defined by its maximal ideal. Since \(A\) is
the localization of an algebra of finite type over \(\mathbb Z\),
\(\widehat S\) is still reduced (EGA IV~7.6.5).

I claim that it is enough to solve the problem after the base
extension \(\widehat S\to S\). Indeed, let \(\widehat T\) be the
sub-torus of \(G_{\widehat S}\) with underlying space
\(E_{\widehat S}\). Its two inverse images in
\(G_{\widehat S\times_S\widehat S}\) are central sub-tori with the
same underlying space; they therefore coincide
(Exposé \SGARefLink{sga3:ix:theorem:5:1:bis}{IX~5.1 bis}). By fpqc descent, \(\widehat T\) comes from a
torus \(T\) of \(G\) that solves the problem (cf.\ \SGARefLink{sga3:xv:lemma:3:4}{3.4}).

\smallskip
\noindent\textit{d) A descent lemma.}\quad
Recall the following properties of finite morphisms, noted in
TDTE I. Let \(S,S'\) be two preschemes and
\(u:S'\to S\) a finite morphism. Then:
\begin{enumerate}
\item[\normalfont\SGAInlineRefTarget{sga3:xv:lemma:3:6:item:1}(1)] The morphism \(u\) is an epimorphism if and
only if the canonical morphism of sheaves
\[
  \mathcal O_S\longrightarrow u_*(\mathcal O_{S'})
\]
is injective.
\item[\normalfont(2)] The morphism \(u\) is an effective
epimorphism (Exposé \SGARefLink{sga3:iv:definition:1:3}{IV~1.3}) if and only if the canonical diagram
\[
  0\longrightarrow\mathcal O_S
  \longrightarrow u_*(\mathcal O_{S'})
  \rightrightarrows
  u_*(\mathcal O_{S'})\otimes_{\mathcal O_S}u_*(\mathcal O_{S'})
\]
is exact.
\item[\normalfont\SGAInlineRefTarget{sga3:xv:lemma:3:6:item:3}(3)] If in addition \(S\) is noetherian and \(u\)
is an epimorphism, \(u\) is the composite of a finite sequence of
finite effective epimorphisms.
\end{enumerate}

We can now prove the following lemma, whose proof uses a non-flat
descent technique.

\SGAStatement{\SGAInlineRefTarget{sga3:xv:lemma:3-7}Lemma 3.7}{Let \(S\) be a locally noetherian
prescheme, \(G\) an \(S\)-group prescheme of finite type, \(S'\) a
prescheme, and \(u:S'\to S\) a finite morphism. For every
\(S\)-prescheme \(T\), let \(\mathrm M(T)\) be the set of subgroups
of \(G_T\) of multiplicative type \textup{(respectively, the set of
central subgroups of \(G_T\) of multiplicative type)}, so
that \(\mathrm M\) is naturally a contravariant functor on
\(\operatorname{Sch}/S\). If \(u\) is an effective epimorphism
\textup{(respectively, an epimorphism)}, then the diagram of
sets
\[
  \tag{*}\label{sga3:xv:lemma:3-7:equation:star}
  \mathrm M(S)\longrightarrow\mathrm M(S')
  \rightrightarrows\mathrm M(S'\times_S S')
\]
is exact.}
\SGARefTarget{sga3:xv:lemma:3:7}

\SGAProof{Proof}
\textit{i) Reduction to the case in which \(u\) is an effective
epimorphism.}\quad
Here we consider the functor of central subgroups of \(G\) of
multiplicative type. The injectivity of
\(\mathrm M(S)\to\mathrm M(S')\) is local on \(S\); once this
injectivity is granted, the exactness of \((*)\) is also local on
\(S\). We may therefore suppose that \(S\) is affine and
noetherian.

Consider the case in which \(u:S'\to S\) is the composite of two
finite epimorphisms
\[
  S'\xrightarrow{v}S''\xrightarrow{w}S.
\]
I claim that if both diagrams
\[
  \tag{*'}\label{sga3:xv:lemma:3-7:equation:star-prime}
  \mathrm M(S'')\longrightarrow\mathrm M(S')
  \rightrightarrows\mathrm M(S'\times_{S''}S')
\]
and
\[
  \tag{*''}\label{sga3:xv:lemma:3-7:equation:star-prime-prime}
  \mathrm M(S)\longrightarrow\mathrm M(S'')
  \rightrightarrows\mathrm M(S''\times_S S'')
\]
are exact, then so is \((*)\).

Indeed, the injectivity of
\(\mathrm M(S)\to\mathrm M(S')\) is clear. If \(T'\) belongs to
\[
  \operatorname{Ker}\bigl(
    \mathrm M(S')\rightrightarrows
    \mathrm M(S'\times_S S')\bigr),
\]
then, a fortiori, \(T'\) belongs to
\[
  \operatorname{Ker}\bigl(
    \mathrm M(S')\rightrightarrows
    \mathrm M(S'\times_{S''}S')\bigr).
\]
By the exactness of \((*')\), it therefore comes from a unique
element \(T''\) of \(\mathrm M(S'')\). It is enough to show that
\(T''\) belongs to
\[
  \operatorname{Ker}\bigl(
    \mathrm M(S'')\rightrightarrows
    \mathrm M(S''\times_S S'')\bigr),
\]
for then the exactness of \((*'')\) descends \(T''\) to \(S\).
Let \(T''_1,T''_2\) be the two inverse images of \(T''\) in
\(\mathrm M(S''\times_S S'')\). These are central subgroups of
\(G_{S''\times_S S''}\) of multiplicative type, so to prove that
\(T''_1=T''_2\) it is enough to see that they have the same fibers
(Exposé \SGARefLink{sga3:ix:theorem:5:1:bis}{IX~5.1 bis}). Consider the commutative diagram
\[
\begin{tikzcd}[column sep=huge,row sep=large]
S' \arrow[d,"v"'] &
  S'\times_S S' \arrow[l,shift left=.55ex]
    \arrow[l,shift right=.55ex] \arrow[d,"v\times v"]\\
S'' &
  S''\times_S S'' . \arrow[l,shift left=.55ex]
    \arrow[l,shift right=.55ex]
\end{tikzcd}
\]
The morphism \(v\) is a finite epimorphism, hence dominant by
property \textup{\SGARefLink{sga3:xv:lemma:3:6:item:1}{(1)}} above, and closed; it is therefore surjective,
and so is \(v\times v\). Let \(x''\) be a point of
\(S''\times_S S''\), and let \(x'\) be a point of
\(S'\times_S S'\) lying above \(x''\). The commutativity of the
diagram shows that the two inverse images of
\((T''_1)_{x''}\) and \((T''_2)_{x''}\) in \(G_{x'}\) coincide.
Thus
\[
  (T''_1)_{x''}=(T''_2)_{x''}
\]
(EGA IV~2.2.15).

The preceding argument and an immediate induction on the number of
factors in a factorization of \(u:S'\to S\) into effective
epimorphisms, supplied by property \textup{\SGARefLink{sga3:xv:lemma:3:6:item:3}{(3)}} above, show that in
order to prove the exactness of \((*)\) for central subgroups of
multiplicative type it is enough to treat the case in which \(u\)
is an effective epimorphism. Finally, another application of
Exposé \SGARefLink{sga3:ix:theorem:5:1:bis}{IX~5.1 bis} shows that it is enough to prove \SGARefLink{sga3:xv:lemma:3-7}{3.7} when
\(\mathrm M\) is the functor of all subgroups of \(G\) of
multiplicative type and \(u\) is an effective epimorphism.

\smallskip
\noindent\textit{ii) Injectivity of
\(\mathrm M(S)\to\mathrm M(S')\).}\quad
Since \(u:S'\to S\) is an epimorphism, the canonical morphism
\[
  \mathcal O_S\longrightarrow u_*(\mathcal O_{S'})
\]
is injective. Moreover, an \(S\)-group of multiplicative type is
flat over \(S\). The injectivity of
\(\mathrm M(S)\to\mathrm M(S')\) is therefore a consequence of the
following more general lemma.

\SGAStatement{\SGAInlineRefTarget{sga3:xv:lemma:3-8}Lemma 3.8}{Let \(f:X\to S\) and \(g:S'\to S\) be
morphisms of preschemes, let \(\mathcal F\) be a quasi-coherent
\(\mathcal O_X\)-module, and put
\[
  X'=X\times_S S',
  \qquad
  \mathcal F'=\mathcal F\otimes_{\mathcal O_S}\mathcal O_{S'}.
\]
Let \(\mathrm Q(\mathcal F)\) be the set of quotient
\(\mathcal O_X\)-modules \(\mathcal G\) of \(\mathcal F\) that are
quasi-coherent and flat over \(S\), and define
\(\mathrm Q(\mathcal F')\) analogously relative to
\(\mathcal F',X'\), and \(S'\). Suppose that \(g\) is quasi-compact
and that
\[
  \mathcal O_S\longrightarrow g_*(\mathcal O_{S'})
\]
is injective. Then the canonical map
\[
  \begin{aligned}
  \mathrm Q(\mathcal F)&\longrightarrow\mathrm Q(\mathcal F'),\\
  \mathcal G&\longmapsto
    \mathcal G\otimes_{\mathcal O_S}\mathcal O_{S'}
  \end{aligned}
\]
is injective.}
\SGARefTarget{sga3:xv:lemma:3:8}

\SGAProof{Proof}
Indeed, we may suppose that \(S\) is affine, and then that \(X\) is
affine. Since \(S'\) is quasi-compact, it has a finite covering by
affine open subsets. Replacing \(S'\) by the disjoint union of these
open subsets, we may therefore suppose that \(S'\) is affine. We are
then reduced to the following lemma, whose proof is immediate.

\SGAStatement{Lemma 3.9}{Let \(A\to A'\) be an injective
homomorphism of rings, let \(M\) be an \(A\)-module, let \(N=M/P\)
be a quotient \(A\)-module flat over \(A\), and put
\[
  M'=M\otimes_A A',
  \qquad
  N'=N\otimes_A A'=M'/P'.
\]
Then \(P\) is the inverse image of \(P'\) under the canonical
homomorphism \(M\to M'\). Thus \(N\) and \(P\) are determined by
\(N'\) and \(P'\).}
\SGARefTarget{sga3:xv:lemma:3:9}

\smallskip
\noindent\textit{iii) Exactness of \SGARefLink{sga3:xv:lemma:3-7}{3.7} \((*)\) at
\(\mathrm M(S')\).}\quad
Let
\[
  T'\in\operatorname{Ker}\bigl(
    \mathrm M(S')\rightrightarrows
    \mathrm M(S'\times_S S')\bigr).
\]
Suppose that iii) has been proved when \(S\) is the spectrum of a
noetherian local ring. For every point \(s\) of \(S\), there then
exists a subgroup \(T_s\) of multiplicative type in
\(G\times_S\operatorname{Spec}(\mathcal O_{S,s})\), obtained by
descent from
\[
  T'\times_{S'}u^{-1}\!\bigl(\operatorname{Spec}(\mathcal O_{S,s})\bigr).
\]
By Exposé VI\textsubscript{B}, \S10, and \SGARefLink{sga3:xv:lemma:3:6}{3.6}, there is an open
neighborhood \(U\) of \(s\) in \(S\) and a subgroup scheme \(T\) of
\(G\) over \(U\), of multiplicative type, that extends \(T_s\).
Let \(U'\) be the inverse image of \(U\) in \(S'\). We then have two
subschemes of \(G'|_{U'}\), namely \(T'|_{U'}\) and
\(T\times_U U'\), which coincide over
\(u^{-1}(\operatorname{Spec}(\mathcal O_{S,s}))\).

Regarding the latter scheme as the projective limit of the
\(u^{-1}(V)\), where \(V\) runs through the open neighborhoods of
\(s\) in \(U\), EGA IV~8 shows that there is an open subset
\(V\subset U\), containing \(s\), over whose inverse image
\(T\times_U U'\) and \(T'\) coincide. Thus, under the stated
hypotheses, \(T'\) descends locally on \(S\); by the uniqueness
proved in ii), it consequently descends globally. In short, it is
enough to prove iii) when \(S\) is the spectrum of a noetherian
local ring.

Let \(\widehat S\) be the spectrum of the completion of the ring of
\(S\), and put
\[
  S''=S'\times_S S',
  \qquad
  \widehat S'=\widehat S\times_S S',
  \qquad
  \widehat S''=\widehat S\times_S S''
      =\widehat S'\times_S S'.
\]
I claim that it is enough to show that
\[
  \tag{\(\widehat *\)}\label{sga3:xv:lemma:3-9:equation:widehat-star}
  \mathrm M(\widehat S)\longrightarrow
  \mathrm M(\widehat S')
  \rightrightarrows
  \mathrm M(\widehat S'')
\]
is exact at \(\mathrm M(\widehat S')\). This follows from the
commutative diagram
\[
\begin{tikzcd}[column sep=large,row sep=large]
\mathrm M(S) \arrow[r] \arrow[d]
  & \mathrm M(S') \arrow[r,shift left=.55ex]
      \arrow[r,shift right=.55ex] \arrow[d]
  & \mathrm M(S'') \arrow[d]\\
\mathrm M(\widehat S) \arrow[r]
      \arrow[d,shift left=.55ex] \arrow[d,shift right=.55ex]
  & \mathrm M(\widehat S') \arrow[r,shift left=.55ex]
      \arrow[r,shift right=.55ex]
      \arrow[d,shift left=.55ex] \arrow[d,shift right=.55ex]
  & \mathrm M(\widehat S'')\\
\mathrm M(\widehat S\times_S\widehat S)
  \arrow[r,"f"]
  & \mathrm M(\widehat S'\times_{S'}\widehat S').
\end{tikzcd}
\]
Here the second row is exact at \(\mathrm M(\widehat S')\) by
hypothesis, the first two columns are exact by fpqc descent, and
\(f\) is injective by ii), applied to the finite epimorphism
\[
  \widehat S'\times_{S'}\widehat S'
  \longrightarrow
  \widehat S\times_S\widehat S
\]
deduced from \(S'\to S\) by the flat base change
\(\widehat S\times_S\widehat S\to S\). The diagram chase is left
to the reader. The characterization of finite effective
epimorphisms recalled in d)\SGARefLink{sga3:xv:theorem:5:2}{(2)} shows that
\(\widehat S'\to\widehat S\), obtained from \(S'\to S\) by the
flat base change \(\widehat S\to S\), is again a finite effective
epimorphism. We are therefore reduced to the case in which \(S\)
is, in addition, the spectrum of a complete noetherian local ring.

Let \(S_0\) be the reduced closed subscheme of \(S\) whose
underlying space is the closed point. Let \(T'\) be as above and
write \(T''\) for its common image in \(\mathrm M(S'')\). Since
\[
  S'_0=S'\times_S S_0
\]
is faithfully flat over \(S_0\), there is a subgroup \(T_0\) of
\(G_{S_0}\), of multiplicative type, whose inverse image in
\(G_{S'_0}\) is \(T'_{S'_0}\), by fpqc descent. Since \(S\) is
complete noetherian local, there exist an \(S\)-group \(T\) of
multiplicative type and an \(S_0\)-morphism
\[
  u_0:T\times_S S_0\longrightarrow T_0
\]
(Exposé \SGARefLink{sga3:x:corollary:3:3}{X~3.3}). The inverse image of \(u_0\) over \(S'_0\) extends
uniquely to an \(S'\)-isomorphism
\[
  u':T_{S'}\xrightarrow{\sim}T',
\]
again by Exposé \SGARefLink{sga3:x:corollary:3:3}{X~3.3}. Notice here that, because \(S'\) is finite
over the complete local scheme \(S\), it is a finite disjoint sum
of complete local schemes.

The two inverse images of \(u'\) over \(S''\) are morphisms from
\(T_{S''}\) to \(T''\) that coincide over \(S_0\times_S S''\);
hence they coincide everywhere by the same result. Since \(T\) is
flat over \(S\) and \(S'\to S\) is a finite effective
epimorphism, TDTE I, page~8, gives an exact diagram
\[
  \operatorname{Hom}_S(T,G)\longrightarrow
  \operatorname{Hom}_{S'}(T_{S'},G_{S'})
  \rightrightarrows
  \operatorname{Hom}_{S''}(T_{S''},G_{S''}).
\]
Thus \(u'\) comes from a morphism of preschemes \(u:T\to G\).
Likewise, \(T\times_S T\) is flat over \(S\), and therefore
\[
  \operatorname{Hom}_S(T\times_S T,G)\longrightarrow
  \operatorname{Hom}_{S'}(T_{S'}\times_{S'}T_{S'},G_{S'})
\]
is injective. It follows at once that \(u\) is a homomorphism of
groups, since \(u'\) is one.

Moreover, \(u\) is a monomorphism. First observe that
\(\operatorname{Ker}(u)\) is flat over \(S\). To prove this we may
suppose that \(S\) is artinian local (EGA
\(0_{\mathrm{III}}\), 10.2.2), hence that \(G\) is separated
(Exposé VI\textsubscript{A}, \SGARefLink{sga3:VIA:0.3}{0.3}); then
\(\operatorname{Ker}(u)\) is of multiplicative type (Exposé
\SGARefLink{sga3:ix:theorem:6:8}{IX~6.8}), and is therefore flat over \(S\). Since
\[
  \operatorname{Ker}(u)\times_S S'
    =\operatorname{Ker}(u')=0,
\]
\SGARefLink{sga3:xv:lemma:3-8}{Lemma~3.8} gives \(\operatorname{Ker}(u)=0\). Finally, a
monomorphism \(u\) is an immersion (\SGARefLink{sga3:viii:remark:7:13}{Exposé VIII, Remarks~7.13}),
and its image \(u(T)\) is an element of \(\mathrm M(S)\) whose
image in \(\mathrm M(S')\) is \(T'\). This proves \SGARefLink{sga3:xv:lemma:3-7}{Lemma~3.7}.

\smallskip
\noindent\textit{e) End of the proof of I).}\quad
By reduction c), we are reduced to the case in which \(S\) is the
spectrum of a complete reduced noetherian local ring \(A\). Let
\(S'\) be the spectrum of the normalization \(A'\) of \(A\).
Nagata's theorem says that \(A'\) is finite over \(A\) (EGA
\(0_{\mathrm{IV}}\), 23.1.5), and \(S'\to S\) is an epimorphism
because \(A\to A'\) is injective. Suppose that \(G_{S'}\) has a
torus \(T'\) whose underlying space is \(E_{S'}\). The two inverse
images of \(T'\) in \(G_{S'\times_S S'}\) are central sub-tori
with the same underlying space, so they coincide (Exposé IX~5
bis). By \SGARefLink{sga3:xv:lemma:3-7}{Lemma~3.7}, \(T'\) therefore comes from a central
sub-torus \(T\) of \(G_S\), which plainly has \(E\) as its
underlying space. It is thus enough to prove the existence of
\(T'\). This reduces us to the normal case and completes the proof
of I).

\medskip
\noindent\textit{II) Proof that \textup{\SGARefLink{sga3:xv:theorem:3-1:item:iii}{(iii)}} implies
\textup{\SGARefLink{sga3:xv:theorem:3-1:item:ii}{(ii)}} when \(S\) is normal.}\quad
We may restrict ourselves to the case in which \(S\) is integral;
let \(t\) be its generic point. For every integer \(n\) that is a
power of \(q\), let \({}_nG\) be the ``kernel'' in \(G\) of
raising to the \(n\)-th power. Since \(E\) is locally closed in
\(G\), by \SGARefLink{sga3:xv:theorem:3-1:item:iii}{(iii)}(b), the intersection of
\(\operatorname{ens}({}_nG)\) with \(E\) is locally closed in
\(\operatorname{ens}(G)\). Let \(E(n)\) denote the reduced
subscheme of \(G\) whose underlying space is
\(\operatorname{ens}({}_nG)\cap E\).

We show that the structural morphism
\[
  E(n)\longrightarrow S
\]
is separated and universally open. Both properties have a
valuative criterion (EGA II~7.2.3 and EGA IV~14.5.8). Let
\(S'\) therefore be an \(S\)-scheme that is the spectrum of a
complete discrete valuation ring with algebraically closed residue
field. We have proved that \SGARefLink{sga3:xv:theorem:3-1:item:iii}{3.1(iii)} implies \SGARefLink{sga3:xv:theorem:3-1:item:iv}{3.1(iv)}, so there is a
sub-torus \(T'\) of \(G_{S'}\) whose underlying space is
\(E_{S'}\). Now \({}_nT'\) is finite étale over \(S'\), hence is
separated and universally open over \(S'\); the same is therefore
true of \(E(n)\times_S S'\), which has the same underlying space.

The fibers of \(E(n)\), moreover, all have the same number of
geometric points, namely \(n^r\) if \(r\) is the rank of the torus
\(T_t\).\footnote{The French edition prints \(r^n\); the intended
cardinality of the \(n\)-torsion of a rank-\(r\) torus is \(n^r\).}
Finally, the generic fiber \(E(n)_t\), being reduced, is equal to
\({}_nT_t\), and hence is étale over \(\operatorname{Spec}(t)\).
Since \(S\) is normal, SGA I~10.11 now shows that \(E(n)\) is an
étale covering of \(S\).

If \(s\) is a point of \(S\), then \(E(n)_s\) is étale over
\(\operatorname{Spec}\kappa(s)\), hence reduced, and consequently
coincides with the group \({}_nT_s\) of multiplicative type. We
show that \(E(n)\) is a group subscheme of \(G\). Let
\[
  m:E(n)\times_S E(n)\longrightarrow G
\]
be the morphism induced by multiplication in \(G\). The underlying
map \(\operatorname{ens}(m)\) factors through \(E(n)\), so \(m\)
itself factors through the scheme \(E(n)\), because its source
\(E(n)\times_S E(n)\) is étale over \(S\), and hence reduced. It
then follows from Exposé \SGARefLink{sga3:x:corollary:4-8:item:a}{X~4.8(a)} that \(E(n)\) is a subgroup of
multiplicative type. As observed already in \SGARefLink{sga3:xv:remark:3-2:item:a}{3.2(a)}, the family of
subgroups \(E(n)\) is necessarily coherent. This proves that
\SGARefLink{sga3:xv:theorem:3-1:item:iii}{(iii)} implies \SGARefLink{sga3:xv:theorem:3-1:item:ii}{(ii)} when \(S\) is normal.

\medskip
\noindent\textit{III) Proof that \textup{\SGARefLink{sga3:xv:theorem:3-1:item:ii}{(ii)}} implies
\textup{\SGARefLink{sga3:xv:theorem:3-1:item:i}{(i)}}.}\quad
In fact, we shall show that there is a unique sub-torus \(T\) of
\(G\), with underlying space \(E\), such that
\[
  {}_nT=M_n
\]
for every \(n\) that is a power of \(q\). Uniqueness follows
immediately from Exposé \SGARefLink{sga3:ix:corollary:4:8:item:b}{IX~4.8(b)}. In proving existence, in view of
this uniqueness, we may successively suppose that:
\begin{enumerate}
\item[\normalfont a)] \(S\) is noetherian.
\item[\normalfont b)] The \(M_n\) are central subgroups. Indeed, it
is enough to replace \(G\) by
\[
  Z_n=\operatorname{Centr}_G(M_n)
\]
for sufficiently large \(n\), by \SGARefLink{sga3:xv:lemma:2:5}{2.5} and \SGARefLink{sga3:xv:lemma:2-5-bis}{2.5 bis}.
\item[\normalfont c)] \(S\) is the spectrum of a local ring.
Suppose that the problem has been solved after every base change
\(\operatorname{Spec}(\mathcal O_{S,s})\to S\), for \(s\in S\),
and let \(T_s\) be the resulting sub-torus of
\(G\times_S\operatorname{Spec}(\mathcal O_{S,s})\). There is then
an open neighborhood \(U\) of \(s\) and a sub-torus \(T\) of
\(G|_U\) extending \(T_s\), by Exposé VI\textsubscript{B}, \S10,
and \SGARefLink{sga3:xv:lemma:3:6}{3.6}. We have proved that \SGARefLink{sga3:xv:theorem:3-1:item:ii}{3.1(ii)} implies \SGARefLink{sga3:xv:theorem:3-1:item:iv}{3.1(iv)}; since \(S\)
is noetherian, \(E\) is therefore constructible by \SGARefLink{sga3:xv:lemma:3:3}{3.3}. After
shrinking \(U\), we may consequently suppose that
\[
  \operatorname{ens}(T)=E\times_S U
\]
(EGA IV~9.5.2). For every \(n\) that is a power of \(q\), the
groups \({}_nT\) and \(M_n\times_S U\) are then two subgroups of
multiplicative type in \(G|_U\), with the same fibers; because
\(M_n\) is central, they coincide by Exposé \SGARefLink{sga3:ix:corollary:5:3:bis}{IX~5.3 bis}. Thus the
problem has a local solution on \(S\), and uniqueness makes these
local solutions into a global one.
\item[\normalfont d)] \(S\) is the spectrum of a complete
noetherian local ring and, if \(s\) is the closed point, \(T_s\)
is trivial. This follows from EGA \(0_{\mathrm{III}}\), 10.3.1,
and fpqc descent.
\item[\normalfont e)] \(S\) is reduced. Apply \SGARefLink{sga3:xv:corollary:2-2}{2.2}.
\item[\normalfont f)] \(S\) is normal. Apply Nagata's finiteness
theorem (EGA \(0_{\mathrm{IV}}\), 23.1.5) and Lemma~3.7 for
central sub-tori.
\end{enumerate}

After these reductions, \(T_s\) is trivial, and hence
\((M_n)_s\) is trivial. It follows from Exposé \SGARefLink{sga3:x:corollary:3:3}{X~3.3} that \(M_n\)
itself is trivial. If \(t\) is any point of \(S\), all the
subgroups \({}_nT_t\) of \(T_t\), for \(n\) a power of \(q\), are
therefore trivial; Exposé \SGARefLink{sga3:x:proposition:1:4}{X~1.4} then shows at once that \(T_t\)
itself is trivial. It remains only to prove the following lemma.

\SGAStatement{Lemma 3.10}{Under the hypotheses of \SGARefLink{sga3:xv:theorem:3-1:item:ii}{3.1(ii)}, suppose
in addition that \(S\) is noetherian and normal and that, for every
point \(s\) of \(S\), \(T_s\) is a trivial torus. Then there exists
a unique sub-torus \(T\) of \(G\), with underlying space \(E\),
such that
\[
  {}_nT=M_n
\]
for every \(n\) that is a power of \(q\). Moreover, \(T\) is
trivial.}
\SGARefTarget{sga3:xv:lemma:3:10}

\SGAProof{Proof}
Uniqueness follows from the fact that \(T\), being smooth over
\(S\), is reduced. To prove existence, we may suppose that \(S\)
is irreducible, with generic point \(\eta\). Let \(r\) be the rank
of \(T_\eta\), put
\[
  T'=\mathbb G_{m,S}^{\,r},
\]
and choose an isomorphism
\[
  u_\eta:T'_\eta\xrightarrow{\sim}T_\eta.
\]
Write \({}_nu_\eta\) for its restriction to \({}_nT'_\eta\).
Since \({}_nT'\) and \(M_n\) are trivial, \({}_nu_\eta\) has a
unique extension to an \(S\)-isomorphism
\[
  {}_nu:{}_nT'\xrightarrow{\sim}M_n.
\]
I claim that, for every point \(s\) of \(S\), there is a
necessarily unique group isomorphism
\[
  u_s:T'_s\xrightarrow{\sim}T_s
\]
that extends \(({}_nu)_s\) for every \(n\) that is a power of
\(q\).

Indeed, let \(S^1\) be an \(S\)-scheme that is the spectrum of a
complete discrete valuation ring with algebraically closed residue
field, whose generic point \(t_1\) maps to \(\eta\) and whose
closed point \(s_1\) maps to \(s\) (EGA II~7.1.9). The proof that
3.1(ii) implies 3.1(iv) gives a sub-torus \(T^1\) of \(G_{S^1}\)
such that
\[
  (M_n)_{S^1}={}_nT^1
\]
for every \(n\) that is a power of \(q\). Since \(S^1\) is normal,
\(T^1\) is isotrivial by Exposé \SGARefLink{sga3:x:theorem:5-16}{X~5.16}. The classification of
isotrivial tori (Exposé \SGARefLink{sga3:x:corollary:1:2}{X~1.2}) and SGA I, Exposé V~8.2, together
with the fact that the generic fiber of \(T^1\) is trivial, show
that \(T^1\) is trivial. We may therefore extend
\[
  u_\eta\times_\eta t_1
\]
to an \(S^1\)-isomorphism of group schemes
\[
  u^1:T'\times_S S^1\xrightarrow{\sim}T^1.
\]
On \({}_nT'\times_S S^1\), the morphisms \(u^1\) and
\({}_nu\times_S S^1\) agree on the generic fiber, and hence agree
everywhere. Restricting \(u^1\) to the closed fiber gives the
required \(u_s\) after the field extension
\(\kappa(s)\to\kappa(s_1)\); it descends to \(\kappa(s)\) by
Exposé \SGARefLink{sga3:ix:corollary:4:8:item:a}{IX~4.8(a)} and fpqc descent.

The scheme \(T'\) is smooth over the normal scheme \(S\), hence is
normal. The extension criterion for rational maps in EGA
IV\textsubscript{4}, 20.4.6, now shows that there is an
\(S\)-morphism
\[
  u:T'\longrightarrow G
\]
whose restriction to each fiber \(T'_s\) is the composite
\[
  T'_s\xrightarrow{u_s}T_s\longrightarrow G_s.
\]
Indeed, for every base change \(S^1\to S\) of the kind used above,
the preceding construction gives precisely the required extension
over \(S^1\).

It remains to show that \(u\) is a homomorphism. Let \(m_{T'}\)
and \(m_G\) denote multiplication in \(T'\) and \(G\),
respectively. The coincidence subscheme of
\[
  m_G\circ(u\times_Su)
  \quad\text{and}\quad
  u\circ m_{T'}
\]
contains every fiber, since every \(u_s\) is a homomorphism. It
therefore has the same underlying space as \(T'\times_S T'\), and
so equals \(T'\times_S T'\), because the latter is smooth over
\(S\) and hence reduced. Finally, \(u\) is a monomorphism because
it is one on every fiber; it is therefore an immersion (Exposé
\SGARefLink{sga3:viii:proposition:7:9}{VIII~7.9}). The image of \(T'\) under \(u\) is a sub-torus of \(G\)
whose underlying set is \(E\).

This completes the proof of \SGARefLink{sga3:xv:theorem:3-1}{Theorem~3.1}.

% END inputs/XV/tex/components/03_expose_XV_section3.tex

% BEGIN inputs/XV/tex/components/04_expose_XV_section4.tex
% Authority: Expo15.pdf, local pp. 26--31 / combined pp. 896--901 /
% printed pp. 398--408.
% Stops immediately before §5.

\SGASection{4}{Characterization of a sub-torus \(T\) by the subgroups
\({}_nT\)}
\SGARefTarget{sga3:xv:section:4}

\SGASubsection{4.1}{Statement of the main theorem}
\SGARefTarget{sga3:xv:subsection:4:1}

\SGAStatement{\SGAInlineRefTarget{sga3:xv:theorem:4-1}Theorem 4.1}{Let \(S\) be a connected locally
noetherian prescheme, let \(G\) be an \(S\)-group prescheme of finite
type over \(S\), let \(q>1\) be an integer invertible on \(S\), and
let \(r\) be a positive integer. For every integer \(n\) that is a
power of \(q\), let \(M(n)\) be a subgroup scheme of \(G\), of
multiplicative type and of type
\((\mathbb Z/n\mathbb Z)^r\). Suppose that:
\begin{enumerate}
\renewcommand{\labelenumi}{\textup{\alph{enumi})}}
\item the family \(M(n)\) is coherent: if \(m\mid n\), then
\[
  {}_mM(n)=M(m);
\]
\item\SGAInlineRefTarget{sga3:xv:theorem:4-1:item:b} there exist a point \(s\in S\) and a sub-torus \(T_s\) of
\(G_s\) such that
\[
  M(n)_s={}_nT_s\qquad\text{for every }n;
\]
\item for every point \(t\in S\), there exists a closed affine
subscheme \(F_t\) of \(G_t\) containing \(M(n)_t\) for every \(n\).
\end{enumerate}
Then there exists a unique sub-torus \(T\) of \(G\) such that
\({}_nT=M(n)\) for every \(n\) that is a power of \(q\).}
\SGARefTarget{sga3:xv:theorem:4:1}

There is an analogous theorem for lifting morphisms.

\SGAStatement{\SGAInlineRefTarget{sga3:xv:theorem:4-1-bis}Theorem 4.1 bis}{Let \(S,G,q\) be as above, let \(T\)
be an \(S\)-torus, and, for every \(n\) that is a power of \(q\), let
\(u(n):{}_nT\to G\) be an \(S\)-homomorphism. Suppose that:
\begin{enumerate}
\renewcommand{\labelenumi}{\textup{\alph{enumi})}}
\item the family \(u(n)\) is coherent: if \(m\mid n\), then
\[
  u(m)=u(n)|_{{}_mT};
\]
\item there exist a point \(s\in S\) and a homomorphism
\(u_s:T_s\to G_s\) such that
\[
  u_s|_{{}_nT_s}=u(n)_s
  \qquad\text{for every }n\text{ that is a power of }q;
\]
\item for every point \(t\in S\), there exists a closed affine
subscheme \(F_t\) of \(G_t\) containing
\(u(n)_t({}_nT_t)\) for every \(n\).
\end{enumerate}
Then there exists a unique homomorphism \(u:T\to G\) whose restriction
to \({}_nT\) equals \(u(n)\) for every \(n\) that is a power of \(q\).}
\SGARefTarget{sga3:xv:theorem:4:1:bis}

\SGAStatement{Remark 4.2}{Using the lower semicontinuity of the
abelian rank of a flat group prescheme of finite type over the
spectrum of a discrete valuation ring (cf. Exposé \SGARefLink{sga3:x:remark:8-7}{X~8.7}), condition
(c) in \SGARefLink{sga3:xv:theorem:4-1}{Theorems 4.1} and \SGARefLink{sga3:xv:theorem:4-1-bis}{4.1 bis} may be weakened by requiring the
closed affine subscheme \(F_t\) only for every maximal point \(t\) of
\(S\).}
\SGARefTarget{sga3:xv:remark:4:2}

Let us deduce \SGARefLink{sga3:xv:theorem:4-1-bis}{Theorem 4.1 bis} from \SGARefLink{sga3:xv:theorem:4-1}{Theorem 4.1}. Put
\(G'=G\times_S T\). For every \(n\) that is a power of \(q\), consider
the homomorphism
\[
  v(n):{}_nT\longrightarrow G'
\]
whose projections to \(G\) and \(T\) are respectively \(u(n)\) and
the canonical immersion \({}_nT\to T\). Thus \(v(n)\) is an immersion;
let \(M(n)\) be its image subgroup. The family \(M(n)\) is coherent
in the sense of \SGARefLink{sga3:xv:theorem:4-1}{Theorem 4.1}. The group \(M(n)_s\) equals
\({}_nT'_s\), where \(T'_s\) is the sub-torus of \(G'_s\) that is the
graph of \(u_s\), and, for every \(t\in S\), the closed affine
subscheme \(F_t\times_tT_t\) of \(G'_t\) contains all the groups
\(M(n)_t\). By \SGARefLink{sga3:xv:theorem:4-1}{Theorem 4.1} there is therefore a sub-torus \(T'\) of
\(G'\) such that \({}_nT'=M(n)\) for every \(n\).

Let \(f:T'\to T\) be the restriction of the projection \(G'\to T\),
and let \(f(n)\) be its restriction to \(M(n)\). Then
\[
  f(n)\circ v(n)=\operatorname{id}_{{}_nT}.
\]
The fibre \(T'_s\) is the graph of \(u_s\) (Exposé \SGARefLink{sga3:ix:corollary:4:8:item:b}{IX~4.8(b)}), so
\(f_s\) is an isomorphism. The groups \(\ker f\) and
\(\operatorname{coker}f\) are of multiplicative type (Exposé \SGARefLink{sga3:ix:proposition:2:7}{IX~2.7})
and of constant type, since \(S\) is connected. They are therefore
trivial, and \(f\) is an isomorphism. If \(v=f^{-1}\), then
\[
  v|_{{}_nT}=v(n).
\]
Composing \(v\) with the projection \(G'\to G\) gives the required
homomorphism \(u:T\to G\). This proves existence; uniqueness follows
from Exposé \SGARefLink{sga3:ix:corollary:4:8:item:a}{IX~4.8(a)}.

\SGASubsection{4.2}{Application}
\SGARefTarget{sga3:xv:subsection:4:2}

We shall generalize Theorem \SGARefLink{sga3:ix:theorem:7:1}{IX~7.1}. Let \(A\) be a complete
noetherian local ring, let \(I\) be its maximal ideal, and put
\[
  S=\operatorname{Spec}A,\qquad
  S_m=\operatorname{Spec}(A/I^m).
\]
For every prescheme \(X\), put \(X_m=X\times_SS_m\).

Let \(G\) be an \(S\)-group prescheme of finite type, let \(T\) be an
\(S\)-torus, let \(q\) be an integer invertible on \(S\), and let
\[
  u_m:T_m\longrightarrow G_m\qquad(m\in\mathbb N)
\]
be a coherent family of homomorphisms. For \(n\) ranging over the
powers of \(q\), denote by \(u_m(n)\) the restriction of \(u_m\) to
\({}_nT_m\), and by
\[
  u(n):{}_nT\longrightarrow G
\]
the unique homomorphism extending the morphisms \(u_m(n)\) for all
\(m\) (\SGARefLink{sga3:xv:proposition:1-6:item:a}{Proposition 1.6(a)}). We call the family \((u_m)_{m\in\mathbb N}\)
\emph{admissible} if, for every point \(t\in S\), there exists a
closed affine subprescheme \(F_t\) of \(G_t\) containing
\(u(n)_t({}_nT_t)\) for every \(n\) that is a power of \(q\). As will
be seen, this property does not depend on \(q\).

\SGAStatement{Proposition 4.3}{With the preceding notation, the
canonical map
\[
 \operatorname{Hom}_{S\text{-}\mathrm{gr}}(T,G)
 \longrightarrow
 \varprojlim_m\operatorname{Hom}_{S_m\text{-}\mathrm{gr}}(T_m,G_m)
\]
identifies its source with the subset of the target consisting of
admissible coherent families.}
\SGARefTarget{sga3:xv:proposition:4:3}

Indeed, apply \SGARefLink{sga3:xv:theorem:4-1-bis}{Theorem 4.1 bis} with \(s\) the closed point of \(S\).

\SGAStatement{Corollary 4.4}{With the preceding notation, suppose in
addition that \(G\) has affine fibres. Then the canonical map
\[
 \operatorname{Hom}_{S\text{-}\mathrm{gr}}(T,G)
 \longrightarrow
 \varprojlim_m\operatorname{Hom}_{S_m\text{-}\mathrm{gr}}(T_m,G_m)
\]
is an isomorphism.}
\SGARefTarget{sga3:xv:corollary:4:4}

Indeed, when \(G\) has affine fibres every coherent family is
admissible.

\SGAStatement{Remark 4.5}{When \(G\) is separated, the ring \(A\) in
Proposition 4.3 may be replaced by a noetherian \(I\)-adic ring: one
uses EGA III~5.4.1 in place of Proposition 1.6(a).}
\SGARefTarget{sga3:xv:remark:4:5}

\SGASubsection{4.3}{Proof of \SGARefLink{sga3:xv:theorem:4-1}{Theorem 4.1}}
\SGARefTarget{sga3:xv:subsection:4:3}

\SGAStatement{Lemma 4.6}{Let \(k\) be a field, let \(G\) be a
\(k\)-algebraic group, let \(q\) be an integer prime to the
characteristic of \(k\), let \(r>0\), and let \(M(n)\), for \(n\)
ranging over the powers of \(q\), be a coherent family of subgroups
of \(G\), of multiplicative type and of type
\((\mathbb Z/n\mathbb Z)^r\). There exists a smallest closed
subscheme \(H\) of \(G\) containing all the \(M(n)\). Moreover, \(H\)
is a smooth, connected, commutative algebraic subgroup of \(G\), and
its formation commutes with extension of the base field.}
\SGARefTarget{sga3:xv:lemma:4:6}

Let \(M\subseteq\operatorname{ens}(G)\) be the union of the underlying
sets of all the \(M(n)\), and let \(H\) be the reduced closed
subscheme of \(G\) whose underlying space is the closure of \(M\).
Each \(M(n)\) is étale over \(k\), hence reduced, and is therefore
contained in \(H\). Thus \(H\) is the smallest closed subscheme
containing all the \(M(n)\).

Let \(\bar k\) be an algebraically closed extension of \(k\). By
construction the \(M(n)\) are schematically dense in \(H\) (Exposé
\SGARefLink{sga3:ix:definition:4:1}{IX~4.1}), hence the \(M(n)_{\bar k}\) are schematically dense in
\(H_{\bar k}\) (Exposé \SGARefLink{sga3:ix:lemma:4:5}{IX~4.5}). Since every \(M(n)_{\bar k}\) is
reduced, \(H_{\bar k}\) is the reduced closed subscheme of
\(G_{\bar k}\) whose underlying space is the closure of \(M_{\bar k}\).
Thus \(H\) is geometrically reduced (EGA IV~4.6.1).

Coherence of the family implies that \(M\) is stable under the group
law. The subset \(M\times_kM\) is dense in
\(\operatorname{ens}(H\times_kH)\), and \(H\times_kH\) is reduced
(EGA IV~4.6.5); it follows that \(H\) is an algebraic subgroup of
\(G\). It is smooth over \(k\), since it is geometrically reduced
(Exposé VI\textsubscript{A}~\SGARefLink{prop:sga3-VIA-1.3.1}{1.3.1}), and it is commutative because all
the \(M(n)\) are commutative.

It remains to prove connectedness. Let \(H^0\) be the identity
component of \(H\), let \(m\) be the number of geometric points of
\(H/H^0\), and let \(q^s\) be the \(q\)-primary part of \(m\). For
every \(n\) that is a power of \(q\), the subgroup
\({}_{q^s}M(n)\) is contained in \(H^0\). Coherence and the prescribed
type of \(M(n)\) give
\[
  {}_{q^s}M(nq^s)=M(n).
\]
Hence \(H^0\) contains every \(M(n)\), and therefore \(H^0=H\).

We can now sharpen condition (c) of \SGARefLink{sga3:xv:theorem:4-1}{Theorem 4.1}. By \SGARefLink{sga3:xv:lemma:4:6}{Lemma 4.6} there
is, for each \(t\in S\), a smallest closed subscheme \(H_t\) of
\(G_t\) containing all the \(M(n)_t\). Its formation commutes with
extension of the residue field. Thus \(H_t\otimes_t\bar t\) is
contained in the affine closed subset \(F_t\), and is affine; hence
\(H_t\) itself is affine. \SGARefLink{sga3:xv:lemma:4:6}{Lemma 4.6} also says that \(H_t\) is a
smooth, connected, commutative algebraic group. The structure theorem
for such affine groups over an algebraically closed field
(\emph{Bible}, 4, \SGARefLink{sga3:xv:section:4}{Theorem 4}) expresses \(H_t\) as a direct product of
a torus \(T_t\) and a unipotent group \(U_t\). Since every \(M(n)_t\)
is contained in \(T_t\), minimality gives \(H_t=T_t\). Consequently
\(H_t\) is a torus.

We begin the proof of \SGARefLink{sga3:xv:theorem:4-1}{Theorem 4.1}.

\medskip\noindent
\textbf{a) Uniqueness.}
This is Exposé \SGARefLink{sga3:ix:corollary:4:8:item:b}{IX~4.8(b)}.

\medskip\noindent
\textbf{b) The case in which \(S\) is the spectrum of a complete
discrete valuation ring \(A\) with algebraically closed residue field
\(k\).}
Let \(K\) be the fraction field of \(A\), let \(s\) and \(t\) be the
closed and generic points of \(S\), let \(J\) be the maximal ideal of
\(A\), and put
\[
  S_m=\operatorname{Spec}(A/J^m),\qquad X_m=X\times_SS_m.
\]

\smallskip\noindent
\textbf{First case: the point in \SGARefLink{sga3:xv:theorem:4-1:item:b}{Theorem 4.1(b)} is the generic point
\(t\).}
Let \(T'\) be the schematic closure of \(T_t\) in \(G\). Its closed
fibre \(T'_s\) is an algebraic subgroup of \(G_s\), of dimension
\(r\), containing every \(M(n)_s\), hence containing \(H_s\). The
torus \(H_s\) contains every \(M(n)_s\), so it has rank at least
\(r\). Consequently \(H_s\) has the same underlying space as the
identity component of \(T'_s\).

The identity component \((T')^0\) of \(T'\) is then a subgroup
prescheme of \(G\), flat and separated (Exposé
VI\textsubscript{B}~\SGARefLink{sga3:xv:theorem:5-2}{5.2}), whose generic fibre \(T_t\) is a torus and
whose reduced closed fibre \(H_s\) is a torus. Hence \((T')^0\) is a
torus (Exposé \SGARefLink{sga3:x:theorem:8-8}{X~8.8}); denote it by \(T\). The groups \({}_nT\) and
\(M(n)\) are smooth over \(S\), hence reduced, and have the same
underlying space. They therefore coincide, and \(T\) is the required
solution.

\smallskip\noindent
\textbf{Second case: the point in \SGARefLink{sga3:xv:theorem:4-1:item:b}{Theorem 4.1(b)} is the closed point
\(s\).}
After replacing \(G\) by the schematic closure in \(G\) of the
smallest algebraic subgroup \(H_t\) containing the family \(M(n)_t\)
(\SGARefLink{sga3:xv:lemma:4:6}{Lemma 4.6}), we may suppose that \(G_t\) is affine.

For every \(m\geq0\), \SGARefLink{sga3:xv:subsection:2:2}{Proposition 2.2} gives a unique sub-torus \(T_m\)
of \(G_m\) lifting \(T_s\) and satisfying
\[
  {}_nT_m=M(n)_m
  \qquad\text{for every }n\text{ that is a power of }q.
\]
Put \(T'=\mathbb G_{m,S}^{\,r}\). Since \(k\) is algebraically
closed, \(T_s\) is split, and there is a \(k\)-isomorphism
\[
  u_s:T'_s\xrightarrow{\sim}T_s.
\]
It lifts uniquely to an \(S_m\)-isomorphism
\(u_m:T'_m\to T_m\) (Exposé \SGARefLink{sga3:ix:corollary:3:3}{IX~3.3}). The coherent family
\((u_m)_{m\in\mathbb N}\) defines a morphism of formal completions
\(\widehat u:\widehat{T'}\to\widehat G\), fitting into
\[
\begin{tikzcd}[column sep=large,row sep=large]
\widehat{T'} \arrow[r,"\widehat u"] \arrow[d,"i"'] &
  \widehat G \arrow[d,"j"'] \\
T' & G ,
\end{tikzcd}
\]
We shall prove that \(\widehat u\) is algebraizable, reducing first
to the case in which \(G\) is affine.

\SGAStatement{Lemma 4.7}{Let \(S=\operatorname{Spec}A\), where \(A\)
is a discrete valuation ring, and let \(X,Y\) be \(S\)-preschemes
that are quasi-compact, quasi-separated, and flat over \(S\). Then the
canonical map
\[
 \Gamma(X,\mathcal O_X)\otimes_A\Gamma(Y,\mathcal O_Y)
 \longrightarrow
 \Gamma(X\times_SY,\mathcal O_{X\times_SY})
\]
is an isomorphism.}
\SGARefTarget{sga3:xv:lemma:4:7}

Let \(f:X\to S\) and \(g:Y\to S\) be the structural morphisms.
Because \(X\) and \(Y\) are quasi-compact and quasi-separated,
EGA I~9.2.2 and EGA IV~1.7.4 show that
\(f_*\mathcal O_X\) and \(g_*\mathcal O_Y\) are quasi-coherent
\(\mathcal O_S\)-algebras. Let \(\widetilde X\) and \(\widetilde Y\)
be the corresponding affine \(S\)-schemes. Since \(Y\) is flat over
\(S\), EGA III~1.4.15 and EGA IV~1.7.21 give an isomorphism
\[
 \Gamma(\widetilde X\times_SY,
        \mathcal O_{\widetilde X\times_SY})
 \xrightarrow{\sim}
 \Gamma(X\times_SY,\mathcal O_{X\times_SY}).
\]
The scheme \(\widetilde X\) is flat over \(S\), because over \(A\)
flatness is equivalent to torsion-freeness. Applying EGA III~1.4.15
again, with \(X\) and \(Y\) interchanged, yields
\[
\begin{aligned}
 \Gamma(\widetilde X,\mathcal O_{\widetilde X})
 \otimes_A\Gamma(\widetilde Y,\mathcal O_{\widetilde Y})
 &\xrightarrow{\sim}
 \Gamma(\widetilde X\times_S\widetilde Y,
        \mathcal O_{\widetilde X\times_S\widetilde Y})\\
 &\xrightarrow{\sim}
 \Gamma(\widetilde X\times_SY,
        \mathcal O_{\widetilde X\times_SY}),
\end{aligned}
\]
which proves the lemma.

The \(S\)-group \(G\) is flat and of finite type over \(S\), hence
quasi-compact and quasi-separated. Applying the lemma to
\(G\times_SG\) gives
\[
 \Gamma(G,\mathcal O_G)\otimes_A\Gamma(G,\mathcal O_G)
 \xrightarrow{\sim}
 \Gamma(G\times_SG,\mathcal O_{G\times_SG}).
\]
Thus the multiplication \(m_G:G\times_SG\to G\) determines a
homomorphism
\[
 \Gamma(G,\mathcal O_G)\longrightarrow
 \Gamma(G,\mathcal O_G)\otimes_A\Gamma(G,\mathcal O_G),
\]
and hence an \(S\)-morphism
\[
  m_{\widetilde G}:\widetilde G\times_S\widetilde G
  \longrightarrow\widetilde G,
\]
where \(\widetilde G\) is the affine \(S\)-scheme corresponding to
\(\Gamma(G,\mathcal O_G)\). The usual identities show that
\(m_{\widetilde G}\) makes \(\widetilde G\) an \(S\)-group scheme and
that the canonical morphism \(v:G\to\widetilde G\) is a homomorphism.

\SGAStatement{Remark 4.8}{The group \(\widetilde G\) plays the role of
a largest affine quotient of \(G\). In fact it is an fpqc quotient of
\(G\), and is therefore of finite type over \(S\) (XVII, Appendix
III, 2).}
\SGARefTarget{sga3:xv:remark:4:8}

In our situation the generic fibre \(G_t\) is affine, so
\(v_t:G_t\to\widetilde G_t\) is an isomorphism (EGA I~9.3.3). Since
\(\widetilde G\) is affine, the coherent family
\[
  w_m=v_m\circ u_m:T'_m\longrightarrow\widetilde G_m
  \qquad(m\in\mathbb N)
\]
comes from a unique \(S\)-homomorphism (Exposé \SGARefLink{sga3:ix:theorem:7:1}{IX~7.1})
\[
  w:T'\longrightarrow\widetilde G.
\]
Let \(T_t\) be the sub-torus
\[
  T_t=v_t^{-1}\bigl(w_t(T'_t)\bigr)
\]
of \(G_t\). Its rank is at most \(r\). We show that it contains every
\(M(n)_t\). Let
\[
 u(n)_m:({}_nT')_m\xrightarrow{\sim}M(n)_m
\]
be the restriction of \(u_m\). This coherent family comes from a
unique \(S\)-isomorphism
\[
  u(n):{}_nT'\xrightarrow{\sim}M(n),
\]
because \(M(n)\) is finite over \(S\). For every \(m\geq0\),
\[
 w_m|_{({}_nT')_m}
  =(v_m\circ u_m)|_{({}_nT')_m}
  =(v\circ u(n))_m.
\]
It follows from \SGARefLink{sga3:xv:proposition:1-6:item:a}{Proposition 1.6(a)} that
\[
  w|_{{}_nT'}=v\circ u(n).
\]
On the generic fibre,
\[
 v_t^{-1}\!\left(w_t(({}_nT')_t)\right)
 =u(n)_t(({}_nT')_t)=M(n)_t.
\]
Thus \(T_t\) contains every \(M(n)_t\), and consequently has rank
\(r\).

We conclude as in the first case. Let \(T''\) be the schematic closure
of \(T_t\) in \(G\). Since \(T_t\) contains every \(M(n)_t\), the
closed fibre \(T''_s\) contains every \(M(n)_s\), and hence contains
\(T_s\) by the density theorem. On the other hand \(T''\) is flat
over \(S\) and \(T_t\) has dimension \(r\), so \(T''_s\) has
dimension \(r\) (Exposé VI\textsubscript{B}~4). Therefore \(T_s\)
has the same underlying space as the identity component of \(T''_s\),
and the identity component of \(T''\) is the required sub-torus of
\(G\).

\medskip\noindent
\textbf{c) End of the proof.}
To prove existence of \(T\), we may assume that \(S\) is reduced
(\SGARefLink{sga3:xv:subsection:2:2}{Proposition 2.2}). By \SGARefLink{sga3:xv:theorem:3-1}{Theorem 3.1}, \SGARefLink{sga3:xv:theorem:3-1:item:ii}{(ii)}\(\Rightarrow\)\SGARefLink{sga3:xv:theorem:3-1:item:i}{(i)}, it is then
enough to show that the set \(U\) of points \(x\in S\) for which there
exists a sub-torus \(T_x\) of \(G_x\) satisfying
\[
  {}_nT_x=M(n)_x
  \qquad\text{for every }n\text{ that is a power of }q
\]
equals \(\operatorname{ens}(S)\). Such a torus is necessarily unique,
and fpqc descent reduces its existence to existence after an extension
of the residue field of \(x\) (Exposé \SGARefLink{sga3:ix:corollary:4:8:item:b}{IX~4.8(b)}).

Since \(S\) is locally noetherian and connected and \(U\) is nonempty
(it contains the point \(s\) of \SGARefLink{sga3:xv:theorem:4-1:item:b}{Theorem 4.1(b)}), it is enough to prove
the following: if \(x,x'\in S\), if \(x\) is a specialization of
\(x'\), and if one of these two points lies in \(U\), then both lie
in \(U\). The usual method (EGA II~7.1.9) reduces this to the case
where \(S\) is the spectrum of a complete discrete valuation ring
with algebraically closed residue field, which was treated in part
(b). This completes the proof of \SGARefLink{sga3:xv:theorem:4-1}{Theorem 4.1}.

% END inputs/XV/tex/components/04_expose_XV_section4.tex

% BEGIN inputs/XV/tex/components/05_expose_XV_section5.tex
% Authority: Expo15.pdf, local pp. 32--38 / combined pp. 902--908 /
% printed pp. 409--421.
% Stops immediately before §6.

\SGASection{5}{Representability of the functor of smooth subgroups
equal to their connected normalizer}
\SGARefTarget{sga3:xv:section:5}

\SGAStatement{\SGAInlineRefTarget{sga3:xv:proposition:5-1}Proposition 5.1}{Let \(S\) be a prescheme, let \(G\)
be an \(S\)-group prescheme of finite presentation, and let \(H\) be
a subgroup prescheme of \(G\), smooth with connected fibres. Then:
\begin{enumerate}
\renewcommand{\labelenumi}{\textup{\alph{enumi})}}
\item the normalizer \(N\) of \(H\) in \(G\) is represented by a
closed subprescheme of \(G\), of finite presentation over \(S\);
\item\SGAInlineRefTarget{sga3:xv:proposition:5-1:item:b} the following conditions are equivalent:
\begin{enumerate}
\renewcommand{\labelenumii}{\textup{\roman{enumii})}}
\item the canonical immersion \(H\to N\) is an open immersion;
\item \(N\) is smooth along the identity section, and its identity
component, which is then representable (Exposé
VI\textsubscript{B}~\SGARefLink{sga3:xv:lemma:3:10}{3.10}), equals \(H\);
\item for every \(s\in S\), one has \(H_s=(N_s)^0\).
\end{enumerate}
\end{enumerate}}
\SGARefTarget{sga3:xv:proposition:5:1}

\SGAProof{Proof}
The group prescheme \(H\) is locally of finite presentation over
\(S\), since it is smooth, and has connected fibres; it is therefore
of finite presentation over \(S\) (Exposé
VI\textsubscript{B}~\SGARefLink{par:sga3-VIA-5.3.3}{5.3.3}). Part (a) follows from Exposé \SGARefLink{sga3:xi:corollary:6:11}{XI~6.11}. The
equivalence in \SGARefLink{sga3:xv:proposition:5-1:item:b}{(b)}, recalled here for reference, was proved in Exposé
VI\textsubscript{B}~\SGARefLink{prop:sga3-VIB-6.5.1}{6.5.1}.

We can now state the main theorem of this section.

\SGAStatement{\SGAInlineRefTarget{sga3:xv:theorem:5-2}Theorem 5.2}{Let \(S\) be a prescheme and \(G\) an
\(S\)-group prescheme of finite presentation over \(S\). Let
\(\mathcal L_G\), or simply \(\mathcal L\), be the \(S\)-functor
defined, for every \(S\)-prescheme \(S'\), by
\[
\mathcal L_G(S')=
\left\{
\begin{array}{c}
\text{subgroup preschemes of \(G_{S'}\), smooth over \(S'\),}\\
\text{with connected fibres, equal to their connected normalizer}
\end{array}
\right\}.
\]
Then \(\mathcal L\) is represented by an \(S\)-prescheme that is the
union of an increasing family \((U_i)_{i\in\mathbb N}\) of open
subpreschemes, each quasi-projective and of finite presentation over
\(S\), and hence in particular separated over \(S\).}
\SGARefTarget{sga3:xv:theorem:5:2}

\medskip\noindent\textbf{Initial reductions.}
For every \(r\geq0\), let \(\mathcal L_r\) be the subfunctor of
\(\mathcal L\) defined by
\[
\mathcal L_r(S')=
\left\{
\begin{array}{c}
H\subseteq G_{S'}:\ H\text{ is smooth with connected fibres,}\\
H\text{ equals its connected normalizer, and }\dim(H/S')=r
\end{array}
\right\}.
\]
The fibre dimension of a smooth group is locally constant. Hence
\(\mathcal L_r\to\mathcal L\) is represented by an immersion that is
both open and closed. It is therefore enough to prove the asserted
representability and finiteness properties for every
\(\mathcal L_r\), for then
\[
  \mathcal L=\coprod_{r\in\mathbb N}\mathcal L_r.
\]

For \(n\geq0\), an \(S\)-prescheme \(S'\), and a subgroup prescheme
\(H\) of \(G_{S'}\), denote by \(H^{(n)}\) the \(n\)-th normal
invariant of the identity section \(S'\to H\) (EGA IV~16.1.2).
Thus \(H^{(n)}\) is a finite-type \(\mathcal O_{S'}\)-module
corresponding to the \(n\)-th infinitesimal neighbourhood of the
identity section of \(H\). If \(H\) is smooth of relative dimension
\(r\), then \(H^{(n)}\) is locally free of a rank
\(\varphi(n,r)\) depending only on \(n,r\). Since \(H\) is a
subprescheme of \(G_{S'}\), there is a canonical epimorphism,
compatible with base change,
\[
 G^{(n)}\otimes_{\mathcal O_S}\mathcal O_{S'}
 \xrightarrow{\sim}G_{S'}^{(n)}
 \longrightarrow H^{(n)}.
\]

Introduce the projective \(S\)-scheme
\[
  P_{\varphi(n,r)}
   =\operatorname{Grass}_{\varphi(n,r)}(G^{(n)})
\]
(EGA I, 2nd ed., 9.7; see also Cartan Seminar 1960/61,
Exposé 14 by A. Grothendieck). The assignment
\[
  H\longmapsto H^{(n)}
\]
therefore defines a canonical morphism
\[
  u_{n,r}:\mathcal L_r\longrightarrow P_{\varphi(n,r)}.
\]
The group \(G\) acts naturally on \(G^{(n)}\), and hence on
\(P_{\varphi(n,r)}\), through
\[
 \operatorname{int}:G\longrightarrow
 \operatorname{Aut}_{S\text{-}\mathrm{gr}}(G),
 \qquad g\longmapsto\operatorname{int}(g).
\]

If \(S'\) is quasi-compact over \(S\) and \(H\in\mathcal L_r(S')\),
then for sufficiently large \(n\) one has (Exposé \SGARefLink{sga3:xi:corollary:6:11}{XI~6.11})
\[
 \operatorname{Norm}_{G_{S'}}(H)
 =\operatorname{Norm}_{G_{S'}}(H^{(n)}).
\]
For each \(n\geq0\), define the subfunctor \(\mathcal L_r^n\) of
\(\mathcal L_r\) by
\[
\mathcal L_r^n(S')=
\left\{
 H\in\mathcal L_r(S')\ \middle|\
 \operatorname{Norm}_{G_{S'}}(H)
 =\operatorname{Norm}_{G_{S'}}(H^{(n)})
\right\}.
\]

\medskip\noindent
\textbf{Representability of \(\mathcal L_r^n\).}
The integer \(r\) remains fixed through the proof, so we write
\(\mathcal L,\mathcal L^n,P_n,u_n\) in place of
\(\mathcal L_r,\mathcal L_r^n,P_{\varphi(n,r)},u_{n,r}\). Let
\[
 v_n:\mathcal L^n\longrightarrow P_n
\]
be the restriction of \(u_n\). The definition of \(\mathcal L^n\)
and \SGARefLink{sga3:xv:proposition:5-1:item:b}{Proposition 5.1(b)}(ii) show that \(v_n\) is a monomorphism.
More precisely:

\SGAStatement{Lemma 5.3}{The morphism \(v_n\) is an immersion of
finite presentation. Consequently \(\mathcal L^n\) is represented by
an \(S\)-prescheme, quasi-projective and of finite presentation over
\(S\).}
\SGARefTarget{sga3:xv:lemma:5:3}

After replacing \(S\) by \(P_n\), the usual method reduces the
assertion to the following. Let \(Q\in P_n(S)\), and define a
subfunctor \(F\) of the final functor \(h_S\) by
\[
F(S')=
\begin{cases}
h_S(S'),&\text{if there exists }H\in\mathcal L^n(S')
                    \text{ with }H^{(n)}=Q_{S'},\\
\varnothing,&\text{otherwise.}
\end{cases}
\]
We must prove that \(F\to S\) is an immersion of finite presentation.

Let \(N\) be the normalizer of \(Q\) in \(G\). It is a subgroup
prescheme of finite presentation over \(S\), namely the inverse image
of \(Q\) under
\[
  G\longrightarrow P_n,\qquad g\longmapsto g\cdot Q.
\]
We claim that \(F\) is equivalently described by
\[
F(S')=
\begin{cases}
h_S(S'),&
\begin{array}{l}
N_{S'}\text{ is smooth along the identity section, of}\\
\text{relative dimension }r,\text{ and }(N_{S'})^{(n)}=Q_{S'},
\end{array}\\[2mm]
\varnothing,&\text{otherwise.}
\end{cases}
\]

Call the first and second descriptions \(F_1\) and \(F_2\).
Suppose \(F_1(S')=h_S(S')\), and let \(H\in\mathcal L^n(S')\) with
\(H^{(n)}=Q_{S'}\). Then
\[
\begin{aligned}
\operatorname{Norm}_{G_{S'}}(H)
 &=\operatorname{Norm}_{G_{S'}}(H^{(n)})\\
 &=\operatorname{Norm}_{G_{S'}}(Q_{S'})
 =N_{S'}.
\end{aligned}
\]
By \SGARefLink{sga3:xv:proposition:5-1:item:b}{Proposition 5.1(b)}(ii), \(N_{S'}\) is smooth along the identity
section and its identity component is \(H\). Thus it has relative
dimension \(r\); since \(H\) is open in \(N_{S'}\),
\[
  (N_{S'})^{(n)}=H^{(n)}=Q_{S'}.
\]
Hence \(F_2(S')=h_S(S')\).

Conversely, suppose \(F_2(S')=h_S(S')\). The identity component of
\(N_{S'}\) is then represented (Exposé VI\textsubscript{B}~\SGARefLink{sga3:xv:lemma:3:10}{3.10}) by
a subgroup prescheme \(H\), smooth over \(S'\), with connected fibres
of dimension \(r\). Since \(H\) is invariant and open in \(N_{S'}\),
\[
\begin{aligned}
N_{S'}&\subseteq\operatorname{Norm}_{G_{S'}}(H)
\subseteq\operatorname{Norm}_{G_{S'}}(H^{(n)})\\
&=\operatorname{Norm}_{G_{S'}}((N_{S'})^{(n)})
=\operatorname{Norm}_{G_{S'}}(Q_{S'})=N_{S'}.
\end{aligned}
\]
All inclusions are equalities. The first says that \(H\) equals its
connected normalizer, and the second says
\(H\in\mathcal L^n(S')\). Hence \(F_1=F_2\).

We retain the second description. Condition (a) is represented by an
immersion of finite presentation, by the following lemma.

\SGAStatement{\SGAInlineRefTarget{sga3:xv:lemma:5-4}Lemma 5.4}{Let \(S\) be a prescheme, let \(X\) be
locally of finite presentation over \(S\), let \(\sigma:S\to X\) be
a section, let \(r\geq0\), and define the subfunctor \(L\) of \(S\)
by
\[
L(S')=
\begin{cases}
h_S(S'),&
\begin{array}{l}
X_{S'}\text{ is smooth along }\sigma_{S'},\text{ of relative}\\
\text{dimension }r\text{ at the points of }\sigma_{S'}(S'),
\end{array}\\[1mm]
\varnothing,&\text{otherwise.}
\end{cases}
\]
Then:
\begin{enumerate}
\renewcommand{\labelenumi}{\textup{\alph{enumi})}}
\item \(L\to S\) is an immersion of finite presentation;
\item let \(\mathcal J\) be the conormal sheaf of \(S\to X\).
If for every \(s\in S\) the vector space
\(\mathcal J\otimes_{\mathcal O_{S,s}}\kappa(s)\) has rank at most
\(r\), then \(L\to S\) is a closed immersion.
\end{enumerate}}
\SGARefTarget{sga3:xv:lemma:5:4}

\SGAProof{Proof}
The functor \(L\) is local on \(S\), so we may suppose \(S\) affine.
After replacing \(X\) by a neighbourhood of \(\sigma(S)\), we may
suppose \(X\) of finite presentation, and then reduce to noetherian
\(S\) by EGA IV~8.9. In part (b), use also compatibility of the
conormal sheaf with base change (EGA IV~16.6.4) and constructibility
of the rank of the fibres of \(\mathcal J\).

For an \(S\)-prescheme \(S'\), let
\(\mathcal J'=\mathcal J_{S'}\), let
\(S^\bullet(\mathcal J')\), canonically isomorphic to
\(S^\bullet(\mathcal J)_{S'}\), be the symmetric algebra of
\(\mathcal J'\) over \(\mathcal O_{S'}\), let
\(\operatorname{Gr}^\bullet(\sigma_{S'})\) be the graded algebra
associated with \(\sigma_{S'}\) (EGA IV~16), and let
\[
 \sigma_{n,S'}:S^n(\mathcal J')
 \longrightarrow\operatorname{Gr}^n(\sigma_{S'})
\]
be the canonical epimorphism. By EGA IV~17.12.3 and
EGA \(0_{\mathrm{IV}}\)~19.5.4, \(X_{S'}\) is smooth along
\(\sigma_{S'}\) of relative dimension \(r\) if and only if:
\begin{enumerate}
\renewcommand{\labelenumi}{\textup{\roman{enumi})}}
\item \(\mathcal J'\) is locally free of rank \(r\);
\item \(\sigma_{n,S'}\) is an isomorphism for every \(n\geq1\).
\end{enumerate}

TDTE IV, \SGARefLink{sga3:xv:lemma:3:6}{Lemma 3.6}, shows that the functor making \(\mathcal J\)
locally free of rank \(r\) is represented by a subprescheme \(S_1\),
closed in \(S\) in case (b). Replacing \(S\) by \(S_1\), assume that
\(\mathcal J\) is locally free of rank \(r\). Proceed by induction on
\(n\). Suppose a closed subprescheme \(S_{n-1}\) represents the
functor making \(\sigma_{q,S'}\) injective for \(q\leq n-1\).
The subfunctor\footnote{The editors replaced ``functor'' by
``subfunctor'' here, and ask whether the intended wording is instead
``which makes \(\sigma_{q,S'}\) injective for \(q\leq n\)''.}
making \(\sigma_{n,S'}\) injective is represented by a closed
subprescheme \(S_n\) of \(S_{n-1}\).

Indeed, after replacing \(S\) by \(S_{n-1}\), all
\(\sigma_{q,S}\) for \(q\leq n-1\) are isomorphisms. The
\((n-1)\)-st normal invariant \(X^{(n-1)}\) of \(\sigma_S\) has a
filtration with successive quotients
 \(\operatorname{Gr}^0(\sigma_S),\ldots,
 \operatorname{Gr}^{n-1}(\sigma_S)\), all locally free; hence
\(X^{(n-1)}\) is flat over \(S\). Compatibility of the normal
invariants with base change gives
\[
\begin{aligned}
\operatorname{Gr}^n(\sigma_{S'})
 &=\ker\!\left(X_{S'}^{(n)}\to X_{S'}^{(n-1)}\right)
  =\operatorname{Gr}^n(\sigma)_{S'},\\
\sigma_{n,S'}&=(\sigma_{n,S})_{S'}.
\end{aligned}
\]
Thus we must represent the functor making \(\sigma_{n,S}\) injective.
Locally we may suppose \(S\) affine and \(S^n(\mathcal J)\) free.
The required functor is represented by the closed subscheme defined
by the ideal generated by the coordinates of
\(\ker\sigma_{n,S}\) in a basis of \(S^n(\mathcal J)\).

Since \(S\) is noetherian, the descending chain \((S_n)\) stabilizes.
Its stable value represents \(L\), proving the lemma.

Return to \(\mathcal L^n\). After replacing \(S\) by the subprescheme
provided by \SGARefLink{sga3:xv:lemma:5-4}{Lemma 5.4}, assume \(N\) smooth along the identity section
of relative dimension \(r\). Then \(\mathcal L^n\) is the coincidence
functor of two sections \(h,g:S\to P_n\), corresponding to \(Q\) and
\(N^{(n)}\). It is therefore represented by the inverse image of the
diagonal of \(P_n\times_SP_n\) under the finite-presentation morphism
\(h\times_Sg\). This proves \SGARefLink{sga3:xv:lemma:5:3}{Lemma 5.3}.

\medskip\noindent
\textbf{The transition morphisms
\(\mathcal L^n\to\mathcal L^m\).}
If \(H\in\mathcal L^n(S')\), then \(H\in\mathcal L^m(S')\) for every
\(m\geq n\). Thus there are natural monomorphisms
\[
  u_n^m:\mathcal L^n\longrightarrow\mathcal L^m
  \qquad(m\geq n).
\]

\SGAStatement{Lemma 5.5}{The morphism \(u_n^m\) is an open immersion.}
\SGARefTarget{sga3:xv:lemma:5:5}

After a base change, it is enough to consider \(H\in\mathcal L^m(S)\)
and
\[
 N=\operatorname{Norm}_G(H)=\operatorname{Norm}_G(H^{(m)}),
 \qquad
 N'=\operatorname{Norm}_G(H^{(n)}).
\]
Let \(D\) be the coincidence functor of \(N,N'\): \(D(S')=h_S(S')\)
if \(N_{S'}=N'_{S'}\), and \(D(S')=\varnothing\) otherwise. We must
show that \(D\to S\) is an open immersion.

The same functor is obtained by requiring that the immersion
\(H\to N'\) become open. One implication follows from
\SGARefLink{sga3:xv:proposition:5-1}{Proposition 5.1}. Conversely, if \(H_{S'}\to N'_{S'}\) is open, then,
because \(H\) has connected fibres, \(H_{S'}\) is the identity
component of \(N'_{S'}\) (Exposé VI\textsubscript{B}~\SGARefLink{sga3:xv:lemma:3:10}{3.10}), hence is
invariant in \(N'_{S'}\). Thus \(N'_{S'}\subseteq N_{S'}\). The
opposite inclusion always holds, so equality follows. Since \(H,N'\)
are of finite presentation and \(H\) is flat over \(S\), Exposé
VI\textsubscript{B}~\SGARefLink{sga3:xv:remark:2:6}{2.6} shows that \(D\to S\) is open.

\medskip\noindent\textbf{End of the proof of \SGARefLink{sga3:xv:theorem:5-2}{Theorem 5.2}.}
The functors \(\mathcal L^n\) are represented, and the transition
morphisms \(u_n^m\) are compatible open immersions. There is therefore
an \(S\)-prescheme \(X\), the union of an increasing sequence of open
subpreschemes \((X_i)_{i\in\mathbb N}\), such that \(X_i\) represents
\(\mathcal L^i\) and \(X_i\to X_j\) corresponds to \(u_i^j\). In the
category of sheaves for the Zariski topology on \(\mathrm{Sch}/S\),
\[
  X=\varinjlim_iX_i,\qquad
  \mathcal L=\varinjlim_i\mathcal L^i.
\]
Thus \(X\) represents \(\mathcal L\), completing the proof.

\SGAStatement{Remark 5.6}{Suppose all residue characteristics of
\(S\) are zero. Then, for every \(r\geq0\),
\(\mathcal L_r=\mathcal L_r^1\), so \(\mathcal L_r\) is represented
by an \(S\)-prescheme quasi-projective and of finite presentation over
\(S\).}
\SGARefTarget{sga3:xv:remark:5:6}

Indeed, for \(H\in\mathcal L_r(S')\), it suffices to prove that
\[
  H\longrightarrow N=\operatorname{Norm}_{G_{S'}}(H^{(1)})
\]
is an open immersion. Because \(H\) is flat and both \(H,N\) are of
finite presentation, this may be checked on the fibres (Exposé
VI\textsubscript{B}~\SGARefLink{sga3:xv:remark:2:6}{2.6}). Cartier's theorem (Exposé
VI\textsubscript{B}~\SGARefLink{sga3:xv:proposition:1-6}{1.6}) gives, for a connected algebraic subgroup
\(H\) of an algebraic group in characteristic zero,
\[
 \operatorname{Norm}_G(H)
 =\operatorname{Norm}_G(\operatorname{Lie}H)
 =\operatorname{Norm}_G(H^{(1)}).
\]
\footnote{For example, one may apply Proposition II.6.1 of
M. Demazure and P. Gabriel, \emph{Groupes algébriques} I,
Masson--North Holland, 1970.}

In nonzero residue characteristic the subfunctors
\(\mathcal L_r^n\) can form a strictly increasing sequence, even when
\(S\) is quasi-compact, and then \(\mathcal L_r\) is not represented
by an \(S\)-prescheme quasi-compact over \(S\). For example, over a
field \(k\) of characteristic \(p>0\), let
\[
 T=\mathbb G_m\times\mathbb G_m,\qquad
 U=\mathbb G_a\times\mathbb G_a,
\]
and let \(G=T\ltimes U\), where
\((t,t')\cdot(u,u')=(tu,t'u')\). For \(n>0\), let \(U_n\subset U\)
be defined by \(u'=u^{p^n}\), and let \(T_n\subset T\) be defined by
\(t'=t^{p^n}\). Then \(T_n\) acts on \(U_n\), and
\(G_n=T_nU_n\) is smooth, connected, and equal to its normalizer in
\(G\). For every \(m\), the distinct groups \(G_n\), \(n\geq m\),
have the same infinitesimal neighbourhood of order \(p^m\).

\SGAStatement{Remark 5.7}{There is a canonical invertible sheaf
\(L\) on \(\mathcal L\) whose restriction to every open subset of
\(\mathcal L\) quasi-compact over \(S\) is \(S\)-ample. One may take
\[
  L=(\det(\operatorname{Lie}H))^{-1},
\]
where \(H\subseteq G_{\mathcal L}\) is the universal subgroup.}
\SGARefTarget{sga3:xv:remark:5:7}

Recall that for a finite-rank locally free sheaf \(\mathcal F\), its
determinant restricts on the open-and-closed locus where
\(\operatorname{rank}\mathcal F=r\) to \(\bigwedge^r\mathcal F\).
It is enough to restrict to \(\mathcal L_r^n\). Let
\(v_n:\mathcal L_r^n\to P_n\) be the canonical immersion, and let
\(Q\) be the universal locally free quotient on the Grassmannian
\(P_n\). By construction,
\[
  v_n^*Q=H^{(n)}.
\]
Since \(\det Q\) is the canonical ample sheaf on \(P_n\) (EGA I,
2nd ed., 9.7), \(\det H^{(n)}\) is ample relative to \(S\)
(EGA II~4.6.13(i) bis).

Let \(\mathcal J=H^{(1)}\) be the conormal sheaf, and let
\(S^q(\mathcal J)\) be the homogeneous degree-\(q\) part of its
symmetric algebra. Smoothness of \(H\) gives
\[
  \det H^{(n)}
  \xrightarrow{\sim}
  \bigotimes_{1\leq q\leq n}\det S^q(\mathcal J).
\]
For every locally free \(\mathcal J\) of rank \(r\) and every \(q>0\)
there is a canonical isomorphism
\[
  \det S^q(\mathcal J)
  \xrightarrow{\sim}
  (\det\mathcal J)^{\otimes s},
\]
where \(s>0\) depends only on \(r,q\). Thus, for some \(s>0\),
\[
  \det H^{(n)}\simeq(\det\mathcal J)^{\otimes s}.
\]
Since \(\det\mathcal J=(\det\operatorname{Lie}H)^{-1}\), EGA
II~4.5.6 proves that \(L\) is \(S\)-ample.

\SGAStatement{Remark 5.8}{Let \(G,H\) be \(S\)-group preschemes of
finite presentation and let \(i:H\to G\) be a monomorphism. Suppose
\(H\) is smooth over \(S\) with connected fibres. The normalizer
\(N=\operatorname{Norm}_G(H)\) is then represented by a closed
subgroup prescheme of \(G\), of finite presentation over \(S\)
(Exposé \SGARefLink{sga3:xi:corollary:6:11}{XI~6.11}). If, in addition, \(N\) is smooth along the identity
section and has the same relative dimension as \(H\), then
\(H=N^0\); in particular \(H\) is a subgroup prescheme of \(G\).}
\SGARefTarget{sga3:xv:remark:5:8}

Indeed, \(N^0\) is represented by an open subgroup prescheme of \(N\),
smooth over \(S\) (Exposé VI\textsubscript{B}~\SGARefLink{sga3:xv:lemma:3:10}{3.10}), and \(i\)
factors through \(N^0\). For every \(s\in S\), the connected smooth
algebraic groups \(H_s,(N^0)_s\) have the same dimension and hence
coincide. Since \(H\) is flat, \(H\to N^0\) is an isomorphism
(EGA IV~17.9.5). Thus the functor \(\mathcal L\) of this section may
equivalently be described as the functor of subgroup preschemes
\(H\) of \(G\) that are smooth over \(S\), have connected fibres,
and equal their connected normalizer.

% END inputs/XV/tex/components/05_expose_XV_section5.tex

% BEGIN inputs/XV/tex/components/06_expose_XV_section6.tex
% Authority: Expo15.pdf, local pp. 38--49 / combined pp. 908--919 /
% printed pp. 421--444.
% Stops immediately before §7.

\SGASection{6}{Functor of Cartan subgroups and functor of parabolic
subgroups}
\SGARefTarget{sga3:xv:section:6}

When \(G\) is a smooth connected algebraic group defined over an
algebraically closed field \(k\), sub-tori of \(G\), maximal
sub-tori, Cartan subgroups (Exposé XII~1), Borel subgroups (Exposé
\SGARefLink{sga3:xiv:4.1}{XIV~4.1}), and parabolic subgroups (Exposé \SGARefLink{sga3:xiv:4.8.1}{XIV~4.8 bis}) have already
been defined. We extend these notions to a group prescheme over an
arbitrary base as follows.

\textbf{\SGAInlineRefTarget{sga3:xv:definition:6-1}Definition 6.1.} \emph{Let \(S\) be a prescheme, let \(G\)
be an \(S\)-group prescheme of finite presentation, let \(H\) be an
\(S\)-group prescheme, and let \(i:H\to G\) be an
\(S\)-monomorphism making \(H\) a subgroup of \(G\). We say that
\(H\) is a sub-torus of \(G\) (respectively, a maximal sub-torus, a
Cartan subgroup, a Borel subgroup, or a parabolic subgroup) if:}

\begin{enumerate}
\renewcommand{\labelenumi}{\textup{\roman{enumi})}}
\item
  \emph{i) \(H\) is smooth over \(S\).}
\item
  \emph{ii) For every geometric point \(\bar s\) above \(S\),
  \(H_{\bar s}\) is a sub-torus of
  \((G_{\bar s})^{0}_{\mathrm{red}}\) (respectively, a maximal
  sub-torus, a Cartan subgroup, a Borel subgroup, or a parabolic
  subgroup).}
\end{enumerate}
\SGARefTarget{sga3:xv:definition:6:1}

\textbf{\SGAInlineRefTarget{sga3:xv:remark:6-1-bis}Remarks 6.1 bis.}

\begin{enumerate}
\renewcommand{\labelenumi}{\textup{\alph{enumi})}}
\item
  If the \(S\)-group \(H\) is a sub-torus of \(G\)
  (respectively, \(\ldots\)), its fibres are connected, and
  consequently \(H\) is of finite presentation over \(S\) (Exposé
  VI\textsubscript{B}~\SGARefLink{par:sga3-VIA-5.3.3}{5.3.3}).
\item
  If \(H\) is a sub-torus of \(G\), then \(H\) is a torus in the
  sense of Exposé IX, as follows immediately from Exposé \SGARefLink{sga3:x:theorem:8-1}{X~8.1}.
  Moreover, the monomorphism \(i:H\to G\) is an immersion (see \SGARefLink{sga3:xv:lemma:8:3}{8.3}
  below).
\item
  If \(G\) is smooth over \(S\), with connected fibres, and if
  \(H\) is a Cartan subgroup of \(G\) (respectively, a Borel
  subgroup or a parabolic subgroup), the monomorphism \(i:H\to G\)
  is an immersion. Thus our definitions coincide with those
  introduced in Exposés XII and XIV. Indeed, \(H\) is then equal to
  its connected normalizer (by \SGARefLink{sga3:xii:theorem:6-6:item:c}{XII~6.6(c)}, \SGARefLink{sga3:xiv:4.8}{XIV~4.8} and \SGARefLink{sga3:xiv:paragraph:4-8-bis}{4.8 bis}),
  and it suffices to apply \SGARefLink{sga3:xv:remark:5:8}{5.8}.
\end{enumerate}

\textbf{Definition 6.1 ter.} \emph{Let \(S\) be a prescheme, let
\(G\) be an \(S\)-group prescheme locally of finite type, and let
\(s\) be a point of \(S\). The \textbf{nilpotent rank} of \(G\) at
\(s\), denoted by \(\rho_n(s)\), is the dimension of the Cartan
subgroups of \((G_s)^0_{\mathrm{red}}\). The reductive rank
\(\rho_r(s)\), unipotent rank \(\rho_u(s)\), and abelian rank
\(\rho_{\mathrm{ab}}(s)\) are defined similarly (see Exposé
\SGARefLink{sga3:x:remark:8-7}{X~8.7}).}
\SGARefTarget{sga3:xv:definition:6:1:ter}

If \(G\) is now a smooth connected algebraic group defined over an
algebraically closed field \(k\), recall that the \textbf{radical}
of \(G\), denoted by \(\operatorname{rad}(G)\), is the largest
invariant smooth connected solvable algebraic subgroup of \(G\);
\(G/\operatorname{rad}(G)\) is then semisimple (use Exposé \SGARefLink{sga3:xii:6.1}{XII~6.1}
to reduce to the affine case). If \(G\) is moreover affine, the
\textbf{unipotent radical} \(\operatorname{rad}^{u}(G)\) is defined
as the largest invariant smooth connected unipotent algebraic
subgroup of \(G\); \(G/\operatorname{rad}^{u}(G)\) is then
reductive.

\textbf{\SGAInlineRefTarget{sga3:xv:proposition:6-2}Proposition 6.2.} \emph{Let \(S\) be a prescheme, let \(G\)
and \(H\) be two \(S\)-group preschemes of finite presentation, and
let \(i:H\to G\) be an \(S\)-monomorphism making \(H\) a subgroup
of \(G\). Let \(P(s)\) be one of the following properties concerning
a point \(s\) of \(S\):}

\begin{enumerate}
\renewcommand{\labelenumi}{\textup{\roman{enumi})}}
\item \emph{\((G_{\bar s})^{0}_{\mathrm{red}}\) is an abelian
variety (respectively, is affine, is a torus, or is unipotent).}
\item \emph{\((H_{\bar s})^{0}_{\mathrm{red}}\) is a maximal torus
of \(G_{\bar s}\).}
\item\SGAInlineRefTarget{sga3:xv:proposition:6-2:item:iii} \emph{\((H_{\bar s})^{0}_{\mathrm{red}}\) is the centralizer
in \((G_{\bar s})^{0}_{\mathrm{red}}\) of a torus of \(G_{\bar s}\)
(respectively, is a Cartan subgroup of \(G_{\bar s}\)).}
\item \emph{\((H_{\bar s})^{0}_{\mathrm{red}}\) is a Borel subgroup
(respectively, a parabolic subgroup) of \(G_{\bar s}\).}
\item \emph{\((H_{\bar s})^{0}_{\mathrm{red}}\) is the radical of
\(G_{\bar s}\) (respectively,
\((G_{\bar s})^{0}_{\mathrm{red}}\) is semisimple).}
\item \emph{\(G_{\bar s}\) is affine and
\((H_{\bar s})^{0}_{\mathrm{red}}\) is the unipotent radical of
\(G_{\bar s}\) (respectively,
\((G_{\bar s})^{0}_{\mathrm{red}}\) is reductive).}
\end{enumerate}

\emph{Then the set \(E\) of points \(s\) of \(S\) for which \(P(s)\)
holds is locally constructible (EGA \(0_{\mathrm{III}}\), 9.1.11).}
\SGARefTarget{sga3:xv:proposition:6:2}

\textbf{\SGAInlineRefTarget{sga3:xv:remark:6-2-1}Remark 6.2.1.} This proposition complements Exposé
VI\textsubscript{B}, \S10. The structure of \(E\) can be described
more precisely by using semicontinuity theorems (see Exposé \SGARefLink{sga3:x:remark:8-7}{X~8.7});
we shall see an example below.

\SGAProof{Proof of \SGARefLink{sga3:xv:proposition:6-2}{6.2}}
Note that if \(S\) is the spectrum of a field, \(E\) is invariant
under extension of that field. A standard reduction (EGA IV~9)
then allows us to reduce to the case where \(S\) is noetherian and
integral, with generic point \(\eta\). One must show that \(E\) or
\(\operatorname{ens}(S)\setminus E\) contains a neighbourhood of
\(\eta\) (EGA IV~9.2.1). One
may assume \(S\) affine of ring \(A\) and field of fractions \(K\). If \(L\) is a finite extension of \(K\), it is immediate that
there exists an \(A\)-subalgebra \(B\) of \(L\), finite over \(A\), having \(L\) as field of fractions. The canonical morphism \(S'
\to S\), where \(S' = \operatorname{Spec} B\), is dominant, of finite presentation, so the image of a non-empty open subset
of \(S'\) contains a non-empty open subset of \(S\) (EGA IV 1.8.4). From the viewpoint that interests us, we may therefore
replace \(S\) by \(S'\), hence replace \(K\) by a finite extension \(L\). Thus we may choose \(L\) so that \((G_{L})_{red}\) and
\((H_{L})_{red}\) are smooth over \(L\) (EGA IV 4.6.6). Possibly restricting \(S'\), we may assume that \(G_{red}\) and
\(H_{\mathrm{red}}\) are group preschemes smooth over \(S\)
(Exposé VI\textsubscript{B}, \S10, and EGA IV~17). In view of the
properties to be proved, we may replace \(G\) and \(H\) by their
reduced identity components (Exposé VI\textsubscript{B}~\SGARefLink{lem:sga3-VIB-10.9}{10.9}), and
hence assume that \(G\) and \(H\) are smooth over \(S\) with
connected fibres. Finally, we may assume that \(H\) is a closed
subgroup prescheme of \(G\) (Exposé VI\textsubscript{B}~10.4).

\emph{Proof of i).} Possibly after a finite extension of \(K\), we
may assume that \(G_{\eta}\) admits a ``Chevalley decomposition'',
i.e. is an extension of an abelian variety \(D_{\eta}\) by a smooth
connected linear algebraic group \(F_{\eta}\) (Séminaire Bourbaki
1956/57, no.~145). By Exposé VI\textsubscript{B}~10.16, there is an
open neighbourhood \(U\) of \(\eta\) over which this generic
extension comes from an extension
\[
  1\longrightarrow F\longrightarrow G|_U\longrightarrow D
  \longrightarrow 1.
\]
One may moreover assume that \(F\) and \(D\) are smooth over \(U\)
with connected fibres, that \(F\) is affine over \(U\), and that
\(D\) is proper over \(U\) (EGA IV~8, 9 and 17). For every point
\(s\) of \(U\), \((D_s,F_s)\) is then the Chevalley decomposition of
\(G_s\). Consequently, \(G_s\) is an abelian variety (respectively,
is affine) if and only if \(F_s\) (respectively, \(D_s\)) is the unit
group; this is a constructible condition (EGA IV~9.2.6.1).

To establish the last two assertions of (i), we may, by the
preceding, assume that \(G\) is affine over \(S\). Let \(q\) be a
prime number invertible on \(S\), and let \({}_qG\) be the
``kernel'' of the \(q\)-th-power morphism in \(G\). It follows
easily from the structure of affine algebraic groups that \(G_s\)
is a torus (respectively, is unipotent) if and only if
\begin{enumerate}
\renewcommand{\labelenumi}{\textup{\alph{enumi})}}
\item \({}_qG_s\) is quasi-finite, which is a constructible
property (EGA IV~9.3.2), and
\item \({}_qG_s\) has \(q^r\) geometric points, where \(r\) is the
relative dimension of \(G\) over \(S\) (respectively,
\({}_qG_s\) has a single point).
\end{enumerate}
The function
\[
 s\longmapsto
 \bigl(\text{number of geometric points of }{}_qG_s\bigr)
\]
is constructible (EGA IV~9.7.9). This completes the proof of (i).

\emph{Proof of iii).} \textbf{a) Case of the centralizer of a
torus.} Suppose
\[
  H_{\eta}=\operatorname{Centr}_{G_{\eta}}(T_{\eta}),
\]
where \(T_{\eta}\) is a torus of \(G_{\eta}\). We shall show that
\(H_s\) is the centralizer of a sub-torus of \(G_s\) for \(s\) in a
neighbourhood of \(\eta\). By (i), after possibly restricting \(S\),
we may assume that \(T_{\eta}\) comes from a sub-torus \(T\) of
\(G\). Then \(Z=\operatorname{Centr}_{G}(T)\) is represented
(Exposé \SGARefLink{sga3:xi:corollary:6:11}{XI~6.11}) by a subgroup prescheme of \(G\). Since \(H\) and
\(Z\) coincide generically, they coincide over a neighbourhood of
\(\eta\). Thus the set \(E\) of points \(s\) of \(S\) for which
\(H_s\) is the centralizer in \(G_s\) of a sub-torus of \(G_s\) is
ind-constructible (EGA IV~1). This will suffice in Lemma~6.6 below
to prove that \(E\) is open in \(S\), and hence, \emph{a fortiori},
locally constructible.

\textbf{b) Case of a Cartan subgroup.} Suppose that \(H_{\eta}\) is
a Cartan subgroup of \(G_{\eta}\). We shall show that \(H_s\) is a
Cartan subgroup of \(G_s\) for every \(s\) in a neighbourhood of
\(\eta\). The group \(H_{\eta}\) is the centralizer in \(G_{\eta}\)
of a torus of \(G_{\eta}\), and is nilpotent (Exposé \SGARefLink{sga3:xii:6.6}{XII~6.6}). By
\SGARefLink{sga3:xii:6.6:item:a}{(a)} and Exposé VI\textsubscript{B}~\SGARefLink{sga3:xv:definition:8-4}{8.4}, \(H_s\) has the same
properties for every \(s\) in some neighbourhood \(U\) of \(\eta\).
For every \(s\in U\), the group \(H_s\) therefore has the same
reductive rank as \(G_s\), and its unique maximal torus is central
(Exposé \SGARefLink{sga3:xii:6.7}{XII~6.7}). Hence \(H_s\) is the centralizer of a maximal torus
of \(G_s\), that is, a Cartan subgroup of \(G_s\).

Suppose now that \(H\) is not a Cartan subgroup of \(G\), and let us show that \(H_{s}\) is not a Cartan subgroup of \(G_{s}\)
for \(s\) in a neighborhood \(U\) of \(\eta\). In view of the assertion provisionally admitted in (a) above, we may restrict
to the case where \(H\) is the centralizer in \(G\) of a sub-torus \(T\). But then \(H_{\eta}\) contains a Cartan subgroup
\(C_{\eta}\) of \(H_{\eta}\). We have just seen that, possibly restricting \(S\), \(C_{\eta}\) extends to a Cartan subgroup \(C\)
of \(G\), which one may assume contained in \(H\). By hypothesis \(H_{\eta}\) strictly contains \(C_{\eta}\), so \(H_{s}\)
strictly contains \(C_{s}\) for \(s\) in a neighborhood \(U\) of \(\eta\) (EGA IV 9.5.2); \emph{a fortiori}, \(H_{s}\) is not

a Cartan subgroup of \(G_{s}\) for \(s\) in \(U\).

\emph{Proof of ii).} Suppose \(H_{\eta}\) is a maximal torus of \(G_{\eta}\) and let \(C_{\eta}\) be its centralizer in \(G_{\eta}\).
By i) and iii), \(H\) is a torus over a neighborhood \(U\) of \(\eta\) and \(C = \operatorname{Centr}_{G|U}(H|U)\) is a Cartan subgroup of
\(G|U\). To prove that \(H\) is a maximal torus of \(G\) over a neighborhood of \(\eta\), we may then replace \(G\) by \(C\), then
by the linear component \(F\) of a Chevalley decomposition of \(C\) (cf.~i)). Let \(q\) be an integer invertible on \(S\),
\({}_{q}F\) the kernel of the \(q\)-th-power morphism in \(F\). Since \(F_{s}\) is affine, nilpotent, smooth and connected,
\(F_{s}\) is the direct product of its maximal torus \(T_{s}\) by a unipotent group (\emph{Bible} 6-04), so
\({}_{q}F_s={}_{q}T_s\). Since \(H_{\eta}\) is a maximal torus,
\({}_{q}H_{\eta}={}_{q}F_{\eta}\), and consequently
\({}_{q}H={}_{q}F\) over a neighborhood \(V\) of \(\eta\).
For every point \(s\) of \(V\), \({}_{q}H_s={}_{q}T_s\), so
\(H_{s} = T_{s}\) is a maximal torus.

Suppose now that \(H_{\eta}\) is not a maximal torus of \(G_{\eta}\). By i), we may restrict to the case where \(H_{\eta}\) is
a torus, then assume that it is contained in a strictly larger torus \(T_{\eta}\). The latter extends to a torus \(T\)
strictly containing \(H\) over a neighborhood \(U\) of \(\eta\). \emph{A fortiori}, \(H_{s}\) is not a maximal torus for \(s \in U\).

\emph{Proof of iv).} Possibly restricting \(S\), we may assume that the center \(Z\) of \(G\) is representable (Exp. VI\textsubscript{B} 10.11)
and flat over \(S\), as well as the quotient \(G/Z\) (\emph{loc. cit.}). The property ``\(H_{s}\) contains \(Z_{s}\)'' is
constructible (EGA IV 9.5.2) and every parabolic subgroup of \(G_{s}\) contains \(Z_{s}\) (Exp. XIV 4.9 a)); this allows us
to replace \(G\) by \(G/Z\), hence assume \(G\) affine over \(S\) (Exp. XII 6.1 and i)). We may further

assume that \(G/H\) is representable, but then \(H_{s}\) is a parabolic subgroup of \(G_{s}\) if and only if \((G/H)_{s}\) is
proper (\emph{Bible} 6 Th. 4 b)), which is an ind-constructible property (EGA IV 9.3.5). So \(E\) is ind-constructible, and
this will suffice for us to prove that \(E\) is open (Lemma 6.6), hence locally constructible.

Let us now examine the case of Borel subgroups. If \(H_{\eta}\) is a Borel subgroup of \(G_{\eta}\), i.e.~a solvable
parabolic subgroup of \(G_{\eta}\), what precedes and Exposé VI\textsubscript{B}, \SGARefLink{prop:sga3-VIB-8.4}{8.4}, entail that these properties remain true at every
point \(s\) of a neighborhood of \(\eta\). If now \(H_{\eta}\) is not a Borel subgroup of \(G_{\eta}\), to prove that the same
holds at points \(s\) of a neighborhood of \(\eta\), we may restrict (in view of what precedes) to the case where \(H_{\eta}\)
is a parabolic subgroup, then assume that \(H_{\eta}\) contains a Borel subgroup \(B_{\eta}\). We have just shown that the
latter extends to a Borel subgroup \(B\) of \(G\) over a neighborhood \(U\) of \(\eta\). Since \(H_{\eta}\) strictly contains
\(B_{\eta}\), \(H_{s}\) strictly contains \(B_{s}\) at every point of an open subset \(V\), and \(H_{s}\) is not a Borel subgroup
of \(G_{s}\) for \(s \in V\).

\emph{Proof of v).} Suppose \(H_{\eta}\) is the radical of \(G_{\eta}\). The group \(H_{\eta}\) is then invariant in \(G_{\eta}\),
solvable (smooth and connected); the same is therefore true for \(H_{s}\) for \(s\) belonging to a neighborhood \(U\) of
\(\eta\) (Exp. VI\textsubscript{B} 8 and 10), so for \(s \in U\), \(H_{s}\) is contained in the radical of \(G_{s}\). Replacing \(G\) by \(G/H\)
(Exp. VI\textsubscript{B} 10), we must

prove that if \(G_{\eta}\) is semisimple, then \(G_{s}\) is semisimple at every point of a neighborhood \(V\) of \(\eta\). Using
i) and ii), one may assume that \(G\) is affine over \(S\) and possesses a maximal torus \(T\). Let \(W\) be the Weyl group of
\(T\) (\SGARefLink{sga3:xii:section:2}{Exp. XII 2}), which is quasi-finite and étale over \(S\), hence finite and étale over an open subset \(V\). It then
follows from the elementary properties of roots (\SGARefLink{sga3:xix:proposition:1-12}{Exp. XIX 1.12}) that \(G\) is semisimple over \(V\).

Suppose now that \(H_{\eta}\) is not the radical of \(G_{\eta}\). Possibly replacing \(K\) by a finite extension \(L\), we may
assume that \(G_{\eta}\) admits a radical \(R_{\eta}\). By what precedes, \(R_{\eta}\) extends to a subgroup prescheme \(R\) of
\(G|_{U}\) such that for every \(s \in U\), \(R_{s}\) is the radical of \(G_{s}\). By hypothesis, \(R_{\eta} \neq H_{\eta}\). So
\(R_{s} \neq H_{s}\) for \(s \in V\). It remains to prove that if \(G_{\eta}\) is not semisimple, neither is \(G_{s}\) at
neighboring points, but this is a particular case of what precedes (take \(H\) = unit group).

\emph{Proof of vi).} The proof is entirely analogous to that of v), in view of i), and is left to the care of the reader.

\textbf{\SGAInlineRefTarget{sga3:xv:corollary:6-3}Corollary 6.3.} \emph{Let \(S_0\) be a quasi-compact
prescheme, let \((S_i)_{i\in L}\) be a projective system of
\(S_0\)-preschemes affine over \(S_0\), and let
\(S=\varprojlim S_i\) (EGA IV~8.2). Let \(G_0\) be a group
prescheme of finite presentation over \(S_0\), set
\[
  G_i=G_0\times_{S_0}S_i,\qquad G=G_0\times_{S_0}S,
\]
and let \(H\) be a subgroup of \(G\). If \(H\) is a sub-torus of
\(G\) (respectively, a maximal sub-torus, a Cartan subgroup, a Borel
subgroup, or a parabolic subgroup), then there exist an index
\(i\in L\) and a subgroup \(H_i\) of \(G_i\) such that
\[
  H=H_i\times_{S_i}S
\]
and \(H_i\) is a sub-torus of \(G_i\) (respectively, \(\ldots\)).}
\SGARefTarget{sga3:xv:corollary:6:3}

Indeed, \(H\) is smooth with connected fibres, hence of finite
presentation over \(S\) (Exposé VI\textsubscript{B}~\SGARefLink{par:sga3-VIA-5.3.3}{5.3.3}). By
Exposé VI\textsubscript{B}, \S10, there exist \(i\in L\) and a
subgroup \(H_i\) of \(G_i\), smooth over \(S_i\), such that
\(H=H_i\times_{S_i}S\). Corollary~6.3 then follows from
Definition~6.1, Proposition~6.2 and EGA IV~9.3.3.

\textbf{Corollary 6.3 bis.} \emph{Let \(S\) be a prescheme and let
\(G\) be an \(S\)-group prescheme of finite presentation. Then the
functions \(\rho_n,\rho_r,\rho_u,\rho_{\mathrm{ab}}\) (cf.~\SGARefLink{sga3:xv:definition:6:1:ter}{6.1 ter})
are locally constructible on \(S\).}
\SGARefTarget{sga3:xv:corollary:6:3:bis}

It suffices to show (EGA IV 9.) that if \(S\) is a noetherian integral scheme with generic point \(\eta\), the functions in
question are constant on a neighborhood of \(\eta\). Possibly replacing \(S\) by a scheme \(S'\) finite over \(S\) and
dominating \(S\), we may assume that \(G_{\eta}\) admits a Cartan subgroup \(C_{\eta}\) with a Chevalley decomposition \(1 \to
L_{\eta} \to C_{\eta} \to A_{\eta} \to 1\). The argument made in \SGARefLink{sga3:xv:proposition:6-2}{6.2} i) proves that this decomposition extends to a
Chevalley decomposition over a neighborhood of \(\eta\):

\[ 1 \longrightarrow L \longrightarrow C \longrightarrow A \longrightarrow 1. \]

Moreover, one may assume that \(C\) is a Cartan subgroup of \(G\) (\SGARefLink{sga3:xv:corollary:6-3}{6.3}) and that the maximal torus \(T_{\eta}\) of \(L_{\eta}\)
extends to a maximal torus \(T\) of \(L\) (\SGARefLink{sga3:xv:corollary:6-3}{6.3}). The corollary follows immediately from this and from the definitions.

\SGARefTarget{sga3:xv:subsection:6:4:0}
\textbf{\SGAInlineRefTarget{sga3:xv:paragraph:6-4-0}6.4.0.} Let \(S\) be a prescheme, \(G\) an \(S\)-prescheme in groups of finite
presentation, and \(\mathcal{L}\) the functor of subgroups of \(G\) that are smooth,
have connected fibres, and are equal to their connected normalizer (cf.~§5).
The functor \(\mathcal{L}\) is representable (\SGARefLink{sga3:xv:theorem:5-2}{5.2} and \SGARefLink{sga3:xv:remark:5:8}{5.8}). The rest of this section
is devoted to certain subfunctors of \(\mathcal{L}\). More precisely, define
\(\mathcal{L}_{C}\) (resp. \(\mathcal{CT}\), resp. \(\mathcal{P}\)) as follows:
for every \(S\)-prescheme \(S'\), the set \(\mathcal{L}_{C}(S')\)
(resp. \(\mathcal{CT}(S')\), resp. \(\mathcal{P}(S')\)) consists of the subgroups
\(H\) of \(G_{S'}\) that are smooth over \(S'\), have connected fibres, are equal
to their connected normalizer, and are such that, for every point \(s'\) of \(S'\),
\(H_{s'}\) contains a Cartan subgroup (\SGARefLink{sga3:xv:definition:6-1}{6.1}) of \(G_{s'}\)
(resp. is the centralizer in \((G_{s'})^{0}_{\mathrm{red}}\) of a subtorus of
\(G_{s'}\), resp. is a parabolic subgroup (\SGARefLink{sga3:xv:definition:6-1}{6.1}) of \(G_{s'}\)).

\SGARefTarget{sga3:xv:theorem:6:4}
\textbf{Theorem 6.4.} \emph{Let \(S\) be a prescheme and \(G\) an \(S\)-prescheme in
groups of finite presentation. The \(S\)-functors \(\mathcal{L}_{C}\),
\(\mathcal{CT}\), and \(\mathcal{P}\) defined in \SGARefLink{sga3:xv:paragraph:6-4-0}{6.4.0} are representable by
\(S\)-preschemes that are quasi-projective and of finite presentation over \(S\).}

\SGARefTarget{sga3:xv:remark:6:5}
\textbf{Remark 6.5.} If \(G\) has smooth fibres---for example, if \(G\) is smooth
over \(S\), or if the residue characteristics of \(S\) are zero
(Exposé VI\textsubscript{B}, \SGARefLink{cor:sga3-VIB-1.6.1}{1.6.1})---then every subgroup \(H\) of \(G\), smooth over
\(S\) with connected fibres, such that \(H_s\) contains a Cartan subgroup of
\((G_s)^0\) for every point \(s\) of \(S\), is necessarily equal to its connected
normalizer (and hence is an element of \(\mathcal{L}_{C}(S)\)). Indeed, by \SGARefLink{sga3:xv:proposition:5-1}{5.1} iii)
we may assume that \(S\) is the spectrum of a field, in which case the property was
noted at the end of the statement of \SGARefLink{sga3:xiii:theorem:2:1}{Exposé XIII, 2.1}.

There are natural monomorphisms
\[
\begin{tikzcd}[column sep=huge,row sep=large]
\mathcal{CT} \arrow[dr,hook] & \\
& \mathcal{L}_{C} \arrow[r,hook] & \mathcal{L}.\\
\mathcal{P} \arrow[ur,hook] &
\end{tikzcd}
\]

Let us show that these monomorphisms are open immersions of finite presentation (which will already prove, in view of
\SGARefLink{sga3:xv:theorem:5-2}{5.2}, that \(\mathcal{L}_{C}\), \(\mathcal{CT}\), \(\mathcal{P}\) are representable by \(S\)-preschemes, union of an increasing
sequence of open subpreschemes, quasi-projective and of finite presentation over \(S\)). Now this will follow, by the
usual technique, from the following lemma:

\SGARefTarget{sga3:xv:lemma:6:6}
\textbf{Lemma 6.6.} \emph{Let \(S\) be a prescheme, \(G\) an \(S\)-prescheme in
groups of finite presentation, and \(H\) a subgroup of \(G\), smooth with connected
fibres. The set of points \(s\) of \(S\) such that \(H_s\) contains a Cartan subgroup
of \(G_s\) (resp. is equal to the centralizer in \((G_s)^0_{\mathrm{red}}\) of a torus
of \(G_s\)) is open in \(S\). If, moreover, \(H\) is equal to its connected
normalizer,\footnote{This hypothesis is in fact superfluous; see Exposé XVII,
Appendix III, 3.} then the set of points \(s\) of \(S\) such that \(H_s\) is a
parabolic subgroup of \(G_s\) is also open in \(S\).}

The assertion to be proved is local on \(S\), which allows us to assume \(S\) affine, then, by EGA IV 8.1 and 6.3, \(S\)
noetherian. Denote by \(E\) the set of points of \(S\) having the property in question. It then follows from the assertions
effectively proved in 6.2 that \(E\) is ind-constructible. But \(S\) is noetherian, so to prove that \(E\) is open, it
suffices to show that \(E\) is stable under generizations (EGA IV 1.10.1). Using EGA II 7.1.9, we are finally led to prove
that if \(S\) is the spectrum of a discrete valuation ring, and if the closed point \(s\) belongs to \(E\), then so does the
generic point \(t\).

We shall need the following lemma:

\SGARefTarget{sga3:xv:lemma:6:7}
\textbf{Lemma 6.7.} \emph{Let \(A\) be a complete noetherian local ring,
\(S=\operatorname{Spec}A\), \(s\) the closed point of \(S\), \(H\) an
\(S\)-prescheme in groups, smooth with connected fibres, and \(T_s\) a subtorus of
\(H_s\). Then:
\begin{enumerate}
\renewcommand{\labelenumi}{\roman{enumi})}
\item there exists a closed subgroup prescheme \(C\) of \(H\), smooth with connected
fibres, such that \(C_s=\operatorname{Centr}_{H_s}(T_s)\);
\item for every point \(t\) of \(S\), \(C_t\) is the centralizer in \(H_t\) of a
subtorus \(T_t\) of \(H_t\).
\end{enumerate}}

\emph{Proof of \SGARefLink{sga3:xv:lemma:6:7}{6.7}.} Let \(T'\) be an \(S\)-torus for which there is an isomorphism
\(u_0:T'_s\xrightarrow{\sim}T_s\) (\SGARefLink{sga3:x:corollary:4-6}{Exposé X, 4.6}). Let \(\mathfrak m\) be the
maximal ideal of \(A\), and put \(A_n=A/\mathfrak m^n\),
\(S_n=\operatorname{Spec}A_n\), \(H_n=H\times_S S_n\), and so on. Since \(H\) is
smooth over \(S\), for every integer \(n>0\) there is an \(S_n\)-group morphism
\[
u_n:T'_n\longrightarrow H_n
\]
lifting \(u_0\) (\SGARefLink{sga3:ix:theorem:3:6}{Exposé IX, 3.6}), and by induction on \(n\) we may assume that
\(u_n\) lifts \(u_{n-1}\). Let \(q\) be a prime invertible on \(S\). For every
integer \(\ell\) that is a power of \(q\), write \({}_{\ell}u_n\) for the restriction
of \(u_n\) to \({}_{\ell}T'_n\). For fixed \(\ell\) and varying \(n\), the morphisms
\({}_{\ell}u_n\) form a projective system and therefore arise from a unique
\(S\)-group morphism
\[
{}_{\ell}u:{}_{\ell}T'\longrightarrow H
\]
(\SGARefLink{sga3:xv:proposition:1-6}{1.6} a)). Since \(H\) is separated (Exposé VI\textsubscript{B}, \SGARefLink{prop:sga3-VIB-5.2}{5.2}) and \(u_0\) is
a monomorphism, \({}_{\ell}u\) is a monomorphism (\SGARefLink{sga3:ix:theorem:6:8}{Exposé IX, 6.8}), and even a closed
immersion because it is finite, \({}_{\ell}T'\) being finite over \(S\). Denote its
image by \(M(\ell)\). The family of subgroups of multiplicative type \(M(\ell)\) is
clearly coherent in the sense of \SGARefLink{sga3:xv:theorem:4-1}{4.1}. Put
\[
C_\ell=\operatorname{Centr}_H(M(\ell)).
\]
This is represented by a closed subgroup prescheme (\SGARefLink{sga3:xv:lemma:2:5}{2.5}), since \(H\) is separated,
and is smooth over \(S\) (\SGARefLink{sga3:xi:corollary:2:4}{Exposé XI, 2.4}). The \(C_\ell\) form a filtered decreasing
family of closed subschemes and are therefore stationary, since \(H\) is noetherian.
Their stationary value is a smooth closed subgroup \(C\) such that
\(C_s=\operatorname{Centr}_{H_s}(T_s)\), by the density theorem (\SGARefLink{sga3:ix:theorem:4:7}{Exposé IX, 4.7}).
It remains to show that for every point \(t\) of \(S\), \(C_t\) is the centralizer
in \(H_t\) of a subtorus \(T_t\); this will also imply that \(C_t\) is connected.
This follows from the more precise lemma below, applied to the family
\(M(\ell)_t\) of subgroups of \(H_t\).

\SGARefTarget{sga3:xv:lemma:6:8}
\textbf{Lemma 6.8.} \emph{Let \(G\) be a connected algebraic group over a field
\(k\), let \(r>0\) be an integer, and let \(q\) be an integer prime to the
characteristic\footnote{Editor's note: the word ``residual'' has been removed.}
of \(k\). Let \(M(\ell)\), where \(\ell\) ranges over the powers of \(q\), be a
coherent family of subgroups of \(G\), of multiplicative type and of type
\((\mathbb Z/\ell\mathbb Z)^r\) (cf.~\SGARefLink{sga3:xv:lemma:4:6}{4.6}). Let \(M\) be the algebraic subgroup of
\(G\) generated by the \(M(\ell)\) (\emph{loc.~cit.}), and let \(T\) be the unique
maximal torus of \(M\) (cf.~\SGARefLink{sga3:xv:lemma:3:4}{3.4}). Then, for all sufficiently large \(\ell\),
\[
\operatorname{Centr}_G(T)=\operatorname{Centr}_G(M)
  =\operatorname{Centr}_G(M(\ell)).
\]}

The last equality is clear. To prove the first, introduce the centre \(Z\) of
\(G\), put \(G'=G/Z\), let \(M'\) (resp. \(T'\)) be the image of \(M\) (resp.
\(T\)) in \(G'\), and let \(K\) be the inverse image of \(T'\) in \(G\), that is,
the algebraic subgroup of \(G\) generated by \(T\) and \(Z\). It plainly suffices
to prove \(K\supset M\), or equivalently \(T'=M'\). Now \(M\) is smooth and
connected (\SGARefLink{sga3:xv:lemma:4:6}{4.6}), while \(G'\) is affine (\SGARefLink{sga3:xii:6.1}{Exposé XII, 6.1}); hence \(M'\) is the
direct product of its maximal torus \(T'\) (\SGARefLink{sga3:xii:6.6}{Exposé XII, 6.6} d)) and a unipotent
group (\emph{Bible} 4, Th.~4), after assuming \(k\) algebraically closed. The image
of \(M(\ell)\) in \(M'\) is therefore contained in \(T'\). Thus the inverse image
of \(T'\) in \(M\) contains \(M(\ell)\) for every \(\ell\), and consequently is
equal to \(M\). Hence \(M'=T'\). This proves \SGARefLink{sga3:xv:lemma:6:8}{6.8}, and therefore \SGARefLink{sga3:xv:lemma:6:7}{6.7}.

We now prove \SGARefLink{sga3:xv:lemma:6:6}{6.6}. We have reduced to the case where \(S\) is the spectrum of a
discrete valuation ring \(A\), which we may moreover assume complete with
algebraically closed residue field. After replacing \(A\), if necessary, by its
normalization in a finite extension of its fraction field, we may assume that
\((G_t)_{\mathrm{red}}\) is smooth.\footnote{Editor's note: details or references
should be supplied here.} To prove \SGARefLink{sga3:xv:lemma:6:6}{6.6} we may clearly replace \(G\) by the identity
component of the schematic closure in \(G\) of \((G_t)^0_{\mathrm{red}}\), and
therefore assume that \(G\) is flat over \(S\), has connected fibres, and has smooth
generic fibre \(G_t\).

\begin{enumerate}
\renewcommand{\labelenumi}{\alph{enumi})}
\item Suppose that \(H_s\) is the centralizer in \((G_s)_{\mathrm{red}}\) of a torus
\(T_s\). We show that \(H_t\) is then the centralizer in \(G_t\) of a subtorus of
\(G_t\). By \SGARefLink{sga3:xv:lemma:6:7}{Lemma 6.7} there is a subgroup scheme \(C\) of \(H\), smooth over \(S\),
whose closed fibre is \(H_s\), such that
\(C_t=\operatorname{Centr}_{H_t}(T_t)\) for a subtorus \(T_t\) of \(H_t\). Since
\(H\) is smooth over \(S\) with connected fibres, dimension considerations give
\(C=H\). With the notation of \SGARefLink{sga3:xv:lemma:6:7}{6.7}, \(H=C=C_\ell\) for all sufficiently large
\(\ell\). Similarly, put
\(C'_\ell=\operatorname{Centr}_G(M(\ell))\) (\SGARefLink{sga3:xv:lemma:2:5}{2.5}), and let \(C'\) be the stationary
value of the \(C'_\ell\) for large \(\ell\) (\SGARefLink{sga3:xv:lemma:2-5-bis}{2.5 bis}). The group scheme \(C'\)
contains \(H\), and
\[
C'_t=\operatorname{Centr}_{G_t}(T_t)\quad\text{(6.8)},\qquad
C'_s=\operatorname{Centr}_{G_s}(T_s).
\]
The hypothesis on \(H_s\) gives \(\dim H_s=\dim C'_s\). Moreover,
\(\dim H_s=\dim H_t\), since \(H\) is smooth over \(S\), and
\(\dim C'_t\leq \dim C'_s\) (Exposé VI\textsubscript{B}, \SGARefLink{prop:sga3-VIB-4.1}{4.1}). Hence
\(\dim H_t=\dim C'_t\). Since \(G_t\) is smooth and connected,
\(C'_t=\operatorname{Centr}_{G_t}(T_t)\) is smooth and connected; finally,
\[
H_t=C'_t=\operatorname{Centr}_{G_t}(T_t).
\]

\item Suppose that \(H_s\) contains a Cartan subgroup \(C_s\) of \(G_s\), that is,
the centralizer in \((G_s)_{\mathrm{red}}\) of a maximal torus \(T_s\) of \(G_s\).
By \SGARefLink{sga3:xv:lemma:6:7}{6.7} there is a subgroup scheme \(C\) of \(H\), smooth over \(S\) with connected
fibres, lifting \(C_s\). It suffices to prove that \(C_t\) contains a Cartan subgroup
of \(G_t\). By part a), applied with \(H=C\), \(C_t\) is the centralizer in \(G_t\)
of a subtorus of \(G_t\), and hence contains a Cartan subgroup of \(G_t\).

\item Suppose that \(H_s\) is a parabolic subgroup of \(G_s\). Let
\(N=\operatorname{Norm}_G(H)\). By hypothesis \(H\) is equal to its connected
normalizer, so \(N_s\) is smooth and consequently
\[
N_s=\operatorname{Norm}_{(G_s)_{\mathrm{red}}}(H_s)=H_s
\qquad\text{(\SGARefLink{sga3:xii:7.10}{Exposé XII, 8 bis}).}
\]
Thus \(N\) is flat over \(S\). We shall see in Exposé XVI that under these
conditions \(G/N\) is representable. Since \(H_s=N_s\) is a parabolic subgroup,
\((G/N)_s\) is proper. Since \(G/N\) is flat over \(S\) with connected fibres,
EGA III, 5.5.1 implies that \(G/N\) is proper over \(S\). Thus
\((G/N)_t=G_t/N_t\) is proper over \(t\), and so is \(G_t/H_t\), because
\(N_t/H_t\) is finite. It remains only to apply the following lemma.
\end{enumerate}

\SGARefTarget{sga3:xv:lemma:6:9}
\textbf{Lemma 6.9.} \emph{Let \(k\) be a field, \(G\) a \(k\)-algebraic group,
\(H\) a smooth connected algebraic subgroup of \(G\), and
\(N=\operatorname{Norm}_G(H)\). If \(\dim H=\dim N\) and \(G/H\) is proper, then
\(H\) is a parabolic subgroup of \(G\).}

Indeed, we may assume that \(k\) is algebraically closed and that \(G\) is smooth
and connected. The centre \(Z\) of \(G\) is contained in \(N\), and
\(\dim H=\dim N\) implies that \(Z'=(Z)^0_{\mathrm{red}}\) is contained in \(H\).
Hence \(G'=G/Z'\) is affine (\SGARefLink{sga3:xii:6.1}{Exposé XII, 6.1}). Replacing \(G\) by \(G'\) and \(H\)
by its image in \(G'\), we reduce to the case in which \(G\) is affine
(\SGARefLink{sga3:xiv:4.9}{Exposé XIV, 4.9}), and \SGARefLink{sga3:xv:lemma:6:9}{Lemma 6.9} then follows from \emph{Bible} 6, Th.~4.
This proves \SGARefLink{sga3:xv:lemma:6:6}{Lemma 6.6}.

To complete the proof of \SGARefLink{sga3:xv:theorem:6:4}{6.4}, we must show that the \(S\)-preschemes representing
\(\mathcal L_C\) (respectively \(\mathcal{CT}\) and \(\mathcal P\)) are of finite
presentation over \(S\). This assertion is local on \(S\), so we may assume first
that \(S\) is affine and then, since \(G\) is of finite presentation, that \(S\) is
noetherian (EGA IV, 8.9). We have just seen that the natural inclusions
\(\mathcal{CT}\to\mathcal L\) and \(\mathcal P\to\mathcal L\) are immersions.
The same is therefore true of
\(\mathcal{CT}\to\mathcal L_C\) and \(\mathcal P\to\mathcal L_C\).
Consequently, it is enough to prove that \(\mathcal L_C\) is represented by an
\(S\)-prescheme of finite presentation.

Resume the notation introduced in \SGARefLink{sga3:xv:theorem:5-2}{5.2}. For every integer \(n\geq 0\), let
\(\mathcal L^n\) be the subfunctor of \(\mathcal L\) defined by
\[
\mathcal L^n(S')
 =\left\{H\in\mathcal L(S')\ \middle|\ 
   \operatorname{Norm}_{G_{S'}}(H)
   =\operatorname{Norm}_{G_{S'}}\bigl(H^{(n)}\bigr)\right\}.
\]
The \(S\)-functor \(\mathcal L^n\) is thus represented by an open subprescheme
of \(\mathcal L\), the disjoint union of the \(\mathcal L_r^n\).
Each \(\mathcal L_r^n\) is of finite presentation over \(S\) (\SGARefLink{sga3:xv:lemma:5:3}{5.3}), and is empty
for
\[
r>\sup_{s\in S}\dim G_s,
\]
which is a finite number because \(S\) is quasi-compact. Hence \(\mathcal L^n\)
is of finite presentation over \(S\). It remains to prove that
\(\mathcal L_C\) is contained in \(\mathcal L^n\) for all sufficiently large
\(n\).

For each point \(s\) of \(S\), let \(d(s)\) be the least integer \(n\), finite or
infinite, such that
\[
\mathcal L_{C,G_s}\subset\mathcal L^n_{G_s}.
\]
It is enough to show that \(d\) is bounded on \(S\). Indeed, if \(M\) is an
upper bound, then \(\mathcal L_C\) is set-theoretically contained in
\(\mathcal L^M\), and hence is contained in \(\mathcal L^M\), since the latter
is open in \(\mathcal L\). A direct constructibility argument reduces us to
showing that, when \(S\) is noetherian and integral with generic point \(t\),
the function \(d\) is bounded on a neighbourhood of \(t\).

\begin{enumerate}
\renewcommand{\labelenumi}{\alph{enumi})}
\item \textbf{Reduction to the case in which \(G\) is smooth over \(S\).}
Proceeding as in \SGARefLink{sga3:xv:proposition:6-2}{6.2}, we see that, after possibly changing \(S\), we may assume
that \(G_{\mathrm{red}}\) is a group prescheme smooth over \(S\); denote it by
\(G'\). Put
\[
X=\mathcal L_{C,G},\qquad X'=\mathcal L_{C,G'},
\]
and let \(H\) (respectively \(H'\)) be the subgroup prescheme of \(G_X\)
(respectively \(G'_{X'}\)) universal for the functor
\(\mathcal L_{C,G}\) (respectively \(\mathcal L_{C,G'}\)).
Since \(H\) is smooth over \(X\), the prescheme
\(H\times_X X_{\mathrm{red}}\) is reduced, and is therefore contained in
\(G'_{X_{\mathrm{red}}}\). It is an element of
\(\mathcal L_{C,G'}(X_{\mathrm{red}})\), and hence defines a canonical
morphism
\[
p:X_{\mathrm{red}}\longrightarrow X'.
\]
It is clear that \(p\) is a monomorphism; we show that it is in fact an
immersion. Let \(N'\) be the normalizer of \(H'\) in \(G_{X'}\).
The set of points \(s\in X'\) for which
\(H'_s\to N'_s\) is an open immersion is an open subset \(U\), and
\[
H'|_U\longrightarrow N'|_U
\]
is an open immersion (Exposé VI\textsubscript{B}, 2.5, and EGA IV, 17.9.5).
It follows from 5.1(iii) that \(U\) is the largest open subset of \(X'\)
over which \(H'\) equals its connected normalizer in \(G_{X'}\).
Thus \(H'|_U\in\mathcal L_{C,G}(U)\), and it follows at once that \(p\)
is an isomorphism of \(X_{\mathrm{red}}\) onto \(U_{\mathrm{red}}\).
If \(\mathcal L_{C,G}\) is known to be of finite type when \(G\) is smooth,
then \(X'\) is of finite type over \(S\); consequently
\(X_{\mathrm{red}}\) is of finite type (since \(S\) is noetherian), and is
therefore contained in \(\mathcal L_G^M\) for all sufficiently large \(M\).
The same then holds for \(X\).

\item \textbf{The case in which \(S\) is the spectrum of an algebraically
closed field \(k\) of characteristic \(p>0\), and \(G\) is a smooth
algebraic group over \(k\).}
Instead of the infinitesimal neighbourhoods \(H^{(n)}\) of a subgroup
prescheme \(H\) of \(G\), we use the radicial subgroups \(F_n(H)\), the
kernels of the iterates of the Frobenius morphism of \(H\)
(Exposé VII\textsubscript{A}, 4). This is legitimate here because \(F_n(H)\)
is contained in the infinitesimal neighbourhood of order \(p^n\) of the
unit section of \(H\). If \(T\) is a subtorus of \(G\), then
\[
F_n(T)={}_{p^n}T,
\]
and duality immediately shows that
\(\operatorname{Centr}_G({}_{p^n}T)=\operatorname{Centr}_G(T)\) for all
sufficiently large \(n\). We have the following more precise proposition.
\end{enumerate}

\SGARefTarget{sga3:xv:proposition:6:10}
\textbf{Proposition 6.10.} \emph{Let \(k\) be a field of characteristic
\(p>0\), let \(G\) be a smooth \(k\)-algebraic group, let \(T\) be a maximal
torus of \(G_{\bar k}\), and let \(m\) be the least integer such that
\[
\operatorname{Centr}_G({}_{p^m}T)^0
 =\operatorname{Centr}_G(T)^0.
\]
Then, for every \(k\)-prescheme \(S\) and every
\(H\in\mathcal L_C(S)\), one has
\[
\operatorname{Norm}_{G_S}(H)
 =\operatorname{Norm}_{G_S}\bigl(F_m(H)\bigr).
\]
In particular, \(\mathcal L_C\) is contained in \(\mathcal L^{p^m}\).}

Since \(H\) is smooth and of finite presentation over \(S\), and has constant
nilpotent rank, namely that of \(G\), we shall see in the next section (\SGARefLink{sga3:xv:theorem:7-3}{7.3})
that the functor \(\mathcal C_H\) of Cartan subgroups of \(H\) is represented
by an \(S\)-prescheme smooth over \(S\). The reader will verify that the proof
given there does not use the finite-type property of \(\mathcal L_{C,H}\).
It follows from \SGARefLink{sga3:xiii:theorem:3:1}{Exposé XIII, 3.1}, that we may consider the open subset \(U\)
of regular points of \(H\).

Let \(S'\) be an \(S\)-prescheme and let \(g\in G(S')\) normalize
\(F_m(H)_{S'}\). To prove that \(g\) normalizes \(H_{S'}\), it is enough to
prove that
\[
\operatorname{int}(g)U_{S'}\subset H_{S'}.
\]
Indeed, \(H_{S'}\cap\operatorname{int}(g)H_{S'}\) would then contain an open
subgroup of \(\operatorname{int}(g)H_{S'}\), and hence would equal it, since
the latter has connected fibres. After replacing \(S'\) by a suitable
covering \(S''\), and then replacing \(S''\) by \(S\), we are reduced to
proving the following: if \(g\in G(S)\) normalizes \(F_m(H)\) and
\(u\in U(S)\), then \(\operatorname{int}(g)u\in H(S)\).

Let \(C\) be the unique Cartan subgroup of \(H\) ``containing'' \(u\)
(\SGARefLink{sga3:xiii:corollary:3:2}{Exposé XIII, 3.2}). It is enough to show that
\[
C'=\operatorname{int}(g)C\subset H.
\]
Since \(H\in\mathcal L_{C,G}(S)\), the group \(C\) is also a Cartan subgroup
of \(G\). The latter admits maximal tori, so \(C\) is the centralizer in
\(G_S^0\) of its unique maximal torus \(T\) (\SGARefLink{sga3:xii:7.1:item:a}{Exposé XII, 7.1(a)}).
By the definition of \(m\), and because any two maximal tori of \(G\) are
locally conjugate for the fpqc topology (\SGARefLink{sga3:xii:7.1}{Exposé XII, 7.1}),
\[
C=\operatorname{Centr}_G(T)^0
 =\operatorname{Centr}_G({}_{p^m}T)^0.
\]
Conjugating by \(g\), we obtain
\[
C'=\operatorname{Centr}_G\bigl(\operatorname{int}(g)({}_{p^m}T)\bigr)^0.
\]
Now \(\operatorname{int}(g)({}_{p^m}T)\) is a subgroup of multiplicative
type of
\[
\operatorname{int}(g)\bigl(F_m(H)\bigr)=F_m(H),
\]
because \(g\) normalizes \(F_m(H)\); it is therefore contained in \(H\).
\SGARefLink{sga3:xiii:theorem:2:1}{Exposé XIII, 2.1} (proved there when the base is a field, but extending at
once to an arbitrary base) then shows that, in order to prove
\(C'\subset H\), it is enough to prove
\[
\operatorname{Lie}C'\subset\operatorname{Lie}H.
\]
But
\[
\operatorname{Lie}C'
 =\operatorname{int}(g)(\operatorname{Lie}C)
 \subset\operatorname{int}(g)(\operatorname{Lie}H).
\]
On the other hand, we may assume \(m\geq1\), and Exposé
VII\textsubscript{A}, \SGARefLink{sga3:VIIA:4.1.2}{4.1.2}, gives
\[
\operatorname{Lie}H
 =\operatorname{Lie}\bigl(F_m(H)\bigr)
 =\operatorname{Lie}\bigl(\operatorname{int}(g)(F_m(H))\bigr)
 =\operatorname{int}(g)(\operatorname{Lie}H).
\]
Thus \(\operatorname{Lie}C'\subset\operatorname{Lie}H\), which completes
the proof of \SGARefLink{sga3:xv:proposition:6:10}{6.10}.

We shall need another description of the integer \(m\) introduced in \SGARefLink{sga3:xv:proposition:6:10}{6.10}.

\SGARefTarget{sga3:xv:lemma:6:11}
\textbf{Lemma 6.11.} \emph{Let \(G\) be a smooth algebraic group over an
algebraically closed field \(k\) of characteristic \(p>0\), let \(T\) be a
maximal torus of \(G\), put \(\mathfrak g=\operatorname{Lie}G\), and let \(R\)
be the set of nonzero characters of \(T\) occurring in the representation of
\(T\) on \(\mathfrak g\) induced by the adjoint representation of \(G\).
For \(r\in R\), let \(e_r\) be the greatest integer \(n\) such that \(p^n\)
divides \(r\) in the character group of \(T\). Then
\[
m=\sup_{r\in R}(e_r+1)
\]
if \(R\neq\varnothing\), and \(m=0\) otherwise.}

Indeed, \(\operatorname{Centr}_G(T)\) is smooth and is contained in
\(\operatorname{Centr}_G({}_{p^m}T)\). Hence
\[
\begin{aligned}
\operatorname{Centr}_G(T)^0
 &=\operatorname{Centr}_G({}_{p^m}T)^0\\
&\Longleftrightarrow\
\operatorname{Lie}\bigl(\operatorname{Centr}_G(T)\bigr)
 =\operatorname{Lie}\bigl(\operatorname{Centr}_G(F_m(T))\bigr)\\
&\Longleftrightarrow\
\mathfrak g^T=\mathfrak g^{{}_{p^m}T}
\qquad\text{(Exposé II, 5.2.3).}
\end{aligned}
\]
With the usual notation,
\[
\mathfrak g=\mathfrak g_0\oplus\bigoplus_{r\in R}\mathfrak g_r.
\]
Thus \(\mathfrak g^T=\mathfrak g_0\), whereas
\[
\mathfrak g^{{}_{p^m}T}
 =\mathfrak g_0\oplus\bigoplus_{r\in R'}\mathfrak g_r,
\]
where \(R'\) is the subset of \(R\) formed by the characters whose restriction
to \({}_{p^m}T\) is zero. A nonzero character \(r\) of \(T\) has trivial
restriction to \({}_{p^m}T\) if and only if \(m\leq e_r\), which proves the
lemma.

\begin{enumerate}
\renewcommand{\labelenumi}{\alph{enumi})}
\setcounter{enumi}{2}
\item \textbf{Return to the proof of 6.4.}
By part a) and the paragraph preceding it, we have reduced to the case in
which \(S\) is a noetherian integral scheme and \(G\) is smooth over \(S\).
We must show that \(d\) is bounded on a neighbourhood of the generic point
\(t\) of \(S\). After possibly changing \(S\), we may assume that \(G\) has
a split maximal torus (\SGARefLink{sga3:xv:proposition:6-2}{6.2})
\[
T=D_S(M).
\]
Let
\[
\mathfrak g=\bigoplus_{\lambda\in M}\mathfrak g^\lambda
\]
be the decomposition of the Lie algebra of \(G\) according to the characters
of \(T\), and let \(R\) be the finite set of nonzero \(\lambda\in M\) for
which \(\mathfrak g^\lambda\neq0\). We distinguish two cases.
\end{enumerate}

\textbf{First case: \(t\), and hence every point of \(S\), has positive
residue characteristic \(p\).}
It is clear from the foregoing that \(d\) is then bounded by \(p^m\), where
\(m\) is defined as in \SGARefLink{sga3:xv:lemma:6:11}{6.11}.

\textbf{Second case: \(t\) has residue characteristic zero.}
For each \(\lambda\in R\), let \(n_\lambda\) be the greatest integer dividing
\(\lambda\) in the group \(M\), and put
\[
n=\prod_{\lambda\in R}n_\lambda.
\]
For every prime number \(q\) dividing \(n\), let \(S_q\) be the closed subset
of \(S\) formed by the points of residue characteristic \(q\), and let \(U\)
be the nonempty open subset, containing \(t\), complementary to the union of
the \(S_q\). If \(s\in U\), then either \(s\) has residue characteristic zero,
in which case \(d(s)=1\) by \SGARefLink{sga3:xv:remark:5:6}{5.6}, or \(s\) has residue characteristic
\(p>0\) not dividing \(n\). In the latter case the integer \(m\) attached to
\(G_s\) in \SGARefLink{sga3:xv:lemma:6:11}{6.11} is at most one. Moreover, Exposé
VII\footnote{Editor's note: the argument should be made explicit.} gives
\[
\operatorname{Norm}_{G_{S'}}\bigl(F_1(H)\bigr)
 =\operatorname{Norm}_{G_{S'}}(\operatorname{Lie}H)
 =\operatorname{Norm}_{G_{S'}}\bigl(H^{(1)}\bigr).
\]
Finally, \SGARefLink{sga3:xv:proposition:6:10}{6.10} shows that, for
\(H\in\mathcal L_{C,G_s}(S')\),
\[
\operatorname{Norm}_{G_{S'}}(H)
 =\operatorname{Norm}_{G_{S'}}\bigl(H^{(1)}\bigr).
\]
Consequently \(d(s)\leq1\), so \(d\) is bounded by \(1\) on \(U\).
This completes the proof of \SGARefLink{sga3:xv:theorem:6:4}{6.4}.

\SGARefTarget{sga3:xv:corollary:6:12}
\textbf{Corollary 6.12.} \emph{Let \(S\) be a prescheme and \(G\) an
\(S\)-prescheme in groups of finite presentation. Suppose that the nilpotent
rank of the fibres of \(G\) (respectively the dimension of their Borel
subgroups) is locally constant on \(S\). Then the functor \(\mathcal C\) of
Cartan subgroups of \(G\) (respectively the functor \(\mathcal B\) of Borel
subgroups of \(G\)) is represented by an \(S\)-prescheme of finite
presentation over \(S\).}

Indeed, after restricting \(S\), we may assume that the nilpotent rank
\(\nu\) of the fibres is constant. Then \(\mathcal C\) is represented by the
open-and-closed subprescheme of the prescheme \(X\) representing
\(\mathcal L_C\) (\SGARefLink{sga3:xv:theorem:6:4}{6.4}) over which the universal subgroup of \(G_X\) for the
functor \(\mathcal L_C\) has relative dimension \(\nu\). The proof is
analogous for \(\mathcal B\), using the representability of \(\mathcal P\).

% END inputs/XV/tex/components/06_expose_XV_section6.tex

% BEGIN inputs/XV/tex/components/07_expose_XV_section7.tex
% Authority: Expo15.pdf, local pp. 50--56 / combined pp. 920--926 /
% printed pp. 445--458.
% Stops immediately before §8.

\SGASection{7}{Cartan subgroups of a smooth group}
\SGARefTarget{sga3:xv:section:7}

\SGAStatement{\SGAInlineRefTarget{sga3:xv:proposition:7-1}Proposition 7.1}{Let \(S\) be a prescheme and let \(G\)
be an \(S\)-group prescheme of finite presentation, smooth over \(S\),
with connected fibres.
\begin{enumerate}
\renewcommand{\labelenumi}{\textup{\roman{enumi})}}
\item\SGAInlineRefTarget{sga3:xv:proposition:7-1:item:i} Let \(\mathcal{CT}\) be the \(S\)-functor defined by
\[
\mathcal{CT}(S')=
\left\{
\begin{array}{c}
\text{subgroup preschemes \(H\) of \(G_{S'}\), smooth over \(S'\),}\\
\text{such that every \(H_{s'}\) is the centralizer in \(G_{s'}\)}\\
\text{of a sub-torus of \(G_{s'}\)}
\end{array}
\right\}.
\]
Then \(\mathcal{CT}\) is represented by an \(S\)-prescheme of finite
presentation, smooth and quasi-projective over \(S\).
\item If \(S\) is artinian and \(H\in\mathcal{CT}(S)\), then \(H\)
is the centralizer in \(G\) of a sub-torus of \(G\). If \(S\) is the
spectrum of a henselian local ring and \(H\in\mathcal{CT}(S)\), then
\(H\) is the centralizer in \(G\) of a subgroup of multiplicative
type that is étale over \(S\).
\item\SGAInlineRefTarget{sga3:xv:proposition:7-1:item:iii} Let \(L\) be an \(S\)-group prescheme smooth and of finite
presentation, let \(i:L\to G\) be an \(S\)-group monomorphism, and
let \(H\in\mathcal{CT}(S)\). Then
\(\operatorname{Transp}_G(H,L)\) and
\(\operatorname{Transpstr}_G(H,L)\)
(Exposé \SGARefLink{sga3:viii:examples:6:5:item:e:01}{VIII~6.5(e)}) are represented by closed subpreschemes of \(G\),
smooth over \(S\).
\item If \(H\in\mathcal{CT}(S)\), then \(H\) is closed in \(G\);
\(N=\operatorname{Norm}_G(H)\) is represented by a closed subgroup
prescheme of \(G\), smooth over \(S\); \(N/H\) is represented by an
\(S\)-group prescheme, separated, étale, and of finite type over
\(S\); and \(G/N\) is represented by an \(S\)-prescheme smooth and
quasi-projective over \(S\).
\item\SGAInlineRefTarget{sga3:xv:proposition:7-1:item:v} Let \(G'\) be an \(S\)-group prescheme of finite presentation
and let \(u:G\to G'\) be faithfully flat. Then \(G'\) satisfies the
same hypotheses as \(G\) (Exposé VI\textsubscript{B}~9). If
\(H\in\mathcal{CT}_G(S)\), its image is represented by a subgroup
prescheme \(H'\in\mathcal{CT}_{G'}(S)\), and \(H\to H'\) is faithfully
flat. If \(H=\operatorname{Centr}_G(T)\) for a torus \(T\), then
\[
 H'=\operatorname{Centr}_{G'}(T'),\qquad T'=u(T).
\]
\item\SGAInlineRefTarget{sga3:xv:proposition:7-1:item:vi} Under the hypotheses of (v), the \(S\)-morphism
\[
 \widetilde u:\mathcal{CT}_G\longrightarrow\mathcal{CT}_{G'},
 \qquad H\longmapsto H'=u(H),
\]
is quasi-compact and faithfully flat. If, moreover, \(\ker u\) is
central, then \(\widetilde u\) is an isomorphism, with inverse
\(H'\mapsto u^{-1}(H')\).
\end{enumerate}}
\SGARefTarget{sga3:xv:proposition:7:1}

\medskip\noindent\textit{Proof of (ii).}
For the first assertion we may suppose that \(S\) is local artinian,
with closed point \(s\). Let \(H\in\mathcal{CT}(S)\), and let \(T_s\)
be the maximal central torus of \(H_s\), which is defined over
\(\kappa(s)\) (cf. \SGARefLink{sga3:xv:lemma:3:4}{Lemma 3.4}). Since
\[
  H_s=\operatorname{Centr}_{G_s}(T_s),
\]
and \(H\) is smooth, \(T_s\) lifts uniquely to a sub-torus \(T\) of
\(H\), central in \(H\) (Exposé \SGARefLink{sga3:ix:theorem:3:6:bis}{IX~3.6 bis} and \SGARefLink{sga3:xv:remark:5:6}{5.6}). Then
\(H'=\operatorname{Centr}_G(T)\) contains \(H\) and has the same
fibre as \(H\). Because \(H\) is flat, \(H=H'\) (Exposé
VI\textsubscript{B}~\SGARefLink{sga3:xv:lemma:2:5}{2.5}).

Suppose now that \(S\) is the spectrum of a henselian local ring,
which the usual reductions allow us to take noetherian. Let \(s\) be
its closed point, let \(T_s\) be the maximal central torus of \(H_s\),
let \(q\) be invertible on \(S\), let \(\ell\) be a power of \(q\),
and let \(T\) be an \(S\)-torus whose closed fibre is isomorphic to
\(T_s\) (Exposé X~4.6). Let \({}_\ell H\) denote the ``kernel'' of
the \(\ell\)-th power morphism in \(H\), and let \(U_\ell\) be the
largest open subset of \({}_\ell H\) étale over \(S\). \SGARefLink{sga3:xv:lemma:1:3}{Lemma 1.3} and
flatness of \(H\) show that \(U_\ell\) contains \({}_\ell T_s\).
Hensel's lemma gives a unique \(S\)-morphism
\[
  u_\ell:{}_\ell T\longrightarrow U_\ell
\]
inducing the canonical immersion
\({}_\ell T_s\to(U_\ell)_s\) on the closed fibre. Uniqueness shows
that \(u_\ell\) is a central homomorphism (Exposé \SGARefLink{sga3:ix:theorem:5:6:item:a}{IX~5.6(a)}).
Proceeding as in \SGARefLink{sga3:xv:lemma:6:6}{6.6} and \SGARefLink{sga3:xv:lemma:6:7}{6.7}, one proves that \(u_\ell\) is an
immersion and, if \(M_\ell=u_\ell({}_\ell T)\), that
\[
  \operatorname{Centr}_G(M_\ell)=H
\]
for all sufficiently large \(\ell\). This is where noetherianity of
\(S\) is used.

\medskip\noindent\textit{Proof of (i).}
If \(H\in\mathcal{CT}(S)\), then \(H\) has connected fibres (Exposé
\SGARefLink{sga3:xii:theorem:6-6:item:b}{XII~6.6(b)}) and equals its connected normalizer (\SGARefLink{sga3:xv:remark:6:5}{6.5}). Hence the
present functor \(\mathcal{CT}\) is the functor of \SGARefLink{sga3:xv:paragraph:6-4-0}{6.4.0}. \SGARefLink{sga3:xv:theorem:6:4}{Theorem 6.4}
therefore represents it by an \(S\)-prescheme of finite presentation,
quasi-projective over \(S\). It remains to prove smoothness.

EGA IV~8 reduces us to affine noetherian \(S\). By Exposé XI~1.5 it
then suffices to prove the following infinitesimal lifting property:
if \(S\) is local artinian, \(S_0\) is defined by a nilpotent ideal,
and \(H_0\in\mathcal{CT}(S_0)\), then \(H_0\) lifts to a smooth
subgroup prescheme \(H\) of \(G\). By part (ii),
\[
 H_0=\operatorname{Centr}_{G_0}(T_0)
\]
for a sub-torus \(T_0\) of \(G_0\). Since \(G\) is smooth, \(T_0\)
lifts to a sub-torus \(T\) of \(G\) (Exposé \SGARefLink{sga3:ix:theorem:3:6:bis}{IX~3.6 bis}), and
\(H=\operatorname{Centr}_G(T)\) works; it is smooth by Exposé \SGARefLink{sga3:xi:corollary:2:4}{XI~2.4}
and \SGARefLink{sga3:xv:lemma:2:5}{Lemma 2.5}.

\medskip\noindent\textit{Proof of (iii).}
Because \(H\) is smooth with connected fibres, Exposé \SGARefLink{sga3:xi:corollary:6:11}{XI~6.11}
represents \(\operatorname{Transp}_G(H,L)\) by a closed subprescheme
of \(G\), of finite presentation over \(S\). Smoothness reduces to
the following lifting problem. Let \(S\) be local artinian, \(S_0\)
a closed subprescheme, and \(g_0\in G(S_0)\) such that
\[
  \operatorname{int}(g_0)H_0\subseteq L_0.
\]
We must lift \(g_0\) to \(g\in G(S)\) with
\(\operatorname{int}(g)H\subseteq L\).

Smoothness of \(G\) gives \(g_1\in G(S)\) lifting \(g_0\). Put
\(H'=\operatorname{int}(g_1)H\). Then
\(H'\in\mathcal{CT}(S)\) and \(H'_0\subseteq L_0\). By (ii),
\(H'=\operatorname{Centr}_G(T')\) for a torus \(T'\subseteq G\).
Since \(L\) is smooth, \(T'_0\subseteq L_0\) lifts to a torus
\(T''\subseteq L\) (Exposé \SGARefLink{sga3:ix:theorem:3:6:bis}{IX~3.6 bis}). The group
\(\operatorname{Centr}_L(T'')\) lies in
\(\operatorname{Centr}_G(T'')\), has the same closed fibre as the
latter, namely \(H'_0\), and is smooth. Hence
\[
 \operatorname{Centr}_L(T'')
 =\operatorname{Centr}_G(T'').
\]
The sub-tori \(T'\) and \(T''\) of \(G\) lift the same \(T'_0\), so
they are conjugate by an element \(h\in G(S)\) reducing to the
identity (Exposé \SGARefLink{sga3:ix:corollary:3:3:bis}{IX~3.3 bis}). Their centralizers are therefore
conjugate. Then \(g=hg_1\) lifts \(g_0\) and satisfies
\(\operatorname{int}(g)H\subseteq L\).

For \(g\in\operatorname{Transp}_G(H,L)(S)\), the equality
\(\operatorname{int}(g)H=L\) holds if and only if
\(\dim H_s=\dim L_s\) for every \(s\in S\). If \(U\) is the
open-and-closed subprescheme over which these dimensions agree, then
\[
 \operatorname{Transpstr}_G(H,L)
 =U\times_S\operatorname{Transp}_G(H,L).
\]

\medskip\noindent\textit{Proof of (iv).}
To prove \(H\) closed in \(G\), reduce first to affine noetherian
\(S\), then, by EGA IV~8, to the spectrum of a complete local ring.
Part (ii) writes \(H\) as the centralizer in \(G\) of a subgroup of
multiplicative type, so \(H\) is closed because \(G\) is separated
(Exposé VI\textsubscript{B}~\SGARefLink{sga3:xv:theorem:5-2}{5.2}).

By (iii),
\[
 N=\operatorname{Norm}_G(H)
  =\operatorname{Transpstr}_G(H,H)
\]
is represented by a subgroup prescheme smooth over \(S\) and closed
in \(G\). The morphism
\[
  G\longrightarrow\mathcal{CT},
  \qquad g\longmapsto\operatorname{int}(g)H
\]
is smooth by (iii), so its image is an open subset \(U\) of
\(\mathcal{CT}\). As in Exposé \SGARefLink{sga3:xi:corollary:5:3}{XI~5.3}, \(G/N\) is represented by
\(U\); it is therefore quasi-projective.

We now treat \(N/H\). By EGA IV~8, assume \(S\) affine noetherian and
then \(S=\operatorname{Spec}A\), with \(A\) local of finite
dimension\footnote{Krull dimension.}. We use induction on
\(\dim S\). When \(\dim S=0\), the assertion is Exposé
VI\textsubscript{A}, §4. If \(N/H\) is representable, it is separated
because \(H\) is closed in \(N\), of finite type and étale because
\(N\) is smooth of finite type and \(H\) is open. It is therefore
quasi-affine (SGA 1, VIII~6.2). Effective descent for quasi-affine
schemes (\emph{loc. cit.}, 7.9) permits replacement of \(A\) by its
completion.

Thus assume \(S\) complete noetherian local, with closed point \(s\),
and put \(U=S\setminus\{s\}\). By induction,
\((N|_U)/(H|_U)\) is represented by a \(U\)-group \(K\). Let \(T_s\)
be the maximal central torus of \(H_s\), choose \(q\) invertible on
\(S\), let \(\ell\) be a power of \(q\), and let \(M_\ell\) be the
unique central subgroup of multiplicative type of \(H\) lifting
\({}_\ell T_s\), as in (ii). Choose \(\ell\) so large that
\[
  \operatorname{Centr}_G(M_\ell)=H,
\]
and put \(N'=\operatorname{Norm}_G(M_\ell)\).

The torus \(T_s\), and hence \((M_\ell)_s\), is characteristic in
\(H_s\). Therefore \(N_s\) normalizes \((M_\ell)_s\), and
\(N_s=N'_s\). The proof of Exposé \SGARefLink{sga3:xi:proposition:5:9}{XI~5.9} represents
\[
  N'/H
  =\operatorname{Norm}_G(M_\ell)/
   \operatorname{Centr}_G(M_\ell)
\]
by an \(S\)-group \(K'\). Since \(N'\) is smooth and
\(N'_s=N_s\), it is an open subgroup of \(N\) containing \(H\)
(Exposé VI\textsubscript{B}~\SGARefLink{sga3:xv:lemma:2:5}{2.5}). Its image over \(U\) is an open
subgroup of \(K\), isomorphic to \(K'|_U\). Glue \(K\) and \(K'\)
along this open subgroup to obtain \(L\), and glue the projections
\(N|_U\to K\) and \(N'\to K'\) to obtain \(p:N\to L\). Then
\((L,p)\) represents \(N/H\).

\medskip\noindent\textit{Proof of (v).}
First suppose \(S=\operatorname{Spec}k\). The image \(H'\) of \(H\)
is a smooth subgroup of \(G'\). The required assertion follows from:

\SGAStatement{Lemma 7.2}{Let \(u:G\to G'\) be an epimorphism of
smooth connected algebraic groups over a field \(k\), let \(T\) be a
torus of \(G\), and let \(T'\) be its image in \(G'\). Then
\[
 u(\operatorname{Centr}_G(T))
 =\operatorname{Centr}_{G'}(T').
\]}
\SGARefTarget{sga3:xv:lemma:7:2}

Put \(H=\operatorname{Centr}_G(T)\),
\(H'=\operatorname{Centr}_{G'}(T')\), and \(H''=u(H)\). We have
\(H''\subseteq H'\). After extending the base field, assume \(k\)
algebraically closed. It suffices to prove that every Cartan subgroup
of \(H'\) lies in \(H''\): then \(H''\) contains the open subset of
regular points of \(H'\), so is open, and equals the connected group
\(H'\).

Let \(C'\) be a Cartan subgroup of \(H'\). It is also a Cartan
subgroup of \(G'\), since \(H'\), being the centralizer of the torus
\(T'\), has the same reductive and nilpotent ranks as \(G'\). Put
\[
  K=(u^{-1}(C'))^0_{\mathrm{red}}.
\]
Because \(T'\) is central in \(H'\), \(C'\) contains \(T'\), and
\(K\) contains \(T\). Let \(C\) be a Cartan subgroup of \(K\)
containing \(T\). The torus \(T\) lies in the unique maximal torus of
\(C\), central in \(C\) (Exposé \SGARefLink{sga3:xii:6.6:item:c}{XII~6.6(c)}); hence
\(C\subseteq H\). Any two Cartan subgroups of \(G\) are conjugate,
and the image of a Cartan subgroup of \(G\) is a Cartan subgroup of
\(G'\) (Exposé \SGARefLink{sga3:xii:6.6}{XII~6.6}). It follows that \(C\) is a Cartan subgroup
of \(G\), and its image is a Cartan subgroup of \(G'\) contained in
\(C'\). Thus \(u(C)=C'\), proving \(C'\subseteq H''\).

This proves (v) over a field. In general, finite presentation and the
usual reductions permit us to assume \(S\) affine noetherian, then
local, and finally the spectrum of a complete noetherian local ring
\(A\), by fpqc descent of subpreschemes of \(G'\).

Use the notation of (ii). Let \(T_s\) be the maximal central torus of
\(H_s\), and let \(M_\ell\subset H\) lift \({}_\ell T_s\), with
\(H=\operatorname{Centr}_G(M_\ell)\). Let \(T'_s=u_s(T_s)\).
Since \(G'\) is separated, \(u(M_\ell)\) is represented by a subgroup
\(M'_\ell\) of multiplicative type (Exposé \SGARefLink{sga3:ix:theorem:6:8}{IX~6.8}). Put
\[
  H'=\operatorname{Centr}_{G'}(M'_\ell),
\]
which is smooth. For every power \(\ell'\) of \(q\), a sufficiently
large \(\ell\) makes \((M'_\ell)_s\) contain
\({}_{\ell'}T'_s\). We may therefore choose \(\ell\) such that
\[
 H'_s=\operatorname{Centr}_{G'_s}(T'_s)=u_s(H_s),
\]
the last equality being \SGARefLink{sga3:xv:lemma:7:2}{Lemma 7.2}.

The restriction \(v:H\to H'\) of \(u\) is flat. Indeed, \(H,H'\) are
flat over \(S\), \(H_s\to H'_s\) is flat, and hence \(v\) is flat
near \(H_s\) (EGA IV~11.3.10 and 11.3.1). It is therefore flat on an
open subgroup of \(H\) (Exposé VI\textsubscript{B}~2.2), and thus
everywhere because the fibres of \(H\) are connected. The
set-theoretic image of \(H\) is consequently an open subset of \(H'\),
necessarily \((H')^0\), and with its induced structure represents
the fpqc image sheaf \(u(H)\). \SGARefLink{sga3:xv:lemma:7:2}{Lemma 7.2} shows that it belongs to
\(\mathcal{CT}_{G'}(S)\).

If \(H=\operatorname{Centr}_G(T)\), then \(T'=u(T)\) is a sub-torus
of \(G'\) (Exposé \SGARefLink{sga3:ix:theorem:6:8}{IX~6.8}). The image \(u(H)\) is contained in
\(\operatorname{Centr}_{G'}(T')\), agrees with it fibrewise by
\SGARefLink{sga3:xv:lemma:7:2}{Lemma 7.2}, and is smooth. Hence
\[
  u(H)=\operatorname{Centr}_{G'}(T').
\]

\medskip\noindent\textit{Proof of (vi).}
Since \(\mathcal{CT}_G,\mathcal{CT}_{G'}\) are smooth over \(S\), it
is enough to prove faithful flatness on geometric fibres. Thus take
\(S=\operatorname{Spec}k\), with \(k\) algebraically closed. Let
\(H'\in\mathcal{CT}_{G'}(k)\), let \(T'\) be its maximal central
torus, choose a sub-torus \(T\subset G\) mapping onto \(T'\), and put
\[
 H=\operatorname{Centr}_G(T),\qquad
 N=\operatorname{Norm}_G(H),\qquad
 N'=\operatorname{Norm}_{G'}(H').
\]
By \SGARefLink{sga3:xv:lemma:7:2}{Lemma 7.2}, \(H'=u(H)\). Part (iv) identifies \(G/N\) and \(G'/N'\)
with open neighbourhoods of \(H\) and \(H'\) in the two
\(\mathcal{CT}\)-spaces. On these neighbourhoods, \(\widetilde u\)
is
\[
  w:G/N\longrightarrow G'/N',
\]
obtained from \(u\) by passage to quotients. This is an epimorphism
of homogeneous spaces under \(G\), hence faithfully flat. Therefore
\(\widetilde u\) is flat and its image is an open subset of
\(\mathcal{CT}_{G'}\) containing all \(k\)-points. Since
\(\mathcal{CT}_{G'}\) is of finite type, \(\widetilde u\) is
surjective and hence faithfully flat. If \(\ker u\) is central, the
remaining assertions are Exposé \SGARefLink{sga3:xii:7.12}{XII~7.12}.

The next theorem generalizes Exposé \SGARefLink{sga3:xii:7.1}{XII~7.1}.

\SGAStatement{\SGAInlineRefTarget{sga3:xv:theorem:7-3}Theorem 7.3}{Let \(S\) be a prescheme and let \(G\)
be an \(S\)-group prescheme of finite presentation, smooth over \(S\),
with connected fibres. Define the functor
\[
 \mathcal C(S')=\{\text{Cartan subgroups of }G_{S'}\}.
\]
\begin{enumerate}
\renewcommand{\labelenumi}{\textup{\roman{enumi})}}
\item The following are equivalent:
\begin{enumerate}
\renewcommand{\labelenumii}{\textup{\alph{enumii})}}
\item \(\mathcal C\) is representable;
\item \(\mathcal C\) is represented by an \(S\)-prescheme smooth,
quasi-projective, and of finite presentation over \(S\), with affine
fibres;
\item \(G\) admits a Cartan subgroup locally for the étale topology;
\item \(G\) admits a Cartan subgroup locally for the faithfully flat
topology;
\item the nilpotent rank of the fibres of \(G\) is locally constant
on \(S\).
\end{enumerate}
\item If these conditions hold, any two Cartan subgroups of \(G\)
are locally conjugate for the étale topology. The regular locus of
the fibres of \(G\) (Exposé \SGARefLink{sga3:xiii:corollary:2:7}{XIII~2.7}) is an open subprescheme
\(G_{\mathrm{reg}}\), of finite presentation over \(S\), and every
section of \(G_{\mathrm{reg}}\) over \(S\) is contained in a unique
Cartan subgroup of \(G\).
\item Let \(G'\) be an \(S\)-group prescheme of finite presentation
and \(u:G\to G'\) faithfully flat. Then, for every Cartan subgroup
\(C\) of \(G\), its image \(C'=u(C)\) is represented by a Cartan
subgroup of \(G'\), and \(C\to C'\) is faithfully flat.
\item Under the conditions of (i) and (iii), the morphism
\[
 \widetilde u:\mathcal C_G\longrightarrow\mathcal C_{G'},
 \qquad C\longmapsto u(C),
\]
is faithfully flat. If \(\ker u\) is central, it is an isomorphism.
\item Further information on transporters and relations with maximal
tori is given in \SGARefLink{sga3:xv:proposition:7-1}{Proposition 7.1} and Exposé \SGARefLink{sga3:xii:7.1}{XII~7.1}.
\end{enumerate}}
\SGARefTarget{sga3:xv:theorem:7:3}

\SGAProof{Proof}
For (i) we prove
\[
 (\mathrm b)\Rightarrow(\mathrm c)\Rightarrow(\mathrm d)
 \Rightarrow(\mathrm e)\Rightarrow(\mathrm b)
 \Rightarrow(\mathrm a)\Rightarrow(\mathrm d).
\]
Assume (b), and let \(s\in S\). Since \(\mathcal C_s\) is smooth
over \(\kappa(s)\), it has points with residue field a finite
separable extension of \(\kappa(s)\). Exposé \SGARefLink{sga3:xi:proposition:1:10}{XI~1.10} gives an open
neighbourhood \(U\) of \(s\) and an étale surjective morphism
\(U'\to U\) such that \(G_{U'}\) admits a Cartan subgroup. Thus
(b)\(\Rightarrow\)(c), while (c)\(\Rightarrow\)(d) is clear.

Assume (d), and let \(s\in S\). There is a flat \(S\)-prescheme
\(S'\), whose image contains \(s\), such that \(G_{S'}\) admits a
Cartan subgroup. For \(s'\in S'\) above \(s\), the nilpotent rank of
the fibres of \(G_{S'}\) is constant on
\(\operatorname{Spec}\mathcal O_{S',s'}\). Hence the nilpotent rank
of the fibres of \(G\) is constant on
\(\operatorname{Spec}\mathcal O_{S,s}\), which is the image of that
local spectrum (EGA IV~2.3.4(ii)). Let the value be \(r\). By 6.3 bis,
the set \(E_r\) where the nilpotent rank equals \(r\) is
ind-constructible; it contains a neighbourhood of \(s\) by
EGA IV~1.10.1. Thus (d)\(\Rightarrow\)(e).

Assume (e). Locally on \(S\), the nilpotent rank is a constant \(r\).
For every \(S'\), the Cartan subgroups of \(G_{S'}\) are exactly the
smooth subgroup preschemes of relative dimension \(r\) whose
geometric fibres are centralizers of tori. By \SGARefLink{sga3:xv:proposition:7-1}{Proposition 7.1},
\(\mathcal C\) is the open-and-closed subprescheme of
\(\mathcal{CT}\) representing the relative-dimension-\(r\)
subfunctor. It has all the properties in (b); affineness of its
fibres is Exposé \SGARefLink{sga3:xii:theorem:7-1:item:d}{XII~7.1(d)}. This proves
(e)\(\Rightarrow\)(b), and (b)\(\Rightarrow\)(a) is immediate.

Finally assume (a). By \SGARefLink{sga3:xv:proposition:6-2:item:iii}{6.2(iii)}, \(\mathcal C\) commutes with filtered
inductive limits of rings, so its representing prescheme is locally
of finite presentation (EGA IV~8.14.2). To prove it smooth, it is
enough to show that, when \(S\) is affine and \(S_0\subseteq S\) is
defined by a nilpotent ideal, every Cartan subgroup \(C_0\) of
\(G_0\) lifts to a Cartan subgroup of \(G\). Existence of \(C_0\)
implies (e) on \(S_0\), hence on \(S\), which has the same underlying
space. Since (e)\(\Rightarrow\)(b), \(\mathcal C\) is smooth.
It is also surjective over \(S\), so (a)\(\Rightarrow\)(d).

For (ii), if \(C,C'\) are Cartan subgroups, then
\(\operatorname{Transp}_G(C,C')\) is represented by a smooth
\(S\)-prescheme (\SGARefLink{sga3:xv:proposition:7-1:item:iii}{Proposition 7.1(iii)}) with nonempty fibres
(Exposé \SGARefLink{sga3:xii:theorem:6-6:item:a}{XII~6.6(a)},(c)). Hensel's lemma (Exposé \SGARefLink{sga3:xi:proposition:1:10}{XI~1.10}) gives local
conjugacy for the étale topology. The assertions about regular points
follow from Exposé \SGARefLink{sga3:xiii:theorem:3:1}{XIII~3.1} and \SGARefLink{sga3:xiii:corollary:3:2}{3.2}.

For (iii), \SGARefLink{sga3:xv:proposition:7-1:item:v}{Proposition 7.1(v)} represents \(u(C)\) by a smooth
subgroup \(C'\) of \(G'\). Its fibres are Cartan subgroups of the
fibres of \(G'\) (Exposé \SGARefLink{sga3:xii:6.6:item:d}{XII~6.6(d)}), so \(C'\) is a Cartan subgroup.

For (iv), faithful flatness is proved as in \SGARefLink{sga3:xv:proposition:7-1:item:vi}{Proposition 7.1(vi)}.
If \(\ker u\) is central and \(C'\) is a Cartan subgroup of \(G'\),
then \(C=u^{-1}(C')\) is smooth with connected fibres
(\SGARefLink{sga3:xv:proposition:7-1:item:vi}{Proposition 7.1(vi)}), and its fibres are Cartan subgroups (Exposé
\SGARefLink{sga3:xii:6.6:item:e}{XII~6.6(f)}). Thus \(C\) is a Cartan subgroup, and \SGARefLink{sga3:xv:proposition:7-1:item:vi}{Proposition
7.1(vi)} completes the proof.

% END inputs/XV/tex/components/07_expose_XV_section7.tex

% BEGIN inputs/XV/tex/components/08_expose_XV_section8_through_theorem8_9.tex
% Authority: Expo15.pdf, local pp. 56--61 / combined pp. 926--931 /
% printed pp. 459--470.
% Stops immediately before the proof of Theorem 8.8(a).

\SGASection{8}{Representability criterion for the functor of sub-tori
of a smooth group}
\SGARefTarget{sga3:xv:section:8}

\medskip\noindent\SGAInlineRefTarget{sga3:xv:paragraph:8-0}\textbf{8.0.}
In this section, if \(S\) is a prescheme and \(G\) an \(S\)-group
prescheme, \(\mathcal T_G\) (or simply \(\mathcal T\), if there is no
ambiguity) denotes the \(S\)-functor
\[
  (\operatorname{Sch}/S)^\circ\longrightarrow\operatorname{Ens}
\]
such that, for every \(S\)-prescheme \(S'\),
\[
  \mathcal T(S')=\{\text{sub-tori of \(G_{S'}\)}\}.
\]
The functor \(\mathcal T_C\) of central sub-tori of \(G\) is defined
similarly.
\SGARefTarget{sga3:xv:paragraph:8:0}

\SGAStatement{\SGAInlineRefTarget{sga3:xv:proposition:8-1}Proposition 8.1}{Let \(S\) be a locally noetherian
prescheme and let \(G\) be an \(S\)-group prescheme of finite type.
Then the following conditions are equivalent:
\begin{enumerate}
\item[\textup{i)}] The functor \(\mathcal T_G\) ``commutes with adic
limits of local artinian rings'': for every \(S\)-prescheme
\(S'=\operatorname{Spec}(A)\), where \(A\) is a complete noetherian
local ring for the topology defined by its maximal ideal
\(\mathfrak m\), the canonical map
\[
  \mathcal T(S')\longrightarrow
  \varprojlim_n\mathcal T(S'_n),
  \qquad S'_n=\operatorname{Spec}(A/\mathfrak m^n),
\]
is bijective.
\item[\textup{ii)}] The assertion in \textup{i)}, restricted to the
case in which \(A\) is a complete discrete valuation ring with
algebraically closed residue field.
\item[\textup{iii)}] The assertion in \textup{ii)}, restricted to the
subfunctor \(\mathcal T^{(1)}\) of \(\mathcal T\) formed by sub-tori
of \(G\) of relative dimension \(1\).
\item[\textup{i bis)}] For every \(S\)-prescheme \(S'\) as in
\textup{i)} and every \(S'\)-torus \(T\), the canonical map
\[
  \operatorname{Hom}_{S'\text{-gr}}(T,G_{S'})
  \longrightarrow
  \varprojlim_n
  \operatorname{Hom}_{S'_n\text{-gr}}(T_{S'_n},G_{S'_n})
\]
is bijective.
\item[\textup{ii bis)}] The assertion in \textup{i bis)}, restricted
to the case in which \(A\) is a complete discrete valuation ring
with algebraically closed residue field.
\item[\textup{iii bis)}] The assertion in \textup{ii bis)}, restricted
to \(T=\mathbb G_{m,S'}\).
\end{enumerate}}
\SGARefTarget{sga3:xv:proposition:8:1}

\medskip\noindent\textbf{Remark 8.2.}
There is an analogous proposition in which one restricts to central
sub-tori of \(G\) and to central homomorphisms from a torus into \(G\).
\SGARefTarget{sga3:xv:remark:8:2}

\medskip\noindent\textit{Proof.}
We shall use the following lemma.

\SGAStatement{Lemma 8.3}{Let \(S\) be a prescheme, \(G\) an
\(S\)-group prescheme, \(T\) an \(S\)-torus, and
\(u:T\to G\) an \(S\)-group morphism. Suppose, moreover, either that
\(G\) is of finite presentation over \(S\), or that \(S\) is locally
noetherian and \(G\) is locally of finite type. Then:
\begin{enumerate}
\renewcommand{\labelenumi}{\textup{\alph{enumi})}}
\item \(\ker u\) is a subgroup of multiplicative type of \(T\);
\item the quotient \(T'=T/\ker u\) is a torus;
\item the canonical monomorphism \(T'\to G\) induced by \(u\) after
passage to the quotient is an immersion.
\end{enumerate}}
\SGARefTarget{sga3:xv:lemma:8:3}

When \(G\) is separated over \(S\), this lemma follows from
Exposé \SGARefLink{sga3:ix:theorem:6:8}{IX~6.8}. In the general case, the usual reduction allows us to
suppose that \(S\) is noetherian. We first show that
\(K=\ker u\) is flat. We may suppose that \(S\) is the spectrum of a
local artinian ring (EGA \(0_{\mathrm{III}}\), 10.2.6); then \(G\) is
separated (Exposé VI\textsubscript{B}~5.2), so \(K\) is of
multiplicative type (Exposé \SGARefLink{sga3:ix:theorem:6:8}{IX~6.8}), and hence is flat over \(S\).
We next show that \(\ker u\) is closed in \(T\). This reduces us to
the case in which \(S\) is the spectrum of a discrete valuation ring
(EGA II~7.2.1). Since \(T\) is flat with connected fibres, \(u\)
factors through the identity component of the schematic closure in
\(G\) of the generic fibre of \(G\). We may thus suppose that \(G\)
is flat with connected fibres; it is then separated
(Exposé VI\textsubscript{B}~\SGARefLink{sga3:xv:theorem:5-2}{5.2}), and \(K\) is closed.
It follows from Exposé \SGARefLink{sga3:x:corollary:4-8:item:b}{X~4.8(b)} that \(K\) is a subgroup of
multiplicative type of \(T\). The quotient \(T'=T/K\) is therefore
representable and is a group of multiplicative type
(Exposé \SGARefLink{sga3:ix:proposition:2:7}{IX~2.7}) whose fibres are tori; hence \(T'\) is a torus.
Finally, the monomorphism \(T'\to G\) is an immersion by
Exposé \SGARefLink{sga3:viii:proposition:7:9}{VIII~7.9}.

\medskip\noindent\textit{Proof of \SGARefLink{sga3:xv:proposition:8-1}{Proposition 8.1}.}
\textup{i)}\(\Rightarrow\)\textup{i bis)}.
Put \(T_n=T_{S'_n}\), \(G_n=G_{S'_n}\), and let
\[
 (u_n)_{n\in\mathbb N}\in
 \varprojlim_n\operatorname{Hom}_{S'_n\text{-gr}}(T_n,G_n).
\]
For every \(n\), \(u_n(T_n)\) is a sub-torus \(T'_n\) of \(G_n\)
(\SGARefLink{sga3:xv:lemma:8:3}{Lemma 8.3}). By hypothesis, there is a unique sub-torus \(T'\) of
\(G\) lifting \(T'_n\) for every \(n\). Since \(T'\) has affine
fibres, the assertion follows from \SGARefLink{sga3:xv:corollary:4:4}{4.4}.

\textup{i bis)}\(\Rightarrow\)\textup{i)}.
Let
\((T_n)_{n\in\mathbb N}\in\varprojlim_n\mathcal T(S'_n)\).
By Exposé \SGARefLink{sga3:x:corollary:4-6}{X~4.6} there are an \(S'\)-torus \(T^0\) and an
\(S'_0\)-isomorphism
\[
  u_0:T^0_0\xrightarrow{\sim}T_0.
\]
Since \(T_n\) is smooth over \(S'_n\), \(u_0\) lifts to an
\(S'_n\)-morphism \(u_n:T^0_n\to T_n\) (Exposé \SGARefLink{sga3:ix:theorem:3:6}{IX~3.6}), and the
family \((u_n)_{n\in\mathbb N}\) may be chosen coherent. It therefore
comes from a morphism \(u:T^0\to G\). By \SGARefLink{sga3:xv:lemma:8:3}{Lemma 8.3}, the image of
\(T^0\) under \(u\) is a sub-torus of \(G\), and it lifts \(T_n\)
for every \(n\).

The implications
\[
 \textup{i)}\Rightarrow\textup{ii)}\Rightarrow\textup{iii)}
 \quad\text{and}\quad
 \textup{i bis)}\Rightarrow\textup{ii bis)}
 \Rightarrow\textup{iii bis)}
\]
are immediate. The implication
\(\textup{iii)}\Rightarrow\textup{iii bis)}\) is proved in the same
way as
\(\textup{i)}\Rightarrow\textup{i bis)}\).
It remains to prove
\[
 \textup{iii bis)}\Rightarrow\textup{ii bis)}
 \Rightarrow\textup{i bis)}.
\]

\textup{ii bis)}\(\Rightarrow\)\textup{i bis)}.
In the terminology of \SGARefLink{sga3:xv:proposition:4:3}{4.3}, assertion \textup{i bis)} holds if and
only if every element
\[
 (u_n)_{n\in\mathbb N}\in
 \varprojlim_n\operatorname{Hom}_{S'_n\text{-gr}}(T_n,G_n)
\]
is admissible. For every point \(t\) of \(S'\) distinct from its
closed point \(s\), there is an \(S\)-scheme \(S''\), the spectrum of
a complete discrete valuation ring with algebraically closed residue
field, whose generic point maps to \(t\) and whose closed point maps
to \(s\) (EGA II~7.1.9). This gives at once a valuative criterion for
a family of morphisms to be admissible, and proves the implication.

\textup{iii bis)}\(\Rightarrow\)\textup{ii bis)}.
Let \(T\) be an \(S\)-torus, where \(S\) is the spectrum of a complete
discrete valuation ring with algebraically closed residue field. The
torus \(T\) is split (Exposé \SGARefLink{sga3:x:corollary:4-6}{X~4.6}), say
\[
  T\simeq\mathbb G_{m,S}^{\,r}.
\]
Let
\[
 (u_n)_{n\in\mathbb N}\in
 \varprojlim_n
 \operatorname{Hom}_{S_n\text{-gr}}
 \bigl((\mathbb G_m)^r_{S_n},G_n\bigr).
\]
By hypothesis, the restriction of \(u_n\) to each factor of
\(\mathbb G_{m,S}^{\,r}\) comes from a group morphism
\(\mathbb G_{m,S}\to G\). Their product, followed by the composition
law \(G^r\to G\), gives a morphism
\[
  v:\mathbb G_{m,S}^{\,r}\longrightarrow G.
\]
Clearly \(u_n=v_n\) for every \(n\). It remains to prove that \(v\)
is a group morphism. Equivalently, two evident morphisms
\[
 f,g:X=\mathbb G_{m,S}^{\,r}\times_S\mathbb G_{m,S}^{\,r}
 \longrightarrow G
\]
must coincide. Let \(Z\) be their coincidence subscheme. Since
\(\mathbb G_{m,S}^{\,r}\) has connected fibres and its image under
\(v\) contains the unit section, the argument of \SGARefLink{sga3:xv:lemma:8:3}{Lemma 8.3} shows
that \(v\) factors through the identity component of \(G\). We may
therefore suppose that \(G\) is separated
(Exposé VI\textsubscript{B}~\SGARefLink{sga3:xv:theorem:5-2}{5.2}), so \(Z\) is closed. On the other
hand, \(v_n\) is a group morphism, hence \(f_n=g_n\) for every \(n\).
Thus \(Z\) contains a neighbourhood of the closed fibre of \(X\)
(EGA I~10.9.4), which is schematically dense in \(X\), since \(X\)
has smooth irreducible fibres (Exposé IX~4.6). Consequently
\(Z=X\), and \(f=g\).

\SGAStatement{\SGAInlineRefTarget{sga3:xv:definition:8-4}Definition 8.4}{Let \(S\) be a locally noetherian
prescheme and \(G\) an \(S\)-group prescheme of finite type. Let
\(\mathcal S\) be the set of \(S\)-schemes \(S'\) that are spectra
of complete discrete valuation rings with algebraically closed
residue field. Write \(s\) and \(t\) for the closed and generic
points of \(S'\), and require that
\((G_t)^0_{\mathrm{red}}\) be smooth and admit a Chevalley
decomposition: an extension of an abelian variety by a smooth
connected linear algebraic subgroup \(L_t\). This decomposition is
then unique. For \(S'\in\mathcal S\), let \(G'\), respectively \(L\),
denote the schematic closure in \(G_{S'}\) of \(G_t\), respectively
of \(L_t\).

We say that ``the abelian part of \(G\) does not degenerate into a
toric part'', or briefly that \(G\) has property \(\mathrm{AT}\), if,
for every \(S'\in\mathcal S\), \(L_s\) has the same reductive rank
as \(G'_s\). Intuitively, suppose that the quotient \(A=G'/L\) is
representable. Then \(A\) is a flat group prescheme and
\((A_t)^0_{\mathrm{red}}\) is an abelian variety; property
\(\mathrm{AT}\) says that \(A_s\) has reductive rank zero, so that
\((A_s)^0_{\mathrm{red}}\) is an extension of an abelian variety by
a unipotent group.

Similarly, if \(G\) has connected fibres, we say that \(G\) has
property \(\mathrm{ATC}\) if, for every \(S'\in\mathcal S\), the
schematic closure \(Z\) of the centre \(Z_t\) of \(G_t\) has
property \(\mathrm{AT}\).}
\SGARefTarget{sga3:xv:definition:8:4}

These technical definitions are justified by the following
proposition.

\SGAStatement{Proposition 8.5}{Let \(S\) be a locally noetherian
prescheme and let \(G\) be an \(S\)-group prescheme of finite type.
\begin{enumerate}
\renewcommand{\labelenumi}{\textup{\alph{enumi})}}
\item\SGAInlineRefTarget{sga3:xv:proposition:8-5:item:a} The functor \(\mathcal T\) of sub-tori of \(G\) commutes with
adic limits of local artinian rings (cf. \SGARefLink{sga3:xv:proposition:8-1}{Proposition 8.1}) if and
only if \(G\) has property \(\mathrm{AT}\).
\item\SGAInlineRefTarget{sga3:xv:proposition:8-5:item:b} If \(G\) has connected fibres, the functor \(\mathcal T_C\) of
central sub-tori of \(G\) commutes with adic limits of local artinian
rings if and only if \(G\) has property \(\mathrm{ATC}\).
\end{enumerate}}
\SGARefTarget{sga3:xv:proposition:8:5}

We need the following technical lemma.

\SGAStatement{Lemma 8.6}{Let \(S\) be the spectrum of a discrete
valuation ring, with closed point \(s\), let \(G\) be an
\(S\)-group prescheme flat and of finite type, and let \(T_s\) be a
sub-torus of \(G_s\). There are an \(S\)-scheme \(S'\), the spectrum
of a discrete valuation ring faithfully flat over \(S\), with closed
point \(s'\), and a subgroup scheme \(C\) of \(G_{S'}\), flat over
\(S'\), commutative, and with connected fibres, such that
\(C_{s'}\) contains \(T_s\times_s s'\).}
\SGARefTarget{sga3:xv:lemma:8:6}

After replacing \(S\) by the spectrum \(S'\) of a discrete valuation
ring faithfully flat over it, we may suppose that
\(T_s=\mathbb G_{m,s}^{\,r}\) and that the transcendence degree of
\(\kappa(s)\) over its prime field is at least \(r\)
(EGA \(0_{\mathrm{III}}\), 10.3.1, and EGA II~7.1.9). There is then
an element \(x\in T_s(s)\) such that every algebraic subgroup of
\(G_s\) containing \(x\) contains \(T_s\)
(cf. Exposé XIII, proof of \SGARefLink{sga3:xiii:theorem:2:1}{2.1},
\(\textup{ii)}\Rightarrow\textup{vii)}\)).
Since \(G\) is flat over \(S\), a quasi-section passes through \(x\)
(EGA IV~14.5.8). After another such replacement of \(S\), we may
therefore suppose that \(x\) lifts to a section
\(\widetilde x\) of \(G\) over \(S\). Let \(C_t\) be the commutative
algebraic subgroup of \(G_t\) generated by \(\widetilde x_t\)
(Exposé VI\textsubscript{B}~7), and let \(C\) be its schematic
closure in \(G\). Clearly \(C_s\) contains \(x\), hence contains
\(T_s\). The identity component of \(C\) is therefore a flat
commutative group scheme with connected fibres having the required
property.

\medskip\noindent\textit{Proof of \SGARefLink{sga3:xv:proposition:8-5:item:a}{Proposition 8.5(a)}.}
Suppose first that \(\mathcal T\) commutes with adic limits of
artinian rings. Let \(S'\in\mathcal S\), and let \(T_s\) be a maximal
torus of \(G'_s\). We must show that \(T_s\subset L_s\). The
formation of \(L\) and \(G'\) plainly commutes with faithfully flat
extensions \(S''\to S'\) of discrete valuation rings. By \SGARefLink{sga3:xv:lemma:8:6}{Lemma 8.6},
we may therefore replace \(S'\) so that there is a flat commutative
subgroup scheme \(C\) of \(G'_{S'}\) for which \(C_s\) contains
\(T_s\). The torus \(T_s\) is central in \(C_s\), and hence lifts
infinitesimally to a central sub-torus (\SGARefLink{sga3:xv:corollary:2:3}{Corollary 2.3}). By the
hypothesis on \(G\), \(T_s\) lifts to a sub-torus \(T\) of \(G\).
Clearly \(T_t\) lies in the linear part \(L_t\) of \(G_t\), and
therefore \(T\subset L\).

Conversely, suppose that \(G\) has property \(\mathrm{AT}\). We prove
condition \SGARefLink{sga3:xv:proposition:8-1}{8.1} \textup{iii bis)}. Let \(S=\operatorname{Spec}(A)\),
where \(A\) is a complete discrete valuation ring with algebraically
closed residue field, let \(\mathfrak m\) be its maximal ideal, put
\(S_n=\operatorname{Spec}(A/\mathfrak m^n)\), and let
\[
 u_n:(\mathbb G_{m,S})_n\longrightarrow
 G_n=G\times_S S_n,\qquad n\in\mathbb N,
\]
be a coherent system of group morphisms. Let \(q\) be a prime
invertible on \(S\). As \(\ell\) ranges over the powers of \(q\),
there is a unique \(S\)-group morphism
\[
 {}_\ell u:{}_\ell\mathbb G_{m,S}\longrightarrow G
\]
lifting the restriction of \(u_n\) to
\({}_\ell(\mathbb G_{m,S})_n\) for every \(n\)
(\SGARefLink{sga3:xv:proposition:1-6:item:a}{Proposition 1.6(a)}). Consequently, if an \(S\)-group morphism
\(u:\mathbb G_{m,S}\to G\) lifting every \(u_n\) exists, its
restriction to \({}_\ell\mathbb G_{m,S}\) is uniquely determined.
The density theorem (Exposé \SGARefLink{sga3:ix:corollary:4:8:item:a}{IX~4.8(a)}) proves both the uniqueness of
\(u\) and that, for its existence, we may make a faithfully flat
extension of the base.

Now \((G_t)^0_{\mathrm{red}}\) admits a Chevalley decomposition,
already defined over a finite extension \(L\) of the fraction field
\(K\) of \(A\). After replacing \(S\) by its normalization in \(L\),
we may suppose that \(S\in\mathcal S\). The morphism
\({}_\ell u_t\) factors through \((G_t)_{\mathrm{red}}\), and
\({}_\ell u\) therefore factors through the schematic closure
\(G''\) of \((G_t)_{\mathrm{red}}\) in \(G'\). The density theorem
then shows that every \(u_n\) factors through \(G''_n\).
Since \(G\) has property \(\mathrm{AT}\), every sub-torus of \(G'_s\),
and hence of \(G''_s\), is contained in \(L_s\); thus \(u_s\)
factors through \(L_s\). In fact every \(u_n\) factors through
\(L_n\). Indeed, \(L_t\) is invariant in \(G''_t\), so \(L\) is
invariant in \(G''\); moreover \(L\) is flat over \(S\), and the
quotient
\[
  H_n=G''_n/L_n
\]
is representable for every \(n\) (Exposé VI\textsubscript{A},
\S4). The image of \((\mathbb G_{m,S})_n\) in \(H_n\) is a
sub-torus (Exposé \SGARefLink{sga3:ix:theorem:6:8}{IX~6.8}) whose closed fibre is zero. It is therefore
zero, and \(u_n\) factors through \(L_n\). Since \(L_t\) is affine,
\SGARefLink{sga3:xv:proposition:4:3}{Proposition 4.3} now shows that the family \((u_n)\) is admissible.

The proof of \SGARefLink{sga3:xv:proposition:8-5:item:b}{Proposition 8.5(b)} is entirely analogous, in view of
\SGARefLink{sga3:xv:remark:8:2}{Remark 8.2} and Exposé \SGARefLink{sga3:ix:theorem:5:6:item:a}{IX~5.6(a)}.

\SGAStatement{Proposition 8.7}{Let \(S\) be a locally noetherian
prescheme and let \(G\) be an \(S\)-group prescheme of finite type.
\begin{enumerate}
\renewcommand{\labelenumi}{\textup{\roman{enumi})}}
\item If \(G\) has connected fibres and has property
\(\mathrm{AT}\), then it has property \(\mathrm{ATC}\).
\item If \(G\) has property \(\mathrm{AT}\), every subgroup
prescheme of \(G\) has property \(\mathrm{AT}\).
\item If \(G\) is flat over \(S\) and the abelian rank of its fibres
is a locally constant function on \(S\), then \(G\) has property
\(\mathrm{AT}\).
\end{enumerate}}
\SGARefTarget{sga3:xv:proposition:8:7}

Assertion \textup{(i)} follows at once from \SGARefLink{sga3:xv:proposition:8:5}{Proposition 8.5} and
Exposé \SGARefLink{sga3:ix:theorem:5:6:item:a}{IX~5.6(a)}.

For \textup{(ii)}, let \(H\) be a subgroup prescheme of \(G\). To
prove that \(H\) has property \(\mathrm{AT}\), we may suppose that
\(S\) is the spectrum of a complete discrete valuation ring. We
must show that every coherent family of \(S_n\)-group morphisms
\[
  u_n:(\mathbb G_{m,S})_n\longrightarrow H_n
\]
is admissible (\SGARefLink{sga3:xv:proposition:8:5}{Proposition 8.5} and \SGARefLink{sga3:xv:proposition:8-1}{8.1} \textup{iii bis)}).
By hypothesis, \(G\) has property \(\mathrm{AT}\), so there is an
\(S\)-group morphism
\[
  u:\mathbb G_{m,S}\longrightarrow G
\]
lifting every \(u_n\). As in the proof of \SGARefLink{sga3:xv:proposition:8:5}{Proposition 8.5}, the
restriction of \(u\) to \({}_\ell\mathbb G_{m,S}\), with \(\ell\)
ranging over powers of a prime \(q\) invertible on \(S\), factors
through \(H\). By density, \(u_t\) factors through \(H_t\) on the
generic fibre. Since \(\mathbb G_{m,S}\) is reduced, \(u\) itself
factors through \(H\).

For \textup{(iii)}, let \(S'\in\mathcal S\), with the notation of
\SGARefLink{sga3:xv:definition:8-4}{Definition 8.4}. The group \(L\) is flat over \(S\), and its generic
fibre \(L_t\) is affine. Under these conditions \(L_s\) is
necessarily affine (Exposé XVII, Appendix III~2). Since \(G\) and
\(L\) are flat, \(G=G'\), and the dimensions of the fibres of \(G\)
and \(L\) are constant on \(S\) (Exposé VI\textsubscript{B}~4).
Consequently
\[
 \operatorname{rk}_{\mathrm{ab}}G_s
 \leq \dim G_s-\dim L_s
 =\dim G_t-\dim L_t
 =\operatorname{rk}_{\mathrm{ab}}G_t.
\]
By hypothesis these inequalities are equalities. Thus
\((G_s/L_s)^0_{\mathrm{red}}\) is an abelian variety and has
reductive rank zero. It follows that \(L_s\) has the same reductive
rank as \(G_s\).

We can now prove the principal theorems of this section.

\SGAStatement{\SGAInlineRefTarget{sga3:xv:theorem:8-8}Theorem 8.8}{Let \(G\) be a group prescheme of finite
type over a locally noetherian prescheme \(S\). Suppose that \(G\)
is flat over \(S\) and has connected fibres.
\begin{enumerate}
\renewcommand{\labelenumi}{\textup{\alph{enumi})}}
\item\SGAInlineRefTarget{sga3:xv:theorem:8-8:item:a} The functor \(\mathcal T_C\) of central sub-tori of \(G\) is
representable if and only if \(G\) has property \(\mathrm{ATC}\).
In that case, \(\mathcal T_C\) is represented by an \(S\)-prescheme
étale and separated over \(S\).
\item\SGAInlineRefTarget{sga3:xv:theorem:8-8:item:b} Under the hypotheses of \textup{(a)}, for every \(S\)-torus
\(T\), the functor
\(\operatorname{Hom}^{\mathrm{cent}}_{S\text{-gr}}(T,G)\) of central
homomorphisms from \(T\) into \(G\) is represented by an
\(S\)-prescheme étale and separated over \(S\).
\end{enumerate}}
\SGARefTarget{sga3:xv:theorem:8:8}

\SGAStatement{Theorem 8.9}{Let \(S\) be a locally noetherian
prescheme and let \(G\) be an \(S\)-group prescheme of finite type,
smooth over \(S\).
\begin{enumerate}
\renewcommand{\labelenumi}{\textup{\alph{enumi})}}
\item\SGAInlineRefTarget{sga3:xv:theorem:8-9:item:a} The functor \(\mathcal T\) of sub-tori of \(G\) is
representable if and only if \(G\) has property \(\mathrm{AT}\).
In that case, \(\mathcal T\) is represented by an \(S\)-prescheme
smooth and separated over \(S\). More precisely, its structural
morphism admits a canonical factorization
\[
  \mathcal T\xrightarrow{u}Y\xrightarrow{v}S,
\]
where \(v\) is smooth and quasi-projective (hence of finite type)
and \(u\) is étale and separated.
\item\SGAInlineRefTarget{sga3:xv:theorem:8-9:item:b} Under the hypotheses of \textup{(a)}, for every \(S\)-torus
\(T\), the functor \(\operatorname{Hom}_{S\text{-gr}}(T,G)\) of
homomorphisms from \(T\) into \(G\) is represented by an
\(S\)-prescheme smooth and separated over \(S\).
\end{enumerate}}
\SGARefTarget{sga3:xv:theorem:8:9}

% END inputs/XV/tex/components/08_expose_XV_section8_through_theorem8_9.tex

% BEGIN inputs/XV/tex/components/09_expose_XV_section8_remainder_and_end.tex
% Authority: Expo15.pdf, local pp. 61--68 / combined pp. 931--938 /
% printed pp. 470--483.
% Completes Exposé XV and stops before Exposé XVI.

\SGAProof{Proof of \SGARefLink{sga3:xv:theorem:8-8:item:a}{Theorem 8.8(a)}}
If the functor \(\mathcal T_C\) is representable, it commutes with
adic limits of artinian rings. Consequently, by \SGARefLink{sga3:xv:remark:8:2}{8.2} and
Proposition 8.5(b), \(G\) has property \(\mathrm{ATC}\). To prove the
converse, we shall use the following result, which will be proved in
EGA VI and may also be found in Murre's exposé,
\emph{Séminaire Bourbaki}, May 1965, no.~294, Theorem~1,
Corollary~2.

\SGAStatement{\SGAInlineRefTarget{sga3:xv:lemma:8-10}Lemma 8.10}{Let \(S\) be a locally noetherian
prescheme and let \(F\) be a contravariant functor from
\(\operatorname{Sch}/S\) to the category of sets. In order that
\(F\) be representable by an \(S\)-prescheme étale and separated
over \(S\), it is necessary and sufficient that \(F\) satisfy the
following five properties:
\begin{enumerate}
\renewcommand{\labelenumi}{\textup{\roman{enumi})}}
\item \(F\) is a sheaf for the fpqc topology (Exposé IV).
\item \(F\) commutes with filtered inductive limits of rings
(Exposé \SGARefLink{sga3:xi:remark:3:2}{XI~3.2}).
\item \(F\) commutes with adic limits of local artinian rings
(\SGARefLink{sga3:xv:proposition:8-1}{8.1}\textup{(i)}).
\item \(F\) satisfies the valuative criterion for separatedness:
for every \(S\)-prescheme \(S'\) that is the spectrum of a discrete
valuation ring, with generic point \(t\), the canonical map
\[
  F(S')\longrightarrow F(t)
\]
is injective.
\item \(F\) is infinitesimally étale (Exposé \SGARefLink{sga3:xi:remark:1:8}{XI~1.8}).
\end{enumerate}}
\SGARefTarget{sga3:xv:lemma:8:10}

We show that the functor \(\mathcal T_C\) of \SGARefLink{sga3:xv:theorem:8-8}{Theorem 8.8} satisfies
the five conditions of \SGARefLink{sga3:xv:lemma:8-10}{Lemma 8.10}.

\begin{enumerate}
\renewcommand{\labelenumi}{\textup{\roman{enumi})}}
\item The functor \(\mathcal T\), and likewise \(\mathcal T_C\), is
a sheaf for the fpqc topology whenever \(G\) is of finite
presentation over \(S\). Indeed, every monomorphism from a torus
into \(G\) is then an immersion (Exposé \SGARefLink{sga3:viii:proposition:7:9}{VIII~7.9}), and the assertion
follows from fpqc descent for subpreschemes.

\item \SGARefLink{sga3:xv:corollary:6-3}{Corollary 6.3} proves that \(\mathcal T\) commutes with filtered
inductive limits of rings when \(G\) is of finite presentation over
\(S\); the same property for \(\mathcal T_C\) follows immediately.

\item By \SGARefLink{sga3:xv:proposition:8:5}{Proposition 8.5}, condition \textup{(iii)} is precisely
equivalent to property \(\mathrm{ATC}\).

\item If \(S\) is the spectrum of a discrete valuation ring, two
sub-tori of \(G\) with the same generic fibre coincide. More
precisely, each is the identity component of the schematic closure
in \(G\) of that common generic fibre.

\item This follows from \SGARefLink{sga3:xv:corollary:2:3}{Proposition 2.3} because \(G\) is flat over
\(S\).
\end{enumerate}

\SGAProof{Proof of \SGARefLink{sga3:xv:theorem:8-8:item:b}{Theorem 8.8(b)}}
Proceeding as in Exposé \SGARefLink{sga3:xi:corollary:4:2}{XI~4.2}, it is enough to prove that the
product group \(T\times_S G\) again has property
\(\mathrm{ATC}\), which is immediate from the definition.

\SGAProof{Proof of \SGARefLink{sga3:xv:theorem:8-9:item:a}{Theorem 8.9(a)}}
After replacing \(G\), if necessary, by its identity component
(Exposé VI\textsubscript{B}~\SGARefLink{sga3:xv:lemma:3:10}{3.10}), we may suppose that \(G\) has
connected fibres. If \(T\) is a sub-torus of \(G\), its centralizer
\(\operatorname{Centr}_G(T)\) is represented (Exposé \SGARefLink{sga3:xi:corollary:6:11}{XI~6.11}) by a
subgroup prescheme of \(G\), smooth over \(S\)
(Exposé \SGARefLink{sga3:xi:corollary:2:4}{XI~2.4}), with connected fibres
(Exposé \SGARefLink{sga3:xii:theorem:6-6:item:b}{XII~6.6(b)}). It is therefore an element of
\(\mathcal{CT}(S)\), where \(\mathcal{CT}\) is the functor defined
in \SGARefLink{sga3:xv:proposition:7-1:item:i}{7.1(i)}. The rule
\[
  T\longmapsto\operatorname{Centr}_G(T)
\]
clearly defines an \(S\)-morphism
\[
  u:\mathcal T\longrightarrow\mathcal{CT}.
\]
Since \(\mathcal{CT}\) is represented by an \(S\)-prescheme smooth
and quasi-projective over \(S\) (\SGARefLink{sga3:xv:proposition:7-1:item:i}{7.1(i)}), it remains only to prove
that \(u\) is representable by a separated étale morphism.

After a suitable base change \(S'\to S\), with
\(S'=\mathcal{CT}\) and hence \(S'\) locally noetherian, we are
reduced to the following problem. Let \(S\) be a locally noetherian
prescheme, let \(G\) be an \(S\)-group prescheme smooth and of
finite type over \(S\), with connected fibres, and let \(H\) be a
subgroup of \(G\), smooth over \(S\), with connected fibres.
Consider the subfunctor \(X\) of \(\mathcal T_{C,H}\) such that,
for every \(S\)-prescheme \(S'\), \(X(S')\) is the set of central
sub-tori \(T\) of \(H_{S'}\) for which
\[
  \operatorname{Centr}_{G_{S'}}(T)=H_{S'}.
\]
We shall show that \(X\) is represented by an \(S\)-prescheme
étale and separated over \(S\).

Indeed, \(G\) has property \(\mathrm{AT}\) by hypothesis, and so
does \(H\) by Proposition 8.7(ii). Since \(H\) has connected
fibres, it also has property \(\mathrm{ATC}\) by
Proposition 8.7(i). Moreover, \(H\) is smooth over \(S\), hence
flat. Theorem 8.8(a) therefore represents \(\mathcal T_{C,H}\) by
an \(S\)-prescheme étale and separated over \(S\). It is enough to
show that the canonical monomorphism
\[
  X\longrightarrow\mathcal T_{C,H}
\]
is represented by an open immersion.

Put \(S'=\mathcal T_{C,H}\), and let \(K\) be the centralizer in
\(G_{S'}\) of the universal central torus of \(H_{S'}\). The group
\(K\) is a smooth group scheme over \(S'\), with connected fibres,
and contains \(H_{S'}\). By definition,
\[
  X=\prod_{K/S'} H_{S'}/K,
\]
which is represented by the open-and-closed subprescheme of \(S'\)
over which \(H_{S'}\) and \(K\) have the same relative dimension.

The proof of \SGARefLink{sga3:xv:theorem:8-9:item:b}{Theorem 8.9(b)} is analogous to that of
\SGARefLink{sga3:xv:theorem:8-8:item:b}{Theorem 8.8(b)}.

\SGAStatement{Corollary 8.11}{Let \(S\) be a prescheme and let
\(G\) be an \(S\)-group prescheme smooth and of finite presentation
over \(S\). If the abelian rank of the fibres of \(G\) is a locally
constant function on \(S\)---in particular, if the fibres of \(G\)
are affine---then the functor of sub-tori of \(G\) is represented
by an \(S\)-prescheme smooth and separated over \(S\). The same
holds for the functor
\(\operatorname{Hom}_{S\text{-gr}}(T,G)\), for every \(S\)-torus
\(T\).}
\SGARefTarget{sga3:xv:corollary:8:11}

The assertion is local on \(S\), so we may suppose that \(S\) is
affine and that the abelian rank of the fibres of \(G\) is constant.
There are then a noetherian affine scheme \(S_0\) and an
\(S_0\)-group prescheme \(G_0\), smooth and of finite type over
\(S_0\), such that
\[
  G\simeq G_0\times_{S_0}S
\]
as \(S\)-group preschemes (Exposé VI\textsubscript{B}, \S10).
The abelian rank of the fibres of an \(S\)-group prescheme of finite
presentation is an ind-constructible function (6.3 bis). A standard
argument (EGA IV~8) therefore allows us, in the present situation,
to choose \(G_0\) so that the abelian rank of its fibres is constant
on \(S_0\). Proposition 8.7(iii) gives property \(\mathrm{AT}\) for
\(G_0\), and \SGARefLink{sga3:xv:theorem:8:8}{Theorem 8.9} then represents its functor of sub-tori by
an \(S_0\)-prescheme smooth and separated over \(S_0\). Base change
gives the assertion for \(G\). The argument for
\(\operatorname{Hom}_{S\text{-gr}}(T,G)\) is analogous.

\medskip\noindent\textbf{Generalization of \SGARefLink{sga3:xv:theorem:8:8}{Theorem 8.9}.}
The functor \(\mathcal T\) of sub-tori of a smooth group \(G\) need
not be representable. We now give sufficient conditions for a
subfunctor of \(\mathcal T\) to be representable.

\SGAStatement{\SGAInlineRefTarget{sga3:xv:proposition:8-12}Proposition 8.12}{Let \(S\) be a locally noetherian
prescheme, let \(G\) be an \(S\)-group prescheme smooth over \(S\)
with connected fibres, and let \(F\) be an \(S\)-subfunctor of the
functor \(\mathcal T\) of sub-tori of \(G\). Suppose that:
\begin{enumerate}
\renewcommand{\labelenumi}{\textup{\roman{enumi})}}
\item\SGAInlineRefTarget{sga3:xv:proposition:8-12:item:i} \(F\) is a sheaf for the fpqc topology (Exposé IV);
\item \(F\) commutes with filtered inductive limits of rings
(Exposé \SGARefLink{sga3:xi:remark:3:2}{XI~3.2});
\item \(F\) commutes with adic limits of local artinian rings
(Exposé \SGARefLink{sga3:xi:remark:3:3}{XI~3.3});
\item\SGAInlineRefTarget{sga3:xv:proposition:8-12:item:iv} \(F\) is infinitesimally smooth over \(S\)
(Exposé \SGARefLink{sga3:xi:remark:1:8}{XI~1.8});
\item\SGAInlineRefTarget{sga3:xv:proposition:8-12:item:v} \(F\) is stable under inner automorphisms of \(G\): for every
\(S\)-prescheme \(S'\),
\[
  T\in F(S'),\quad g\in G(S')
  \quad\Longrightarrow\quad
  \operatorname{int}(g)T\in F(S').
\]
\end{enumerate}
Then \(F\) is represented by an \(S\)-prescheme smooth and
separated over \(S\).}
\SGARefTarget{sga3:xv:proposition:8:12}

\SGAStatement{Proposition 8.12 bis}{Let \(S\) and \(G\) be as in
\SGARefLink{sga3:xv:proposition:8-12}{Proposition 8.12}, let \(T\) be an \(S\)-torus, and let \(F\) be a
subfunctor of \(\operatorname{Hom}_{S\text{-gr}}(T,G)\). Suppose
that \(F\) satisfies conditions \textup{(i)}--\textup{(iv)} of
\SGARefLink{sga3:xv:proposition:8-12}{Proposition 8.12} and is stable under inner automorphisms of \(G\):
for every \(S\)-prescheme \(S'\),
\[
  u\in F(S'),\quad g\in G(S')
  \quad\Longrightarrow\quad
  \operatorname{int}(g)u\in F(S').
\]
Then \(F\) is represented by an \(S\)-prescheme smooth and
separated over \(S\).}
\SGARefTarget{sga3:xv:proposition:8:12:bis}

\SGAProof{Proof of \SGARefLink{sga3:xv:proposition:8-12}{Proposition 8.12}}
The proof of \SGARefLink{sga3:xv:proposition:8:12}{Proposition 8.12 bis} is entirely analogous and is left
to the reader. Let
\[
  u:F\longrightarrow\mathcal{CT}
\]
be the \(S\)-morphism which assigns to every
\(T\in F(S')\) the element
\(\operatorname{Centr}_{G_{S'}}(T)\) of
\(\mathcal{CT}(S')\) (\SGARefLink{sga3:xv:proposition:7-1:item:i}{7.1(i)}). Since \(\mathcal{CT}\) is represented
by an \(S\)-prescheme smooth and quasi-projective over \(S\), we
are reduced to proving that \(u\) is representable.

After the base change \(\mathcal{CT}\to S\), the problem is the
following. Given \(S\) and \(G\) as above and a smooth subgroup
\(H\) of \(G\) with connected fibres, represent the functor \(F_H\)
defined by
\[
  F_H(S')=
  \left\{T\in F(S')\ \middle|\
  H_{S'}=\operatorname{Centr}_{G_{S'}}(T)\right\}.
\]
We show that \(F_H\) is étale and separated over \(S\) by verifying
the five conditions of \SGARefLink{sga3:xv:lemma:8-10}{Lemma 8.10}.

The first three conditions follow directly from
\SGARefLink{sga3:xv:proposition:8-12:item:i}{Proposition 8.12(i)}--(iii). We have already observed that
\(\mathcal T_{C,H}\) is separated and infinitesimally étale.
Thus \(F_H\), being a subfunctor of \(\mathcal T_{C,H}\), is
separated and infinitesimally unramified (Exposé \SGARefLink{sga3:xi:remark:1:8}{XI~1.8}). It remains
to prove that \(F_H\) is infinitesimally smooth.

Let \(S\) be the spectrum of a local artinian ring, let
\(S_0\subseteq S\) be defined by a nilpotent ideal, and let
\(T_0\in F_H(S_0)\). By \SGARefLink{sga3:xv:proposition:8-12:item:iv}{Proposition 8.12(iv)}, \(T_0\) lifts to an
element \(T'\in F(S)\). Since \(H\) is smooth, \(T_0\) also lifts
to a sub-torus \(T''\) of \(H_S\) (Exposé \SGARefLink{sga3:ix:theorem:3:6:bis}{IX~3.6 bis}). This
sub-torus is conjugate to \(T'\) by an element \(g\in G(S)\)
(\emph{loc. cit.}), and hence \(T''\in F(S)\) by
\SGARefLink{sga3:xv:proposition:8-12:item:v}{Proposition 8.12(v)}. The centralizer
\(\operatorname{Centr}_{G_S}(T'')\) is smooth over \(S\) and agrees
with \(H_{S_0}\) over \(S_0\). Therefore \(T''\) lies in the centre
of \(H_S\) (Exposé \SGARefLink{sga3:ix:theorem:5:6:item:a}{IX~5.6(a)}), and its centralizer in \(G\) equals
\(H_S\). Thus \(T''\in F_H(S)\), and we take \(T=T''\).

\SGAStatement{Corollary 8.13}{Let \(S\) be a prescheme and let
\(G\) be an \(S\)-group prescheme smooth and of finite presentation
over \(S\), with connected fibres. Let \(\mathcal T_D\) be the
subfunctor of \(\mathcal T\) for which \(\mathcal T_D(S')\) is the
set of sub-tori \(T\) of \(G_{S'}\) such that, for every
\(s'\in S'\), \(T_{s'}\) is contained in the derived subgroup of
\(G_{s'}\) (Exposé VI\textsubscript{B}~7). Then \(\mathcal T_D\)
is represented by an \(S\)-prescheme smooth and separated over
\(S\).}
\SGARefTarget{sga3:xv:corollary:8:13}

\SGAStatement{Corollary 8.13 bis}{Let \(S\) and \(G\) be as in
\SGARefLink{sga3:xv:corollary:8:13}{Corollary 8.13}, and let \(T\) be an \(S\)-torus. Let
\[
  \operatorname{Hom}^{\mathrm{der}}_{S\text{-gr}}(T,G)
\]
be the subfunctor of \(\operatorname{Hom}_{S\text{-gr}}(T,G)\)
whose \(S'\)-points are the \(S'\)-morphisms
\(u:T_{S'}\to G_{S'}\) such that, for every \(s'\in S'\), the
morphism \(u_{s'}\) factors through the derived group of
\(G_{s'}\). Then this functor is represented by an \(S\)-prescheme
smooth and separated over \(S\).}
\SGARefTarget{sga3:xv:corollary:8:13:bis}

The two corollaries have analogous proofs; we prove
\SGARefLink{sga3:xv:corollary:8:13:bis}{Corollary 8.13 bis}. By the usual procedure, we reduce to the case
in which \(S\) is noetherian. We verify the five conditions of
\SGARefLink{sga3:xv:proposition:8:12}{Proposition 8.12 bis}.

Conditions \textup{(i)} and \textup{(iv)} follow immediately from
the corresponding properties of the functor
\[
  \operatorname{Hom}_{S\text{-gr}}(T,G).
\]
Condition \textup{(v)} holds because the derived group of an
algebraic group is invariant (Exposé VI\textsubscript{B}, \S7).

For condition \textup{(ii)}, a standard reduction (EGA IV~8)
reduces us to the following assertion. Let \(S\) be an integral
noetherian scheme with generic point \(\eta\), and let
\(u:T\to G\) be an \(S\)-group morphism whose generic fibre factors
through the derived group of \(G_\eta\). Then there is a
neighbourhood \(U\) of \(\eta\) such that, for every \(s\in U\),
\(u_s\) factors through the derived group of \(G_s\). This is an
immediate consequence of Exposé VI\textsubscript{B}~\SGARefLink{prop:sga3-VIB-10.12}{10.12}.

For condition \textup{(iii)}, retain the notation of \SGARefLink{sga3:xv:proposition:4:3}{Proposition
4.3}. We must show that an element
\[
  (u_m)_{m\in\mathbb N}\in
  \varprojlim_m
  \operatorname{Hom}^{\mathrm{der}}_{S_m\text{-gr}}(T_m,G_m)
\]
is admissible in the sense of \SGARefLink{sga3:xv:proposition:4:3}{Proposition 4.3} and comes from an
element of
\(\operatorname{Hom}^{\mathrm{der}}_{S\text{-gr}}(T,G)\).
As in \SGARefLink{sga3:xv:proposition:8-1}{8.1}, we reduce immediately to the case in which \(S\) is the
spectrum of a complete discrete valuation ring. Let \(t\) be its
generic point, \(s\) its closed point, and let \(D\) be the
schematic closure in \(G\) of the derived group \(D_t\) of \(G_t\).

\begin{enumerate}
\renewcommand{\labelenumi}{\textup{\alph{enumi})}}
\item The group \(D_s\) contains the derived group of \(G_s\).
Indeed, \(D_t\) is invariant in \(G_t\), and \(G\) is flat over
\(S\), so \(D\) is invariant in \(G\). Moreover, the morphism
\[
  G\times_S G\longrightarrow G,
  \qquad (x,y)\longmapsto xyx^{-1}y^{-1},
\]
factors through \(D_t\) on the generic fibre and therefore factors
through \(D\). Consequently \(G_s/D_s\) is commutative.

\item If
\(u_m\in
\operatorname{Hom}^{\mathrm{der}}_{S_m\text{-gr}}(T_m,G_m)\),
then \(u_m\) factors through \(D_m\). Its closed fibre \(u_s\)
factors through the derived group of \(G_s\), hence through \(D_s\)
by \textup{(a)}. Since \(D_m\) is flat over \(S_m\) and invariant
in \(G_m\), the quotient
\[
  H_m=G_m/D_m
\]
is representable (Exposé VI\textsubscript{A}, \S4). The image of
\(T_m\) in \(H_m\) is a torus (Exposé \SGARefLink{sga3:ix:theorem:6:8}{IX~6.8}) with trivial closed
fibre, hence is trivial. Thus \(u_m\) factors through \(D_m\).

\item The family \((u_m)_{m\in\mathbb N}\) is admissible and lifts
to a morphism \(T\to G\) belonging to
\(\operatorname{Hom}^{\mathrm{der}}_{S\text{-gr}}(T,G)\).
By what precedes and \SGARefLink{sga3:xv:theorem:4-1-bis}{Proposition 4.1 bis}, it is enough to prove that
\(D_t\) is affine. This follows from the next lemma.
\end{enumerate}

\SGAStatement{Lemma 8.14}{Let \(G\) be a smooth connected
algebraic group over a field \(k\). Then the derived group \(D\) of
\(G\) is affine.}
\SGARefTarget{sga3:xv:lemma:8:14}

Formation of \(D\) commutes with extension of the base field
(Exposé VI\textsubscript{B}~7), so we may suppose that \(k\) is
algebraically closed. The group \(G\) is then an extension of an
abelian variety \(A\) by a linear group \(L\). Since \(A\) is
commutative, \(D\) is necessarily contained in \(L\), and is
therefore affine.

\medskip\noindent\textbf{Maximal tori.}

\SGAStatement{\SGAInlineRefTarget{sga3:xv:theorem:8-15}Theorem 8.15}{Let \(S\) be a locally noetherian
prescheme and let \(G\) be an \(S\)-group prescheme smooth and of
finite type over \(S\). The following conditions are equivalent:
\begin{enumerate}
\renewcommand{\labelenumi}{\textup{\roman{enumi})}}
\item The \(S\)-functor \(\mathcal T_M\), whose \(S'\)-points are
the maximal tori of \(G_{S'}\) (Exposé \SGARefLink{sga3:xii:1.3}{XII~1.3}), is representable.
\item The functor \(\mathcal T_M\) is represented by an
\(S\)-prescheme smooth and quasi-projective over \(S\), with affine
fibres.
\item The group \(G\) admits a maximal torus locally for the étale
topology.
\item The group \(G\) admits a maximal torus locally for the
faithfully flat topology.
\item\SGAInlineRefTarget{sga3:xv:theorem:8-15:item:v} The group \(G\) has property \(\mathrm{AT}\), and the
reductive rank of its fibres is a locally constant function on
\(S\).
\end{enumerate}}
\SGARefTarget{sga3:xv:theorem:8:15}

\SGAProof{Proof}
\textup{(ii)}\(\Rightarrow\)\textup{(i)} is clear.

\textup{(i)}\(\Rightarrow\)\textup{(iii)}.
Since \(G\) is of finite presentation over \(S\), \SGARefLink{sga3:xv:corollary:6-3}{Corollary 6.3}
shows that \(\mathcal T_M\) commutes with filtered inductive limits
of rings. If representable, it is therefore locally of finite
presentation (EGA IV~8.14). Moreover, \(\mathcal T_M\) is formally
smooth (Exposé XI~2.1 bis). Hence its representing prescheme is
smooth over \(S\), and \textup{(iii)} follows from Hensel's lemma
(Exposé \SGARefLink{sga3:xi:proposition:1:10}{XI~1.10}).

\textup{(iii)}\(\Rightarrow\)\textup{(iv)} is clear.

\textup{(iv)}\(\Rightarrow\)\textup{(v)}.
Let \(s\in S\). There is an \(S\)-prescheme \(S'\), flat over \(S\)
and with image containing \(s\), such that \(G_{S'}\) has a maximal
torus \(T'\). If \(s'\in S'\) lies over \(s\), the reductive rank of
the fibres of \(G_{S'}\) is constant on
\(\operatorname{Spec}\mathcal O_{S',s'}\). Thus the reductive rank
of the fibres of \(G\) is constant on its image
\(\operatorname{Spec}\mathcal O_{S,s}\)
(EGA IV~2.3.4(ii)). Since this rank is a locally constructible
function on \(S\) (6.3 bis), it is constant on a neighbourhood of
\(s\) (EGA IV~1.10.1).

To prove property \(\mathrm{AT}\), let \(S_1\) be an \(S\)-scheme
which is the spectrum of a discrete valuation ring, whose closed
point \(s_1\) maps to \(s\). The prescheme
\[
  S'_1=S_1\times_S S'
\]
is faithfully flat over \(S_1\), and \(G_{S'_1}\) has a maximal
torus. Let \(A\) and \(A'\) be the rings of \(S_1\) and \(S'_1\),
respectively. Writing \(A'\) as the filtered inductive limit of its
finitely generated \(A\)-subalgebras and applying \SGARefLink{sga3:xv:corollary:6-3}{Corollary 6.3},
we obtain an \(S_1\)-scheme \(S''_1\) such that
\(G_{S''_1}\) has a maximal torus and
\(S''_1\to S_1\) is of finite type and surjective. By
EGA II~7.1.9, we may reduce to the case in which \(S'_1\) is the
spectrum of a discrete valuation ring. Then \(G_{S'_1}\), and
hence \(G_{S_1}\), has property \(\mathrm{AT}\). Since this holds
for every such \(S_1\), \(G\) has property \(\mathrm{AT}\).

\textup{(v)}\(\Rightarrow\)\textup{(i)}.
By \SGARefLink{sga3:xv:theorem:8:8}{Theorem 8.9} the functor \(\mathcal T\) of sub-tori of \(G\) is
representable. If the reductive rank of \(G\) is the constant
function \(r\), then \(\mathcal T_M\) is represented by the
open-and-closed subprescheme of \(\mathcal T\) representing the
subfunctor of tori of rank \(r\).

\textup{(iii)}\(\Rightarrow\)\textup{(ii)}.
Under condition \textup{(iii)}, the results of Exposé \SGARefLink{sga3:xii:7.1}{XII~7.1}
identify \(\mathcal T_M\) canonically with the functor of Cartan
subgroups of \(G\). The assertion follows from Theorem 7.3(i).

\medskip\noindent\textbf{Remark 8.16.}
The prescheme \(\mathcal T_M\) of maximal tori of \(G\) can in fact
be shown to be affine over \(S\),\footnote{Cf. M. Raynaud,
\emph{Faisceaux amples sur les schémas en groupes et les espaces
homogènes} (thesis, forthcoming). Editor's note: see
\emph{Lecture Notes in Mathematics} 119 (1970), Springer,
especially \SGARefLink{sga3:ix:corollary:2:9}{IX~2.9}.} generalizing Exposé \SGARefLink{sga3:xii:5.4}{XII~5.4}.
\SGARefTarget{sga3:xv:remark:8:16}

\SGAStatement{Corollary 8.17}{Let \(S\) be a prescheme and let
\(G\) be an \(S\)-group prescheme smooth and of finite presentation
over \(S\). If the abelian and reductive ranks of the fibres of
\(G\) are locally constant functions on \(S\), then \(G\) satisfies
the equivalent conditions \textup{(i)}--\textup{(iv)} of
\SGARefLink{sga3:xv:theorem:8-15}{Theorem 8.15}.}
\SGARefTarget{sga3:xv:corollary:8:17}

We may suppose that the abelian and reductive ranks of the fibres
of \(G\) are constant. Proceeding as in \SGARefLink{sga3:xv:corollary:8:11}{Corollary 8.11} and using
\SGARefLink{sga3:xv:corollary:6:3:bis}{6.3 bis}, we reduce to the case in which \(S\) is noetherian. Then
\(G\) has property \(\mathrm{AT}\) by \SGARefLink{sga3:xv:proposition:8:7}{Proposition 8.7}, and the
conclusion follows from \SGARefLink{sga3:xv:theorem:8-15:item:v}{Theorem 8.15(v)}.

\SGAStatement{\SGAInlineRefTarget{sga3:xv:corollary:8-18}Corollary 8.18}{Let \(S\) be a prescheme and let
\(G\) be an \(S\)-group prescheme smooth over \(S\). The following
conditions are equivalent:
\begin{enumerate}
\renewcommand{\labelenumi}{\textup{\alph{enumi})}}
\item The unipotent rank \(\rho_u\) and the abelian rank
\(\rho_{\mathrm{ab}}\) of the fibres of \(G\) are locally constant
functions on \(S\).
\item The unipotent rank \(\rho_u\) is locally constant, and \(G\)
admits maximal tori locally for the fpqc topology.
\item The reductive rank \(\rho_r\) and the abelian rank
\(\rho_{\mathrm{ab}}\) of the fibres of \(G\) are locally constant
functions on \(S\).
\end{enumerate}}
\SGARefTarget{sga3:xv:corollary:8:18}

\medskip\noindent\textbf{Remark 8.19.}
Under the hypotheses of \SGARefLink{sga3:xv:corollary:8-18}{Corollary 8.18}, a more precise argument
using lower semicontinuity of the abelian rank, announced in
Exposé \SGARefLink{sga3:x:remark:8-7}{X~8.7}, shows that if two of the three ranks
\(\rho_u,\rho_r,\rho_{\mathrm{ab}}\) are locally constant, then so
is the third.
\SGARefTarget{sga3:xv:remark:8:19}

\SGAProof{Proof of \SGARefLink{sga3:xv:corollary:8-18}{Corollary 8.18}}
After replacing \(G\), if necessary, by its identity component, we
may suppose that \(G\) is of finite presentation over \(S\)
(Exposé VI\textsubscript{B}~\SGARefLink{par:sga3-VIA-5.3.3}{5.3.3}).

\textup{(a)}\(\Rightarrow\)\textup{(c)}.
Let \(s\in S\). Since \(\rho_{\mathrm{ab}}\) is locally constant,
\SGARefLink{sga3:xv:corollary:8:11}{Corollary 8.11} allows us, after an étale extension covering \(S\),
to suppose that \(G\) has a sub-torus \(T\) whose fibre \(T_s\) is
a maximal torus of \(G_s\). Put
\[
  C=\operatorname{Centr}_G(T).
\]
This is a subgroup prescheme of \(G\), smooth over \(S\), with
connected fibres. For every \(t\in S\), the group \(C_t\) contains
a Cartan subgroup \(C'_t\) of \(G_t\). After shrinking \(S\), we
may suppose that \(\rho_u\), \(\rho_{\mathrm{ab}}\), and the
relative dimension of \(C\) over \(S\) are constant. We then have
\[
\begin{aligned}
  \dim C_s=\dim C_t
  &\geq \dim C'_t\\
  &\geq \rho_u(t)+\rho_{\mathrm{ab}}(t)+\dim T_t\\
  &=\rho_u(s)+\rho_{\mathrm{ab}}(s)+\dim T_s\\
  &=\rho_\nu(s)=\dim C_s.
\end{aligned}
\]
It follows that \(C_t=C'_t\). Thus \(T\) is a maximal torus of
\(G\), and in particular \(\rho_r(t)=\rho_r(s)\).

\textup{(c)}\(\Rightarrow\)\textup{(b)} follows from
\SGARefLink{sga3:xv:corollary:8:17}{Corollary 8.17}.

\textup{(b)}\(\Rightarrow\)\textup{(a)}.
Since \(G\) admits maximal tori locally for the fpqc topology, it
also admits Cartan subgroups locally for that topology. Hence, by
Exposé \SGARefLink{sga3:xii:7.3}{XII~7.3}, the nilpotent rank
\[
  \rho_\nu=\rho_u+\rho_r+\rho_{\mathrm{ab}}
\]
is locally constant. The existence of maximal tori makes
\(\rho_r\) locally constant, while \(\rho_u\) is locally constant
by hypothesis; therefore \(\rho_{\mathrm{ab}}\) is locally
constant.\footnote{The final sentence of the printed French
interchanges \(\rho_u\) and \(\rho_{\mathrm{ab}}\). The displayed
identity and condition \textup{(b)} give the implication in the
form written here.}

% END inputs/XV/tex/components/09_expose_XV_section8_remainder_and_end.tex


